% =====================================================================
%  PHỤ LỤC
%  Các mục được đánh số A, B, C…; bảng đánh số theo mục (Bảng A.1, B.1…).
% =====================================================================
\chapter*{PHỤ LỤC}
\addcontentsline{toc}{chapter}{PHỤ LỤC}

\setcounter{section}{0}
\renewcommand{\thesection}{\Alph{section}}
\renewcommand{\theHsection}{\Alph{section}}
\counterwithin{table}{section}
\counterwithin{figure}{section}
\counterwithin{code}{section}

Phần phụ lục bổ sung nội dung cho Chương~\ref{chap:khao-sat}, Chương~\ref{chap:dac-ta} và Chương~\ref{chap:backend}. Phụ lục~\ref{pl:tu-dien} mô tả thuộc tính các khái niệm nghiệp vụ; Phụ lục~\ref{pl:dac-ta-bo-sung} đặc tả các thao tác thành phần của use case tầng 2 và 3 bổ sung; kết hợp với 12 đặc tả trọng tâm và 12 mô tả ngắn trong Chương~\ref{chap:dac-ta}, tài liệu vẫn bao phủ bộ 53 use case. Phụ lục~\ref{pl:thong-diep} liệt kê các thông điệp hệ thống được dẫn chiếu trong đặc tả. Phụ lục~\ref{pl:truy-vet} trình bày ma trận liên kết yêu cầu với nghiệp vụ, use case, quy tắc nghiệp vụ và yêu cầu phi chức năng. Ba phụ lục tiếp theo là phần chi tiết của Chương~\ref{chap:backend}: Phụ lục~\ref{pl:hien-thuc-uc} trình bày ma trận và biểu đồ tuần tự hiện thực hóa các use case, Phụ lục~\ref{pl:dac-ta-api} đặc tả chi tiết các API bổ sung kèm test case dự kiến, Phụ lục~\ref{pl:dac-ta-bang} đặc tả các bảng cơ sở dữ liệu. Phụ lục~\ref{pl:so-do} chứa mã nguồn PlantUML của các sơ đồ trong báo cáo.

% ---------------------------------------------------------------------
%  PHỤ LỤC A. TỪ ĐIỂN DỮ LIỆU NGHIỆP VỤ
% ---------------------------------------------------------------------
\section{Từ điển dữ liệu nghiệp vụ}\label{pl:tu-dien}

Phụ lục này mô tả thuộc tính của các khái niệm trong Hình~\ref{fig:mo-hinh-khai-niem}. Kiểu dữ liệu được nêu ở mức nghiệp vụ, chưa gắn với hệ quản trị cơ sở dữ liệu cụ thể; lược đồ vật lý sẽ được xác định trong giai đoạn thiết kế. Mỗi khái niệm còn có hai thuộc tính kỹ thuật — thời điểm tạo và thời điểm cập nhật gần nhất — không lặp lại trong các bảng. Phụ lục cũng bổ sung các khái niệm hỗ trợ phát sinh từ đặc tả use case: thiết lập nhắc nhở, thiết bị nhận thông báo, mã OTP, công thức yêu thích, cấu hình hệ thống và nhật ký quản trị. Ký hiệu “Liệt kê” cho biết thuộc tính chỉ nhận các giá trị cố định được nêu ở cột mô tả.

\begin{ddtable}{Từ điển dữ liệu nhóm khái niệm tài khoản và nhóm gia đình}{tab:dd-tai-khoan}
\ddentity{Người dùng}
\dd{Mã người dùng}{Mã định danh}{Có}{Duy nhất, do hệ thống sinh}
\dd{Email}{Chuỗi, tối đa 100}{Có}{Duy nhất, đúng định dạng email (\br{BR01})}
\dd{Mật khẩu đã băm}{Chuỗi}{Có}{Kết quả băm có muối, không lưu mật khẩu rõ (\nfr{NFR14})}
\dd{Họ tên}{Chuỗi, 2–50}{Có}{Tên hiển thị trong nhóm}
\dd{Ảnh đại diện}{Tệp ảnh}{Không}{JPEG hoặc PNG, tối đa 5 MB}
\dd{Vai trò hệ thống}{Liệt kê}{Có}{Người dùng, Quản trị viên}
\dd{Trạng thái}{Liệt kê}{Có}{Chờ xác thực, Hoạt động, Khóa tạm, Bị khóa}
\dd{Số lần đăng nhập sai}{Số nguyên}{Có}{Đặt lại khi đăng nhập thành công (\br{BR03})}
\dd{Khóa đến}{Ngày giờ}{Không}{Thời điểm hết khóa tạm}
\dd{Ngôn ngữ}{Liệt kê}{Có}{Tiếng Việt (mặc định), Tiếng Anh}
\dd{Không gian gần nhất}{Tham chiếu Nhóm}{Không}{Không gian được mở khi đăng nhập}
\ddentity{Thiết lập nhắc nhở (thuộc Người dùng)}
\dd{Ngưỡng nhắc N}{Số nguyên, 1–7}{Có}{Mặc định lấy từ cấu hình hệ thống (\br{BR12})}
\dd{Bật nhắc hạn dùng}{Logic}{Có}{Mặc định bật}
\dd{Bật thông báo phân công}{Logic}{Có}{Mặc định bật}
\dd{Bật thông báo hoạt động nhóm}{Logic}{Có}{Thành viên mới, rời nhóm; mặc định bật}
\ddentity{Thiết bị nhận thông báo}
\dd{Mã thiết bị}{Chuỗi}{Có}{Mã do FCM cấp, duy nhất}
\dd{Người dùng}{Tham chiếu Người dùng}{Có}{Chủ sở hữu thiết bị}
\dd{Nền tảng}{Liệt kê}{Có}{Android, iOS}
\ddentity{Mã OTP}
\dd{Email}{Chuỗi}{Có}{Email nhận mã}
\dd{Mục đích}{Liệt kê}{Có}{Đăng ký, Khôi phục mật khẩu}
\dd{Mã đã băm}{Chuỗi}{Có}{Không lưu mã rõ}
\dd{Hết hạn lúc}{Ngày giờ}{Có}{Thời điểm tạo cộng thời hạn OTP (\br{BR02})}
\dd{Số lần nhập sai}{Số nguyên, 0–5}{Có}{Đạt 5 thì mã bị vô hiệu}
\ddentity{Nhóm}
\dd{Mã nhóm}{Mã định danh}{Có}{Duy nhất}
\dd{Tên nhóm}{Chuỗi, 3–50}{Có}{Với không gian cá nhân, mặc định là “Cá nhân” (\br{BR05})}
\dd{Loại}{Liệt kê}{Có}{Cá nhân, Gia đình (\br{BR04})}
\dd{Ảnh nhóm}{Tệp ảnh}{Không}{JPEG hoặc PNG}
\ddentity{Thành viên nhóm}
\dd{Nhóm}{Tham chiếu Nhóm}{Có}{Cặp (Nhóm, Người dùng) là duy nhất}
\dd{Người dùng}{Tham chiếu Người dùng}{Có}{Mỗi người thuộc tối đa một nhóm loại Gia đình (\br{BR04})}
\dd{Vai trò}{Liệt kê}{Có}{Trưởng nhóm, Thành viên; mỗi nhóm đúng một trưởng nhóm (\br{BR05})}
\dd{Ngày tham gia}{Ngày giờ}{Có}{}
\ddentity{Mã mời}
\dd{Mã}{Chuỗi, 8}{Có}{Chữ in hoa và chữ số, bỏ 0, O, 1, I (\br{BR06})}
\dd{Nhóm}{Tham chiếu Nhóm}{Có}{Nhóm loại Gia đình}
\dd{Người tạo}{Tham chiếu Người dùng}{Có}{Trưởng nhóm tại thời điểm tạo}
\dd{Hết hạn lúc}{Ngày giờ}{Có}{Mặc định sau 48 giờ}
\dd{Trạng thái}{Liệt kê}{Có}{Hiệu lực, Đã thu hồi, Hết hạn}
\end{ddtable}

\begin{ddtable}{Từ điển dữ liệu nhóm khái niệm mua sắm}{tab:dd-mua-sam}
\ddentity{Danh sách mua sắm}
\dd{Mã danh sách}{Mã định danh}{Có}{Duy nhất}
\dd{Nhóm}{Tham chiếu Nhóm}{Có}{Không gian sở hữu danh sách}
\dd{Tên}{Chuỗi, 1–100}{Có}{Gợi ý theo ngày dự kiến mua}
\dd{Ngày dự kiến mua}{Ngày}{Có}{Không ở quá khứ khi tạo (\br{BR08})}
\dd{Ghi chú}{Chuỗi, tối đa 500}{Không}{}
\dd{Trạng thái}{Liệt kê}{Có}{Đang lập, Đang mua, Hoàn tất; suy ra từ mặt hàng (\br{BR08})}
\dd{Người tạo}{Tham chiếu Người dùng}{Có}{}
\dd{Nguồn gốc}{Liệt kê}{Có}{Tạo mới, Sao chép, Sinh từ kế hoạch}
\ddentity{Mặt hàng mua sắm}
\dd{Mã mặt hàng}{Mã định danh}{Có}{Duy nhất}
\dd{Danh sách}{Tham chiếu Danh sách mua sắm}{Có}{Quan hệ hợp thành}
\dd{Thực phẩm}{Tham chiếu Thực phẩm}{Không}{Bỏ trống khi dùng tên tự do}
\dd{Tên tự do}{Chuỗi, 1–100}{Không}{Bắt buộc khi không tham chiếu thực phẩm}
\dd{Số lượng cần mua}{Số thực dương}{Có}{Tối đa hai chữ số thập phân (\br{BR14})}
\dd{Đơn vị}{Tham chiếu Đơn vị đo}{Có}{}
\dd{Ghi chú}{Chuỗi, tối đa 200}{Không}{}
\dd{Trạng thái}{Liệt kê}{Có}{Cần mua, Đã phân công, Đã mua, Đã nhập tủ lạnh, Đã bỏ qua (Hình~\ref{fig:trang-thai-mat-hang})}
\dd{Người được giao}{Tham chiếu Thành viên nhóm}{Không}{Tối đa một người (\br{BR10})}
\dd{Người mua}{Tham chiếu Người dùng}{Không}{Ghi khi đánh dấu đã mua}
\dd{Thời điểm mua}{Ngày giờ}{Không}{}
\dd{Số lượng thực tế}{Số thực dương}{Không}{Mặc định bằng số lượng cần mua}
\dd{Đơn giá}{Số nguyên không âm}{Không}{Đơn vị đồng, dùng cho báo cáo (\br{BR21})}
\end{ddtable}

\begin{ddtable}{Từ điển dữ liệu nhóm khái niệm danh mục và tủ lạnh}{tab:dd-tu-lanh}
\ddentity{Loại thực phẩm}
\dd{Mã loại}{Mã định danh}{Có}{Duy nhất}
\dd{Tên loại}{Chuỗi, 2–50}{Có}{Duy nhất, không phân biệt hoa thường (\br{BR22})}
\dd{Biểu tượng}{Tệp ảnh}{Không}{Dùng trong danh sách và bộ lọc}
\dd{Vị trí lưu trữ mặc định}{Liệt kê}{Có}{Một trong năm vị trí của Bảng~\ref{tab:vi-tri-luu-tru}}
\dd{Trạng thái}{Liệt kê}{Có}{Đang dùng, Ngừng sử dụng}
\ddentity{Đơn vị đo}
\dd{Ký hiệu}{Chuỗi, 1–20}{Có}{Duy nhất, ví dụ g, kg, lạng, ml, quả, bó}
\dd{Tên đầy đủ}{Chuỗi, 1–50}{Có}{}
\dd{Nhóm đại lượng}{Liệt kê}{Có}{Khối lượng, Thể tích, Số đếm (\br{BR16})}
\dd{Hệ số quy đổi}{Số thực dương}{Có}{So với đơn vị cơ sở của nhóm; ví dụ lạng = 100}
\dd{Trạng thái}{Liệt kê}{Có}{Đang dùng, Ngừng sử dụng}
\ddentity{Thực phẩm (danh mục)}
\dd{Mã thực phẩm}{Mã định danh}{Có}{Duy nhất}
\dd{Tên}{Chuỗi, 1–100}{Có}{Duy nhất trong cùng loại (\br{BR22})}
\dd{Tên gọi khác}{Danh sách chuỗi}{Không}{Hỗ trợ tìm kiếm, ví dụ “thịt heo” cho “thịt lợn”}
\dd{Loại}{Tham chiếu Loại thực phẩm}{Có}{}
\dd{Đơn vị mặc định}{Tham chiếu Đơn vị đo}{Có}{}
\dd{Hạn dùng gợi ý}{Số nguyên theo từng vị trí}{Không}{Số ngày bảo quản tại mỗi vị trí (Bảng~\ref{tab:vi-tri-luu-tru})}
\dd{Phạm vi}{Liệt kê}{Có}{Dùng chung, Riêng của một nhóm (do người dùng tạo)}
\dd{Trạng thái}{Liệt kê}{Có}{Đang dùng, Chờ duyệt, Ngừng sử dụng}
\ddentity{Thực phẩm trong tủ (lô)}
\dd{Mã lô}{Mã định danh}{Có}{Duy nhất}
\dd{Nhóm}{Tham chiếu Nhóm}{Có}{Không gian sở hữu}
\dd{Thực phẩm}{Tham chiếu Thực phẩm}{Có}{}
\dd{Số lượng ban đầu}{Số thực dương}{Có}{}
\dd{Số lượng còn lại}{Số thực không âm}{Có}{Không vượt số lượng ban đầu (\br{BR14})}
\dd{Đơn vị}{Tham chiếu Đơn vị đo}{Có}{}
\dd{Vị trí lưu trữ}{Liệt kê}{Có}{Ngăn mát, Ngăn đông, Ngăn rau củ, Cánh tủ, Kệ bếp}
\dd{Ngày mua}{Ngày}{Có}{Không sau ngày hiện tại}
\dd{Ngày hết hạn}{Ngày}{Có}{Không trước ngày mua (\br{BR11})}
\dd{Trạng thái}{Liệt kê}{Có}{Còn hạn, Sắp hết hạn, Đã hết hạn, Đã dùng hết, Đã loại bỏ (\br{BR12})}
\dd{Mặt hàng nguồn}{Tham chiếu Mặt hàng mua sắm}{Không}{Lưu vết nguồn gốc và đơn giá}
\dd{Ghi chú}{Chuỗi, tối đa 200}{Không}{}
\ddentity{Nhật ký sử dụng}
\dd{Lô}{Tham chiếu Thực phẩm trong tủ}{Có}{Quan hệ hợp thành}
\dd{Loại}{Liệt kê}{Có}{Tiêu thụ, Loại bỏ}
\dd{Số lượng}{Số thực dương}{Có}{Theo đơn vị của lô}
\dd{Lý do}{Liệt kê}{Không}{Hết hạn, Hỏng, Khác; bắt buộc khi loại là Loại bỏ}
\dd{Bữa ăn liên quan}{Tham chiếu Bữa ăn dự kiến}{Không}{Ghi khi tiêu thụ qua \uc{UC19}}
\dd{Người thực hiện}{Tham chiếu Người dùng}{Có}{}
\dd{Thời điểm}{Ngày giờ}{Có}{}
\end{ddtable}

\begin{ddtable}{Từ điển dữ liệu nhóm khái niệm công thức, kế hoạch và hệ thống}{tab:dd-cong-thuc}
\ddentity{Công thức}
\dd{Mã công thức}{Mã định danh}{Có}{Duy nhất}
\dd{Tên món}{Chuỗi, 2–100}{Có}{(\br{BR17})}
\dd{Ảnh}{Tệp ảnh}{Không}{Tối đa 5 MB, nén tự động}
\dd{Số khẩu phần chuẩn}{Số nguyên, 1–20}{Có}{}
\dd{Thời gian chuẩn bị, nấu}{Số nguyên phút}{Không}{0–600}
\dd{Bữa phù hợp}{Tập liệt kê}{Không}{Sáng, Trưa, Tối}
\dd{Các bước thực hiện}{Danh sách văn bản}{Không}{Có thứ tự; mỗi bước tối đa 500 ký tự}
\dd{Loại công thức}{Liệt kê}{Có}{Cá nhân, Mẫu (do quản trị viên quản lý)}
\dd{Không gian sở hữu}{Tham chiếu Nhóm}{Không}{Bắt buộc với công thức cá nhân: công thức thuộc bộ sưu tập của không gian hiện hành, nên các thành viên nhóm dùng chung trong kế hoạch bữa ăn; bỏ trống với công thức mẫu}
\dd{Trạng thái}{Liệt kê}{Có}{Nháp, Công bố, Ẩn (với công thức mẫu, \uc{UC23}); công thức cá nhân luôn ở trạng thái Công bố}
\dd{Người tạo}{Tham chiếu Người dùng}{Có}{}
\dd{Công thức gốc}{Tham chiếu Công thức}{Không}{Khi được sao chép từ công thức mẫu}
\ddentity{Nguyên liệu công thức}
\dd{Công thức}{Tham chiếu Công thức}{Có}{Quan hệ hợp thành}
\dd{Thực phẩm}{Tham chiếu Thực phẩm}{Có}{Mỗi thực phẩm xuất hiện một lần trong công thức}
\dd{Số lượng}{Số thực dương}{Có}{Ứng với số khẩu phần chuẩn}
\dd{Đơn vị}{Tham chiếu Đơn vị đo}{Có}{}
\dd{Bắt buộc}{Logic}{Có}{Chỉ nguyên liệu bắt buộc được tính vào $R$ (\br{BR19})}
\ddentity{Công thức yêu thích}
\dd{Người dùng}{Tham chiếu Người dùng}{Có}{Cặp (Người dùng, Công thức) là duy nhất}
\dd{Công thức}{Tham chiếu Công thức}{Có}{}
\ddentity{Bữa ăn dự kiến}
\dd{Mã}{Mã định danh}{Có}{Duy nhất}
\dd{Nhóm}{Tham chiếu Nhóm}{Có}{Không gian sở hữu}
\dd{Ngày}{Ngày}{Có}{Từ hôm nay đến 28 ngày tiếp theo khi tạo (\br{BR18})}
\dd{Bữa}{Liệt kê}{Có}{Sáng, Trưa, Tối}
\dd{Công thức}{Tham chiếu Công thức}{Không}{Bỏ trống khi là món tự do}
\dd{Tên món tự do}{Chuỗi, 1–100}{Không}{Bắt buộc khi không có công thức}
\dd{Số khẩu phần}{Số nguyên, 1–20}{Có}{}
\dd{Trạng thái}{Liệt kê}{Có}{Dự kiến, Đã nấu}
\dd{Thứ tự}{Số nguyên}{Có}{Thứ tự hiển thị trong bữa}
\ddentity{Thông báo}
\dd{Mã thông báo}{Mã định danh}{Có}{Duy nhất}
\dd{Người nhận}{Tham chiếu Người dùng}{Có}{}
\dd{Loại}{Liệt kê}{Có}{Nhắc hạn dùng, Phân công, Hoạt động nhóm, Hệ thống}
\dd{Tiêu đề, nội dung}{Chuỗi}{Có}{Nội dung tóm tắt, không chứa dữ liệu nhạy cảm}
\dd{Đối tượng liên quan}{Tham chiếu}{Không}{Danh sách, lô hoặc nhóm để mở đúng màn hình}
\dd{Khóa chống trùng}{Chuỗi}{Không}{Với nhắc hạn dùng là người nhận + ngày, bảo đảm mỗi người nhận tối đa một thông báo tổng hợp mỗi ngày (\br{BR13})}
\dd{Đã đọc}{Logic}{Có}{Mặc định chưa đọc}
\dd{Thời điểm gửi}{Ngày giờ}{Có}{}
\ddentity{Cấu hình hệ thống}
\dd{Khóa tham số}{Chuỗi}{Có}{Duy nhất, ví dụ \texttt{reminder.defaultDays}}
\dd{Giá trị}{Chuỗi theo kiểu của tham số}{Có}{Trong miền hợp lệ của Bảng~\ref{tab:tham-so-he-thong}}
\dd{Người cập nhật}{Tham chiếu Người dùng}{Có}{Quản trị viên}
\ddentity{Nhật ký quản trị}
\dd{Người thực hiện}{Tham chiếu Người dùng}{Có}{Quản trị viên}
\dd{Hành động}{Chuỗi}{Có}{Ví dụ khóa tài khoản, sửa danh mục, đổi cấu hình}
\dd{Giá trị cũ, giá trị mới}{Văn bản}{Không}{Phục vụ truy vết}
\dd{Thời điểm}{Ngày giờ}{Có}{}
\end{ddtable}

% ---------------------------------------------------------------------
%  PHỤ LỤC B. ĐẶC TẢ CÁC USE CASE BỔ SUNG
% ---------------------------------------------------------------------
\section{Đặc tả các use case bổ sung}\label{pl:dac-ta-bo-sung}

Phụ lục này ghi các use case tầng 2 và 3 chưa được đặc tả trong Chương~\ref{chap:dac-ta}, theo cùng mẫu bảng đặc tả và cùng hệ thống mã tham chiếu. Khi một use case gộp nhiều thao tác, mỗi thao tác là một biến thể (A, B, C) với bước đánh số có tiền tố là chữ của biến thể, trừ biến thể A. Use case tầng 2 chỉ cần cài đặt luồng chính; use case tầng 3 chỉ cài đặt ở mức tối giản, bản đặc tả được giữ để tham khảo và không phải tiêu chí chấp nhận. Luồng ngoại tuyến nếu xuất hiện trong các bảng dưới đây là định hướng phát triển.

% =====================================================================
\subsection{Gói Xác thực và quản lý tài khoản}

\begin{ucspec}{UC01}{Xác thực mã OTP}
\ucfield{Tác nhân}{Khách (chính); Dịch vụ thư điện tử (phụ)}
\ucfield{Mô tả}{Xác minh quyền sở hữu email bằng mã dùng một lần gửi qua thư; được bao gồm bởi \uc{UC01} và \uc{UC03}.}
\ucfield{Điều kiện trước}{Một mã OTP còn hiệu lực đã được sinh cho email của Khách.}
\ucflow{Luồng thực thi chính}
\ucstep{1}{Hệ thống}{Hiển thị màn hình nhập mã gồm 6 ô số, đồng hồ đếm ngược thời hạn mã và nút “Gửi lại mã” bị vô hiệu trong 60 giây đầu.}
\ucstep{2}{Khách}{Nhập 6 chữ số; ứng dụng tự gửi khi nhập đủ hoặc tự điền khi hệ điều hành đọc được mã từ thư.}
\ucstep{3}{Hệ thống}{So khớp mã, kiểm tra thời hạn và số lần nhập sai (\br{BR02}).}
\ucstep{4}{Hệ thống}{Đánh dấu mã đã dùng và trả kết quả thành công cho use case gọi.}
\ucflow{Luồng thực thi mở rộng}
\ucstep{2a}{Khách}{Nhấn “Gửi lại mã” sau 60 giây: hệ thống kiểm tra giới hạn 5 lần mỗi giờ, sinh mã mới, vô hiệu mã cũ và gửi thư; nếu chưa đủ thời gian chờ, hiển thị \msg{MSG30}.}
\ucstep{3a}{Hệ thống}{Mã sai: tăng số lần nhập sai, hiển thị \msg{MSG04}; quay lại bước 2.}
\ucstep{3b}{Hệ thống}{Mã hết hạn hoặc đã nhập sai đủ 5 lần: vô hiệu mã, hiển thị \msg{MSG04} và yêu cầu gửi lại mã.}
\ucfield{Điều kiện sau}{Thành công: mã được đánh dấu đã dùng, use case gọi được tiếp tục. Thất bại: không có thay đổi trạng thái tài khoản.}
\ucfield{Quy tắc nghiệp vụ}{\br{BR02}}
\end{ucspec}

\begin{ucspec}{UC03}{Đặt lại và đổi mật khẩu}
\ucfield{Tác nhân}{Khách (chính); Dịch vụ thư điện tử (phụ, thông qua \uc{UC01}); Người dùng}
\ucfield{Mô tả}{\textbf{A.} Cho phép người quên mật khẩu đặt lại mật khẩu mới sau khi xác thực email bằng mã OTP. \textbf{B.} Cho phép người dùng đổi mật khẩu khi đang đăng nhập.}
\ucfield{Điều kiện trước}{\textbf{A.} Khách chưa đăng nhập. \textbf{B.} Người dùng đang đăng nhập.}
\ucflow{Biến thể A – Khôi phục mật khẩu: luồng thực thi chính}
\ucstep{1}{Khách}{Chọn “Quên mật khẩu” tại màn hình đăng nhập và nhập email.}
\ucstep{2}{Hệ thống}{Nếu email thuộc một tài khoản hoạt động, sinh mã OTP mục đích Khôi phục và gửi thư; luôn hiển thị cùng một thông điệp “Nếu email tồn tại, mã đã được gửi” để không tiết lộ email nào đã đăng ký.}
\ucstep{3}{Hệ thống}{Thực hiện \uc{UC01}.}
\ucstep{4}{Khách}{Nhập mật khẩu mới và xác nhận mật khẩu.}
\ucstep{5}{Hệ thống}{Kiểm tra độ mạnh mật khẩu (\br{BR01}) và sự trùng khớp.}
\ucstep{6}{Hệ thống}{Cập nhật mật khẩu, thu hồi mọi phiên đăng nhập hiện có (\br{BR03}), đặt lại bộ đếm đăng nhập sai và chuyển về màn hình đăng nhập với \msg{MSG10}.}
\ucflow{Biến thể A – Khôi phục mật khẩu: luồng thực thi mở rộng}
\ucstep{1a}{Hệ thống}{Email sai định dạng: hiển thị \msg{MSG02}.}
\ucstep{5a}{Hệ thống}{Mật khẩu không đủ mạnh hoặc không khớp: hiển thị \msg{MSG02} hoặc \msg{MSG29}; quay lại bước 4.}
\ucflow{Biến thể B – Đổi mật khẩu: luồng thực thi chính}
\ucstep{B1}{Người dùng}{Chọn “Đổi mật khẩu”, nhập mật khẩu hiện tại, mật khẩu mới và xác nhận mật khẩu mới.}
\ucstep{B2}{Hệ thống}{Kiểm tra mật khẩu hiện tại đúng, mật khẩu mới đạt \br{BR01}, khác mật khẩu hiện tại và trùng khớp với xác nhận.}
\ucstep{B3}{Hệ thống}{Cập nhật mật khẩu, thu hồi các phiên trên thiết bị khác, cấp phiên mới cho thiết bị hiện tại, hiển thị \msg{MSG10}.}
\ucflow{Biến thể B – Đổi mật khẩu: luồng thực thi mở rộng}
\ucstep{B2a}{Hệ thống}{Mật khẩu hiện tại sai: hiển thị \msg{MSG02} và tăng bộ đếm sai theo \br{BR03}.}
\ucstep{B2b}{Hệ thống}{Mật khẩu mới không hợp lệ hoặc không khớp: hiển thị \msg{MSG02} hoặc \msg{MSG29}.}
\ucfield{Điều kiện sau}{\textbf{A.} Mật khẩu của tài khoản được thay mới; các phiên cũ không còn hiệu lực. \textbf{B.} Mật khẩu được thay mới; chỉ thiết bị hiện tại còn phiên hợp lệ.}
\ucfield{Quy tắc nghiệp vụ}{\br{BR01}, \br{BR02}, \br{BR03}}
\end{ucspec}

\begin{ucspec}{UC02}{Đăng xuất}
\ucfield{Tác nhân}{Người dùng}
\ucfield{Mô tả}{Kết thúc phiên làm việc trên thiết bị hiện tại.}
\ucfield{Điều kiện trước}{Người dùng đang đăng nhập.}
\ucflow{Luồng thực thi chính}
\ucstep{1}{Người dùng}{Mở màn hình Tài khoản, chọn “Đăng xuất”.}
\ucstep{2}{Hệ thống}{Hiển thị hộp thoại xác nhận; nếu còn thay đổi ngoại tuyến chưa đồng bộ, cảnh báo số thay đổi sẽ bị mất.}
\ucstep{3}{Người dùng}{Xác nhận đăng xuất.}
\ucstep{4}{Hệ thống}{Thu hồi mã làm mới trên máy chủ, hủy đăng ký mã thiết bị nhận thông báo, xóa mã và dữ liệu lưu đệm trên thiết bị, chuyển về màn hình chào.}
\ucflow{Luồng thực thi mở rộng}
\ucstep{2a}{Hệ thống}{Có thay đổi chưa đồng bộ và thiết bị có mạng: thử đồng bộ trước khi đăng xuất.}
\ucstep{4a}{Hệ thống}{Không kết nối được máy chủ: vẫn xóa dữ liệu trên thiết bị; mã làm mới sẽ tự hết hạn theo \br{BR03}.}
\ucfield{Điều kiện sau}{Thiết bị không còn lưu thông tin phiên và không nhận thông báo của tài khoản.}
\ucfield{Quy tắc nghiệp vụ}{\br{BR03}}
\end{ucspec}




% =====================================================================
\subsection{Gói Quản lý nhóm gia đình}

\begin{ucspec}{UC05}{Xem thông tin nhóm và thành viên}
\ucfield{Tác nhân}{Thành viên nhóm}
\ucfield{Mô tả}{Cho phép thành viên xem tên nhóm, ảnh nhóm, danh sách thành viên và vai trò của từng người.}
\ucfield{Điều kiện trước}{Người dùng đã đăng nhập.}
\ucflow{Luồng thực thi chính}
\ucstep{1}{Thành viên nhóm}{Mở màn hình Gia đình.}
\ucstep{2}{Hệ thống}{Hiển thị tên, ảnh nhóm, số thành viên và danh sách thành viên kèm ảnh đại diện, vai trò, ngày tham gia; trưởng nhóm được đặt đầu danh sách.}
\ucstep{3}{Hệ thống}{Hiển thị tóm tắt hoạt động gần đây của nhóm như thành viên mới, danh sách vừa hoàn tất.}
\ucflow{Luồng thực thi mở rộng}
\ucstep{2a}{Hệ thống}{Người xem là trưởng nhóm: hiển thị thêm nút “Mời thành viên” (\uc{UC06}) và trình đơn trên từng thành viên gồm “Xóa khỏi nhóm” (\uc{UC07}), “Chuyển quyền trưởng nhóm” (\uc{UC07}).}
\ucstep{2b}{Hệ thống}{Người dùng chưa thuộc nhóm gia đình nào: hiển thị hai lựa chọn “Tạo nhóm” (\uc{UC05}) và “Tham gia nhóm” (\uc{UC06}).}
\ucfield{Điều kiện sau}{Không thay đổi dữ liệu.}
\ucfield{Quy tắc nghiệp vụ}{\br{BR05}, \br{BR23}}
\end{ucspec}

\begin{ucspec}{UC07}{Rời nhóm, xóa thành viên và chuyển quyền}
\ucfield{Tác nhân}{Trưởng nhóm; Thành viên nhóm}
\ucfield{Mô tả}{\textbf{A.} Cho phép trưởng nhóm loại một thành viên ra khỏi nhóm. \textbf{B.} Cho phép trưởng nhóm trao vai trò trưởng nhóm cho một thành viên khác. \textbf{C.} Cho phép thành viên rời khỏi nhóm gia đình và trở về không gian cá nhân.}
\ucfield{Điều kiện trước}{\textbf{A.} Người thực hiện là trưởng nhóm; thành viên bị xóa không phải chính trưởng nhóm. \textbf{B.} Người thực hiện là trưởng nhóm; nhóm có ít nhất hai thành viên. \textbf{C.} Người dùng là thành viên của một nhóm gia đình.}
\ucflow{Biến thể A – Xóa thành viên khỏi nhóm: luồng thực thi chính}
\ucstep{1}{Trưởng nhóm}{Tại \uc{UC05}, mở trình đơn của một thành viên và chọn “Xóa khỏi nhóm”.}
\ucstep{2}{Hệ thống}{Hiển thị \msg{MSG19}, nêu số mặt hàng đang giao cho thành viên này sẽ được trả về trạng thái Cần mua.}
\ucstep{3}{Trưởng nhóm}{Xác nhận.}
\ucstep{4}{Hệ thống}{Xóa tư cách thành viên, trả các mặt hàng đang giao về Cần mua (\br{BR07}), thông báo cho người bị xóa và chuyển không gian của họ về không gian cá nhân.}
\ucflow{Biến thể A – Xóa thành viên khỏi nhóm: luồng thực thi mở rộng}
\ucstep{3a}{Trưởng nhóm}{Hủy: không có thay đổi.}
\ucstep{4a}{Hệ thống}{Người thực hiện không còn là trưởng nhóm: hiển thị \msg{MSG13}.}
\ucflow{Biến thể B – Chuyển quyền trưởng nhóm: luồng thực thi chính}
\ucstep{B1}{Trưởng nhóm}{Chọn “Chuyển quyền trưởng nhóm” trên một thành viên, hoặc được chuyển tới đây từ \uc{UC07}.}
\ucstep{B2}{Hệ thống}{Hiển thị hộp thoại xác nhận nêu rõ sau khi chuyển, người thực hiện sẽ mất quyền phân công và quản lý thành viên.}
\ucstep{B3}{Trưởng nhóm}{Xác nhận.}
\ucstep{B4}{Hệ thống}{Trong một giao dịch, đổi vai trò người nhận thành Trưởng nhóm và người chuyển thành Thành viên (\br{BR05}); thu hồi mã mời đang hiệu lực; thông báo cho cả nhóm.}
\ucflow{Biến thể B – Chuyển quyền trưởng nhóm: luồng thực thi mở rộng}
\ucstep{B3a}{Trưởng nhóm}{Hủy: không có thay đổi.}
\ucstep{B4a}{Hệ thống}{Người nhận vừa rời nhóm: hiển thị lỗi và yêu cầu chọn người khác.}
\ucflow{Biến thể C – Rời nhóm gia đình: luồng thực thi chính}
\ucstep{C1}{Thành viên nhóm}{Tại màn hình Gia đình, chọn “Rời nhóm”.}
\ucstep{C2}{Hệ thống}{Kiểm tra vai trò và số thành viên của nhóm (\br{BR07}).}
\ucstep{C3}{Hệ thống}{Hiển thị \msg{MSG19} xác nhận rời nhóm.}
\ucstep{C4}{Thành viên nhóm}{Xác nhận.}
\ucstep{C5}{Hệ thống}{Xóa tư cách thành viên, trả các mặt hàng đang giao về Cần mua, chuyển không gian về cá nhân, thông báo cho trưởng nhóm.}
\ucflow{Biến thể C – Rời nhóm gia đình: luồng thực thi mở rộng}
\ucstep{C2a}{Hệ thống}{Người rời là trưởng nhóm và nhóm còn thành viên khác: hiển thị \msg{MSG24} và chuyển sang biến thể B (chuyển quyền); sau khi chuyển quyền thành công, tiếp tục bước C3.}
\ucstep{C2b}{Hệ thống}{Người rời là thành viên duy nhất: cảnh báo nhóm sẽ bị giải tán cùng toàn bộ dữ liệu chung; nếu xác nhận, giải tán nhóm.}
\ucfield{Điều kiện sau}{\textbf{A.} Người bị xóa không còn truy cập được dữ liệu của nhóm; dữ liệu họ đã đóng góp vẫn thuộc về nhóm. \textbf{B.} Nhóm vẫn có đúng một trưởng nhóm là người nhận quyền. \textbf{C.} Người dùng không còn thuộc nhóm; nhóm vẫn có đúng một trưởng nhóm hoặc đã bị giải tán.}
\ucfield{Quy tắc nghiệp vụ}{\br{BR05}, \br{BR06}, \br{BR07}, \br{BR10}, \br{BR23}}
\end{ucspec}



% =====================================================================
\subsection{Gói Quản lý danh sách mua sắm}

\begin{ucspec}{UC08}{Quản lý danh sách và mặt hàng mua sắm}
\ucfield{Tác nhân}{Người dùng; Trưởng nhóm (chính); Dịch vụ thông báo đẩy (phụ)}
\ucfield{Mô tả}{\textbf{A.} Cho phép tạo danh sách mới từ một danh sách cũ, phục vụ các lần mua sắm lặp lại. \textbf{B.} Cho phép trưởng nhóm giao một hoặc nhiều mặt hàng cho một thành viên phụ trách mua; người được giao nhận thông báo đẩy. \textbf{C.} Cho phép người dùng xem lại các danh sách đã hoàn tất theo thời gian.}
\ucfield{Điều kiện trước}{\textbf{A.} Không gian hiện hành có ít nhất một danh sách mua sắm. \textbf{B.} Trưởng nhóm đang xem một danh sách chưa hoàn tất của nhóm; người được giao có thể là một thành viên khác hoặc chính trưởng nhóm. \textbf{C.} Người dùng đang đăng nhập.}
\ucflow{Biến thể A – Sao chép danh sách mua sắm đã có: luồng thực thi chính}
\ucstep{1}{Người dùng}{Chọn “Tạo từ danh sách cũ” tại \uc{UC08}, hoặc “Mua lại” trên một danh sách tại \uc{UC08}.}
\ucstep{2}{Hệ thống}{Hiển thị các mặt hàng của danh sách nguồn, mặc định chọn tất cả, kèm ô ngày dự kiến mua mới.}
\ucstep{3}{Người dùng}{Bỏ chọn mặt hàng không cần, chỉnh số lượng, chọn ngày và nhấn “Tạo”.}
\ucstep{4}{Hệ thống}{Tạo danh sách mới với nguồn gốc Sao chép; các mặt hàng ở trạng thái Cần mua, không giữ phân công, đơn giá và người mua cũ.}
\ucflow{Biến thể A – Sao chép danh sách mua sắm đã có: luồng thực thi mở rộng}
\ucstep{3a}{Hệ thống}{Ngày dự kiến mua ở quá khứ: hiển thị \msg{MSG02}.}
\ucflow{Biến thể B – Phân công người mua: luồng thực thi chính}
\ucstep{B1}{Trưởng nhóm}{Nhấn giữ một mặt hàng để vào chế độ chọn nhiều, chọn thêm các mặt hàng khác, nhấn “Phân công”.}
\ucstep{B2}{Hệ thống}{Hiển thị danh sách thành viên kèm số mặt hàng mỗi người đang được giao trong danh sách này.}
\ucstep{B3}{Trưởng nhóm}{Chọn một thành viên và nhấn “Giao việc”.}
\ucstep{B4}{Hệ thống}{Kiểm tra quyền trưởng nhóm và kiểm tra các mặt hàng được chọn đều ở trạng thái Cần mua hoặc Đã phân công (\br{BR10}).}
\ucstep{B5}{Hệ thống}{Ghi nhận người được giao, chuyển các mặt hàng sang Đã phân công, đồng bộ tới các thành viên.}
\ucstep{B6}{Hệ thống}{Tạo thông báo “Bạn được giao mua k mặt hàng trong danh sách <tên>” và yêu cầu Dịch vụ thông báo đẩy gửi tới thiết bị của người được giao.}
\ucstep{B7}{Hệ thống}{Hiển thị \msg{MSG10}; mỗi mặt hàng hiển thị ảnh đại diện của người được giao.}
\ucflow{Biến thể B – Phân công người mua: luồng thực thi mở rộng}
\ucstep{B3a}{Trưởng nhóm}{Chọn “Bỏ phân công”: hệ thống chuyển các mặt hàng về Cần mua và thông báo cho người được giao trước đó.}
\ucstep{B4a}{Hệ thống}{Người thực hiện không còn là trưởng nhóm (do vừa chuyển quyền trên thiết bị khác): hiển thị \msg{MSG13}; kết thúc use case.}
\ucstep{B4b}{Hệ thống}{Có mặt hàng đã ở trạng thái Đã mua hoặc Đã bỏ qua: bỏ qua các mặt hàng đó, chỉ phân công các mặt hàng hợp lệ và thông báo số mặt hàng bị bỏ qua.}
\ucstep{B6a}{Hệ thống}{Gửi thông báo đẩy thất bại: thông báo vẫn được lưu trong hộp thư trong ứng dụng, hệ thống thử gửi lại tối đa 3 lần (\nfr{NFR12}).}
\ucflow{Biến thể C – Xem lịch sử mua sắm: luồng thực thi chính}
\ucstep{C1}{Người dùng}{Tại màn hình Mua sắm, chọn thẻ “Lịch sử”.}
\ucstep{C2}{Hệ thống}{Hiển thị các danh sách Hoàn tất của không gian hiện hành, nhóm theo tháng, mỗi danh sách gồm ngày mua, số mặt hàng, tổng chi tiêu (nếu có đơn giá) và những người đã mua.}
\ucstep{C3}{Người dùng}{Chọn một danh sách.}
\ucstep{C4}{Hệ thống}{Hiển thị chi tiết ở chế độ chỉ đọc kèm nút “Mua lại”.}
\ucflow{Biến thể C – Xem lịch sử mua sắm: luồng thực thi mở rộng}
\ucstep{C2a}{Hệ thống}{Chưa có danh sách nào hoàn tất: hiển thị trạng thái rỗng với hướng dẫn.}
\ucstep{C4a}{Người dùng}{Nhấn “Mua lại”: thực hiện \uc{UC08}.}
\ucfield{Điều kiện sau}{\textbf{A.} Danh sách mới ở trạng thái Đang lập; danh sách nguồn không thay đổi. \textbf{B.} Các mặt hàng được chọn có đúng một người được giao; người được giao nhận được thông báo trong hộp thư và trên thiết bị. \textbf{C.} Không thay đổi dữ liệu.}
\ucfield{Quy tắc nghiệp vụ}{\br{BR07}, \br{BR08}, \br{BR09}, \br{BR10}}
\end{ucspec}



% =====================================================================
\subsection{Gói Quản lý thực phẩm trong tủ lạnh}

\begin{ucspec}{UC12}{Xem và lọc thực phẩm trong tủ lạnh}
\ucfield{Tác nhân}{Người dùng}
\ucfield{Mô tả}{\textbf{A.} Cho phép người dùng xem toàn bộ tồn kho của không gian hiện hành theo vị trí lưu trữ và trạng thái hạn dùng. \textbf{B.} Hiển thị tập trung các lô Sắp hết hạn và Đã hết hạn cùng các món ăn gợi ý để sử dụng kịp thời.}
\ucfield{Điều kiện trước}{\textbf{A.} Người dùng đang đăng nhập. \textbf{B.} Người dùng đang đăng nhập; mở từ thẻ trên màn hình chính hoặc từ thông báo nhắc nhở.}
\ucflow{Biến thể A – Xem danh sách thực phẩm trong tủ lạnh: luồng thực thi chính}
\ucstep{1}{Người dùng}{Mở màn hình Tủ lạnh.}
\ucstep{2}{Hệ thống}{Hiển thị dải thống kê nhanh (tổng số lô, số lô sắp hết hạn, số lô đã hết hạn) và các thẻ vị trí Ngăn mát, Ngăn đông, Ngăn rau củ, Cánh tủ, Kệ bếp.}
\ucstep{3}{Hệ thống}{Trong mỗi vị trí, hiển thị thực phẩm gộp theo tên với tổng lượng còn lại và ngày hết hạn gần nhất, sắp xếp mặc định theo hạn gần nhất; màu và nhãn theo \nfr{NFR04}.}
\ucstep{4}{Người dùng}{Chạm vào một thực phẩm để xem các lô chi tiết.}
\ucflow{Biến thể A – Xem danh sách thực phẩm trong tủ lạnh: luồng thực thi mở rộng}
\ucstep{2a}{Hệ thống}{Tủ lạnh trống: hiển thị trạng thái rỗng với nút “Thêm thực phẩm” và “Cất hàng vừa mua”.}
\ucstep{3a}{Người dùng}{Đổi cách sắp xếp theo tên hoặc theo ngày nhập, hoặc áp bộ lọc (\uc{UC20}).}
\ucstep{4a}{Người dùng}{Trên một lô, chọn “Sửa” (\uc{UC13}), “Đã dùng” (\uc{UC13}) hoặc “Bỏ đi” (\uc{UC13}).}
\ucflow{Biến thể B – Xem thực phẩm sắp hết hạn: luồng thực thi chính}
\ucstep{B1}{Người dùng}{Mở “Sắp hết hạn”.}
\ucstep{B2}{Hệ thống}{Hiển thị hai nhóm “Sắp hết hạn” và “Đã hết hạn”, mỗi lô kèm số ngày còn lại hoặc số ngày quá hạn, sắp xếp tăng dần theo ngày hết hạn.}
\ucstep{B3}{Hệ thống}{Hiển thị ba món gợi ý hàng đầu từ \uc{UC18} có sử dụng các thực phẩm trong nhóm Sắp hết hạn.}
\ucflow{Biến thể B – Xem thực phẩm sắp hết hạn: luồng thực thi mở rộng}
\ucstep{B2a}{Người dùng}{Vuốt trên một lô để chọn nhanh “Đã dùng” (\uc{UC13}) hoặc “Bỏ đi” (\uc{UC13}).}
\ucstep{B3a}{Người dùng}{Chọn một món gợi ý: thực hiện \uc{UC16} hoặc thêm vào lịch theo \uc{UC17}.}
\ucfield{Điều kiện sau}{\textbf{A.} Không thay đổi dữ liệu. \textbf{B.} Không thay đổi dữ liệu nếu người dùng chỉ xem.}
\ucfield{Quy tắc nghiệp vụ}{\br{BR12}, \br{BR19}, \br{BR23}}
\end{ucspec}


\begin{ucspec}{UC13}{Điều chỉnh lô thực phẩm (sửa, sử dụng, loại bỏ)}
\ucfield{Tác nhân}{Người dùng}
\ucfield{Mô tả}{\textbf{A.} Cho phép người dùng sửa số lượng còn lại, ngày hết hạn, vị trí và ghi chú của một lô. \textbf{B.} Cho phép người dùng ghi nhận lượng thực phẩm đã lấy ra sử dụng; hệ thống trừ tồn kho theo nguyên tắc hết hạn trước dùng trước. \textbf{C.} Cho phép người dùng ghi nhận một lô hoặc một phần lô bị bỏ đi do hỏng hoặc hết hạn.}
\ucfield{Điều kiện trước}{\textbf{A.} Lô có lượng còn lại lớn hơn 0. \textbf{B.} Thực phẩm có ít nhất một lô với lượng còn lại lớn hơn 0 trong không gian hiện hành. \textbf{C.} Lô có lượng còn lại lớn hơn 0.}
\ucflow{Biến thể A – Cập nhật thông tin thực phẩm: luồng thực thi chính}
\ucstep{1}{Người dùng}{Chọn “Sửa” trên một lô tại \uc{UC12}.}
\ucstep{2}{Hệ thống}{Hiển thị biểu mẫu với dữ liệu hiện có.}
\ucstep{3}{Người dùng}{Sửa thông tin và nhấn “Lưu”.}
\ucstep{4}{Hệ thống}{Kiểm tra như bước B6 của \uc{UC11}; lưu, tính lại trạng thái hạn dùng (\br{BR12}); hiển thị \msg{MSG10}.}
\ucflow{Biến thể A – Cập nhật thông tin thực phẩm: luồng thực thi mở rộng}
\ucstep{3a}{Người dùng}{Đổi vị trí lưu trữ: hệ thống hỏi có muốn cập nhật ngày hết hạn theo hạn dùng gợi ý của vị trí mới tính từ hôm nay hay không, ví dụ khi chuyển thịt sang ngăn đông.}
\ucstep{4a}{Hệ thống}{Dữ liệu không hợp lệ: hiển thị \msg{MSG14} hoặc \msg{MSG16}.}
\ucflow{Biến thể B – Ghi nhận sử dụng thực phẩm: luồng thực thi chính}
\ucstep{B1}{Người dùng}{Trên màn hình Tủ lạnh, chọn một thực phẩm và nhấn “Đã dùng”.}
\ucstep{B2}{Hệ thống}{Hiển thị tổng lượng còn lại, các lô theo thứ tự ngày hết hạn tăng dần và ô nhập lượng sử dụng với các nút nhanh “Một nửa”, “Hết”.}
\ucstep{B3}{Người dùng}{Nhập lượng sử dụng hoặc chọn nút nhanh, nhấn “Xác nhận”.}
\ucstep{B4}{Hệ thống}{Kiểm tra lượng sử dụng dương và không vượt tổng lượng còn lại của các lô chưa hết hạn (\br{BR14}).}
\ucstep{B5}{Hệ thống}{Trừ dần từ lô có ngày hết hạn sớm nhất (\br{BR15}); lô có lượng còn lại bằng 0 chuyển sang Đã dùng hết.}
\ucstep{B6}{Hệ thống}{Ghi nhật ký tiêu thụ cho từng lô bị trừ, cập nhật hiển thị và thông báo \msg{MSG10}.}
\ucflow{Biến thể B – Ghi nhận sử dụng thực phẩm: luồng thực thi mở rộng}
\ucstep{B3a}{Người dùng}{Chọn đích danh một lô (ví dụ hộp sữa đã mở): hệ thống chỉ trừ trên lô được chọn.}
\ucstep{B4a}{Hệ thống}{Lượng sử dụng vượt lượng còn lại: hiển thị \msg{MSG17}; quay lại bước B3.}
\ucstep{B4b}{Hệ thống}{Chỉ còn lô Đã hết hạn: hỏi người dùng có chắc chắn sử dụng lô quá hạn hay muốn chuyển sang \uc{UC13} Loại bỏ thực phẩm.}
\ucflow{Biến thể C – Loại bỏ thực phẩm: luồng thực thi chính}
\ucstep{C1}{Người dùng}{Chọn “Bỏ đi” trên một lô.}
\ucstep{C2}{Hệ thống}{Hiển thị lượng loại bỏ mặc định bằng toàn bộ lượng còn lại và danh sách lý do Hết hạn, Hỏng, Khác; lý do mặc định là Hết hạn nếu lô ở trạng thái Đã hết hạn.}
\ucstep{C3}{Người dùng}{Điều chỉnh lượng, chọn lý do và xác nhận.}
\ucstep{C4}{Hệ thống}{Trừ lượng còn lại, ghi nhật ký loại bỏ; nếu lượng còn lại bằng 0 thì lô chuyển sang Đã loại bỏ; hiển thị giá trị ước tính bị lãng phí nếu lô có đơn giá.}
\ucflow{Biến thể C – Loại bỏ thực phẩm: luồng thực thi mở rộng}
\ucstep{C3a}{Hệ thống}{Lượng loại bỏ vượt lượng còn lại: hiển thị \msg{MSG17}.}
\ucfield{Điều kiện sau}{\textbf{A.} Lô được cập nhật; nếu ngày hết hạn được kéo dài ra ngoài ngưỡng, lô trở về trạng thái Còn hạn. \textbf{B.} Tồn kho của thực phẩm giảm đúng bằng lượng sử dụng; nhật ký tiêu thụ được ghi nhận để phục vụ \uc{UC21}. \textbf{C.} Nhật ký loại bỏ được ghi nhận và được tính vào \uc{UC21}.}
\ucfield{Quy tắc nghiệp vụ}{\br{BR11}, \br{BR12}, \br{BR14}, \br{BR15}, \br{BR16}, \br{BR21}}
\end{ucspec}




% =====================================================================
\subsection{Gói Quản lý công thức nấu ăn}

\begin{ucspec}{UC15}{Quản lý công thức nấu ăn}
\ucfield{Tác nhân}{Người dùng}
\ucfield{Mô tả}{\textbf{A.} Cho phép người dùng sửa hoặc xóa công thức trong bộ sưu tập cá nhân. \textbf{B.} Cho phép người dùng thêm một công thức vào hoặc bỏ khỏi danh sách yêu thích.}
\ucfield{Điều kiện trước}{\textbf{A.} Công thức thuộc bộ sưu tập của người dùng (không phải công thức mẫu). \textbf{B.} Người dùng đang xem một công thức.}
\ucflow{Biến thể A – Cập nhật, xóa công thức: luồng thực thi chính}
\ucstep{1}{Người dùng}{Tại \uc{UC16}, chọn “Sửa”.}
\ucstep{2}{Hệ thống}{Mở biểu mẫu của \uc{UC15} với dữ liệu hiện có.}
\ucstep{3}{Người dùng}{Sửa và lưu.}
\ucstep{4}{Hệ thống}{Kiểm tra theo \br{BR17}, lưu và hiển thị \msg{MSG10}; các bữa ăn dự kiến tham chiếu công thức dùng nội dung mới.}
\ucflow{Biến thể A – Cập nhật, xóa công thức: luồng thực thi mở rộng}
\ucstep{1a}{Người dùng}{Chọn “Xóa”: hệ thống hiển thị \msg{MSG19}, nêu số món dự kiến đang dùng công thức; khi được xác nhận, các món Dự kiến dùng công thức chuyển thành món tự do giữ nguyên tên, các món Đã nấu giữ tham chiếu lịch sử và công thức bị ẩn khỏi bộ sưu tập.}
\ucflow{Biến thể B – Đánh dấu công thức yêu thích: luồng thực thi chính}
\ucstep{B1}{Người dùng}{Chạm biểu tượng trái tim.}
\ucstep{B2}{Hệ thống}{Đảo trạng thái yêu thích, cập nhật biểu tượng ngay; công thức yêu thích được hiển thị ở thẻ “Yêu thích” và nhận hệ số $F = 1$ khi gợi ý món (\br{BR19}).}
\ucflow{Biến thể B – Đánh dấu công thức yêu thích: luồng thực thi mở rộng}
\ucstep{B2a}{Hệ thống}{Mất kết nối mạng: lưu thay đổi trên thiết bị và đồng bộ sau.}
\ucfield{Điều kiện sau}{\textbf{A.} Công thức được cập nhật hoặc bị xóa mà không làm mất dữ liệu lịch sử. \textbf{B.} Trạng thái yêu thích của công thức đối với người dùng được cập nhật.}
\ucfield{Quy tắc nghiệp vụ}{\br{BR17}, \br{BR18}, \br{BR19}}
\end{ucspec}



% =====================================================================
\subsection{Gói Quản lý kế hoạch bữa ăn}



% =====================================================================
\subsection{Gói Tìm kiếm và phân loại thực phẩm}

\begin{ucspec}{UC20}{Lọc thực phẩm theo danh mục}
\ucfield{Tác nhân}{Người dùng}
\ucfield{Mô tả}{Cho phép người dùng thu hẹp danh sách thực phẩm hoặc kết quả tìm kiếm theo loại, vị trí lưu trữ và trạng thái hạn dùng.}
\ucfield{Điều kiện trước}{Người dùng đang xem kết quả tìm kiếm (\uc{UC20}) hoặc danh sách thực phẩm trong tủ (\uc{UC12}).}
\ucflow{Luồng thực thi chính}
\ucstep{1}{Người dùng}{Chạm biểu tượng bộ lọc.}
\ucstep{2}{Hệ thống}{Hiển thị bảng lọc gồm các loại thực phẩm (kèm số lượng mỗi loại), vị trí lưu trữ và trạng thái hạn dùng.}
\ucstep{3}{Người dùng}{Chọn một hoặc nhiều giá trị và nhấn “Áp dụng”.}
\ucstep{4}{Hệ thống}{Kết hợp các điều kiện (cùng nhóm là “hoặc”, khác nhóm là “và”), cập nhật danh sách và hiển thị các thẻ bộ lọc đang áp dụng.}
\ucflow{Luồng thực thi mở rộng}
\ucstep{3a}{Người dùng}{Chạm “x” trên một thẻ bộ lọc để bỏ điều kiện đó.}
\ucstep{4a}{Hệ thống}{Không có kết quả: hiển thị \msg{MSG18} và nút “Xóa bộ lọc”.}
\ucfield{Điều kiện sau}{Bộ lọc được ghi nhớ trong phiên làm việc của màn hình.}
\ucfield{Quy tắc nghiệp vụ}{\br{BR12}}
\end{ucspec}

% =====================================================================
\subsection{Gói Thống kê báo cáo}



% =====================================================================
\subsection{Gói Quản trị hệ thống}

\begin{ucspec}{UC24}{Quản lý tài khoản người dùng}
\ucfield{Tác nhân}{Quản trị viên}
\ucfield{Mô tả}{Cho phép quản trị viên tra cứu tài khoản, khóa, mở khóa tài khoản vi phạm và cấp, thu hồi vai trò quản trị.}
\ucfield{Điều kiện trước}{Quản trị viên đã đăng nhập.}
\ucflow{Luồng thực thi chính}
\ucstep{1}{Quản trị viên}{Mở “Tài khoản người dùng”.}
\ucstep{2}{Hệ thống}{Hiển thị danh sách tài khoản có phân trang với email, họ tên, trạng thái, ngày tạo, lần đăng nhập gần nhất; có tìm kiếm theo email, họ tên và lọc theo trạng thái.}
\ucstep{3}{Quản trị viên}{Chọn một tài khoản và chọn “Khóa tài khoản”, nhập lý do.}
\ucstep{4}{Hệ thống}{Hiển thị \msg{MSG19}; khi được xác nhận, chuyển tài khoản sang Bị khóa, thu hồi mọi phiên, ghi nhật ký quản trị và gửi thư thông báo cho chủ tài khoản.}
\ucflow{Luồng thực thi mở rộng}
\ucstep{2a}{Hệ thống}{Quản trị viên chỉ xem được thông tin tài khoản, không xem được nội dung dữ liệu của các nhóm (\br{BR23}).}
\ucstep{3a}{Quản trị viên}{Chọn “Mở khóa”: tài khoản trở về Hoạt động, ghi nhật ký.}
\ucstep{3b}{Quản trị viên}{Chọn “Cấp quyền quản trị” hoặc “Thu hồi quyền quản trị”: cập nhật vai trò hệ thống, ghi nhật ký; không cho phép tự thu hồi quyền của chính mình nếu là quản trị viên cuối cùng.}
\ucfield{Điều kiện sau}{Trạng thái hoặc vai trò của tài khoản được cập nhật và lưu vết.}
\ucfield{Quy tắc nghiệp vụ}{\br{BR03}, \br{BR23}}
\end{ucspec}

\begin{ucspec}{UC23}{Quản lý danh mục dùng chung}
\ucfield{Tác nhân}{Quản trị viên}
\ucfield{Mô tả}{\textbf{A.} Cho phép quản trị viên thêm, sửa, ngừng sử dụng các loại thực phẩm dùng để phân loại và lọc. \textbf{B.} Cho phép quản trị viên quản lý đơn vị đo, nhóm đại lượng và hệ số quy đổi về đơn vị cơ sở. \textbf{C.} Cho phép quản trị viên biên soạn, cập nhật, ẩn các công thức mẫu dùng chung cho mọi người dùng.}
\ucfield{Điều kiện trước}{Quản trị viên đã đăng nhập.}
\ucflow{Biến thể A – Quản lý loại mặt hàng: luồng thực thi chính}
\ucstep{1}{Quản trị viên}{Mở “Loại mặt hàng”.}
\ucstep{2}{Hệ thống}{Hiển thị danh sách loại với tên, biểu tượng, vị trí lưu trữ mặc định, số thực phẩm thuộc loại.}
\ucstep{3}{Quản trị viên}{Thêm hoặc sửa một loại, nhập tên, chọn biểu tượng và vị trí mặc định, nhấn “Lưu”.}
\ucstep{4}{Hệ thống}{Kiểm tra tên duy nhất (\br{BR22}), lưu và hiển thị \msg{MSG10}.}
\ucflow{Biến thể A – Quản lý loại mặt hàng: luồng thực thi mở rộng}
\ucstep{3a}{Quản trị viên}{Chọn “Xóa” trên một loại đang có thực phẩm: hệ thống hiển thị \msg{MSG22} và chỉ cho phép chuyển loại sang Ngừng sử dụng.}
\ucstep{4a}{Hệ thống}{Tên trùng: hiển thị \msg{MSG02}.}
\ucflow{Biến thể B – Quản lý đơn vị đo lường: luồng thực thi chính}
\ucstep{B1}{Quản trị viên}{Mở “Đơn vị đo lường”.}
\ucstep{B2}{Hệ thống}{Hiển thị đơn vị nhóm theo đại lượng Khối lượng, Thể tích, Số đếm, mỗi đơn vị kèm hệ số quy đổi.}
\ucstep{B3}{Quản trị viên}{Thêm hoặc sửa một đơn vị, ví dụ “lạng”, nhóm Khối lượng, hệ số 100; nhấn “Lưu”.}
\ucstep{B4}{Hệ thống}{Kiểm tra ký hiệu duy nhất, hệ số dương, đơn vị cơ sở của mỗi nhóm có hệ số bằng 1 (\br{BR16}); lưu và hiển thị \msg{MSG10}.}
\ucflow{Biến thể B – Quản lý đơn vị đo lường: luồng thực thi mở rộng}
\ucstep{B3a}{Quản trị viên}{Sửa hệ số hoặc nhóm đại lượng của đơn vị đang được tham chiếu: hệ thống cảnh báo ảnh hưởng tới phép quy đổi và chỉ cho phép sau khi xác nhận lần hai.}
\ucstep{B3b}{Quản trị viên}{Chọn “Xóa” trên một đơn vị đang được tham chiếu: hiển thị \msg{MSG22}; đơn vị được chuyển sang Ngừng sử dụng, bị ẩn khỏi lựa chọn mới nhưng dữ liệu cũ giữ nguyên.}
\ucstep{B4a}{Hệ thống}{Dữ liệu không hợp lệ: hiển thị \msg{MSG02}.}
\ucflow{Biến thể C – Quản lý công thức mẫu: luồng thực thi chính}
\ucstep{C1}{Quản trị viên}{Mở “Công thức mẫu”.}
\ucstep{C2}{Hệ thống}{Hiển thị danh sách công thức mẫu với trạng thái Công bố hoặc Nháp, số lượt lưu về bộ sưu tập.}
\ucstep{C3}{Quản trị viên}{Thêm hoặc sửa một công thức theo biểu mẫu của \uc{UC15}, chọn trạng thái và nhấn “Lưu”.}
\ucstep{C4}{Hệ thống}{Kiểm tra theo \br{BR17}, yêu cầu mọi nguyên liệu tham chiếu thực phẩm dùng chung; lưu và hiển thị \msg{MSG10}.}
\ucflow{Biến thể C – Quản lý công thức mẫu: luồng thực thi mở rộng}
\ucstep{C3a}{Quản trị viên}{Chọn “Ẩn”: công thức không còn xuất hiện trong gợi ý; các bản sao người dùng đã lưu không bị ảnh hưởng.}
\ucstep{C4a}{Hệ thống}{Dữ liệu không hợp lệ: hiển thị \msg{MSG01} hoặc \msg{MSG02}.}
\ucfield{Điều kiện sau}{\textbf{A.} Danh mục loại được cập nhật, có hiệu lực cho bộ lọc và biểu mẫu nhập liệu. \textbf{B.} Danh mục đơn vị và hệ số quy đổi được cập nhật. \textbf{C.} Bộ công thức mẫu được cập nhật, công thức Công bố tham gia \uc{UC18} cho mọi người dùng.}
\ucfield{Quy tắc nghiệp vụ}{\br{BR16}, \br{BR17}, \br{BR22}}
\end{ucspec}



% ---------------------------------------------------------------------
%  PHỤ LỤC C. DANH MỤC THÔNG ĐIỆP HỆ THỐNG
% ---------------------------------------------------------------------
\section{Danh mục thông điệp hệ thống}\label{pl:thong-diep}

Bảng~\ref{tab:thong-diep} liệt kê các thông điệp hiển thị cho người dùng. Mỗi thông điệp có mã \msg{MSGxx} để dẫn chiếu trong đặc tả use case. Việc chuẩn hóa giúp giao diện nhất quán, hỗ trợ bản địa hóa theo \nfr{NFR07} và thuận tiện kiểm thử. Thông điệp gồm bốn loại: \emph{Lỗi} khi không thể thực hiện thao tác, \emph{Cảnh báo} khi thao tác có thể gây hệ quả cần lưu ý, \emph{Xác nhận} khi cần người dùng đồng ý và \emph{Thông tin} để thông báo kết quả. Nội dung trong ngoặc nhọn là tham số được thay bằng giá trị cụ thể khi hiển thị.

{\small
\begin{longtable}{|C{1.4cm}|L{\dimexpr\textwidth-8\tabcolsep-5\arrayrulewidth-1.4cm-1.9cm-3.6cm\relax}|C{1.9cm}|L{3.6cm}|}
\caption{Danh mục thông điệp hệ thống}\label{tab:thong-diep}\\
\hline
\hd{Mã} & \hd{Nội dung thông điệp} & \hd{Loại} & \hd{Use case sử dụng}\\ \hline
\endfirsthead
\multicolumn{4}{l}{\footnotesize\itshape Bảng \thetable\ (tiếp theo)}\\ \hline
\hd{Mã} & \hd{Nội dung thông điệp} & \hd{Loại} & \hd{Use case sử dụng}\\ \hline
\endhead
MSG01 & Vui lòng nhập đầy đủ các trường bắt buộc. & Lỗi & UC01, UC03, UC25, UC32, UC49, UC52\\ \hline
MSG02 & \{Tên trường\} không hợp lệ: \{yêu cầu cụ thể\}. & Lỗi & UC01, UC04, UC06, UC07, UC09, UC16, UC18, UC32, UC49–UC52\\ \hline
MSG03 & Email này đã được sử dụng. Bạn có thể đăng nhập hoặc khôi phục mật khẩu. & Lỗi & UC01\\ \hline
MSG04 & Mã xác thực không đúng hoặc đã hết hiệu lực. & Lỗi & UC02\\ \hline
MSG05 & Email hoặc mật khẩu không đúng. & Lỗi & UC03\\ \hline
MSG06 & Bạn đã đăng nhập sai quá nhiều lần. Vui lòng thử lại sau \{n\} phút. & Lỗi & UC03\\ \hline
MSG07 & Mã mời không tồn tại, đã hết hạn hoặc đã bị thu hồi. & Lỗi & UC11\\ \hline
MSG08 & Bạn đang là thành viên của nhóm \{tên nhóm\}. Vui lòng rời nhóm hiện tại trước khi tham gia hoặc tạo nhóm mới. & Lỗi & UC09, UC11\\ \hline
MSG09 & Nhóm đã đạt số thành viên tối đa (\{n\} người). & Lỗi & UC10, UC11\\ \hline
MSG10 & Đã lưu thay đổi. & Thông tin & Dùng chung cho các thao tác thành công\\ \hline
MSG11 & “\{tên mặt hàng\}” đã có trong danh sách, số lượng đã được cộng dồn thành \{số lượng\}. & Thông tin & UC17\\ \hline
MSG12 & Mặt hàng này đang được giao cho \{tên thành viên\}. Bạn có muốn đánh dấu mua thay không? & Xác nhận & UC20\\ \hline
MSG13 & Bạn không có quyền thực hiện thao tác này. & Lỗi & UC13, UC19\\ \hline
MSG14 & Ngày hết hạn không được trước ngày mua. & Lỗi & UC21, UC25, UC27\\ \hline
MSG15 & Tủ lạnh đã có đủ nguyên liệu cho các bữa ăn đã chọn. & Thông tin & UC22\\ \hline
MSG16 & Số lượng phải là số dương, tối đa hai chữ số thập phân. & Lỗi & UC17, UC25, UC27\\ \hline
MSG17 & Lượng nhập vượt quá lượng còn lại (\{lượng còn lại\}). & Lỗi & UC28, UC29\\ \hline
MSG18 & Không tìm thấy kết quả phù hợp. & Thông tin & UC42, UC43\\ \hline
MSG19 & Bạn có chắc chắn muốn \{thao tác\} \{đối tượng\}? \{Hệ quả nếu có\}. & Xác nhận & UC10, UC13, UC15, UC33, UC48, UC49\\ \hline
MSG20 & Không có kết nối mạng. Thay đổi đã được lưu trên thiết bị và sẽ đồng bộ khi có mạng. & Cảnh báo & UC16 và các thao tác hỗ trợ ngoại tuyến\\ \hline
MSG21 & Tài khoản của bạn đã bị khóa. Vui lòng liên hệ quản trị viên. & Lỗi & UC03\\ \hline
MSG22 & Dữ liệu danh mục đang được sử dụng nên không thể xóa, chỉ có thể chuyển sang Ngừng sử dụng. & Cảnh báo & UC50, UC51\\ \hline
MSG23 & Không tìm thấy sản phẩm ứng với mã vạch này. Vui lòng nhập thông tin thủ công. & Thông tin & UC26\\ \hline
MSG24 & Bạn là trưởng nhóm. Hãy chuyển quyền cho một thành viên khác trước khi rời nhóm. & Cảnh báo & UC15\\ \hline
MSG25 & Khoảng thời gian không hợp lệ. Vui lòng chọn tối đa 12 tháng. & Lỗi & UC44, UC45, UC46\\ \hline
MSG26 & Chưa có dữ liệu trong khoảng thời gian đã chọn. & Thông tin & UC44, UC45\\ \hline
MSG27 & Giá trị của \{tham số\} phải nằm trong khoảng \{miền hợp lệ\}. & Lỗi & UC53\\ \hline
MSG28 & Chỉ có thể lập kế hoạch từ hôm nay đến \{n\} ngày tới. & Lỗi & UC37, UC39\\ \hline
MSG29 & Mật khẩu xác nhận không khớp. & Lỗi & UC01, UC04, UC07\\ \hline
MSG30 & Vui lòng chờ \{n\} giây trước khi yêu cầu gửi lại mã. & Thông tin & UC02\\ \hline
MSG31 & Chưa có món nào phù hợp với thực phẩm hiện có. Dưới đây là các món gần đủ nguyên liệu nhất. & Thông tin & UC40\\ \hline
MSG32 & Mặt hàng đã được cất vào tủ lạnh nên không thể sửa hoặc hoàn tác. & Lỗi & UC17\\ \hline
MSG33 & Không thể kết nối tới máy chủ. Vui lòng kiểm tra kết nối mạng và thử lại. & Lỗi & UC01 và các thao tác bắt buộc trực tuyến\\ \hline
MSG34 & Tài khoản chưa được xác thực. Vui lòng nhập mã xác thực đã được gửi tới \{email\} để kích hoạt tài khoản. & Lỗi & UC02, UC03\\ \hline
MSG35 & Dữ liệu vừa được \{tên\} thay đổi. Vui lòng xem lại trước khi tiếp tục. & Cảnh báo & UC17, UC19, UC20, UC21, UC27, UC41 và các thao tác đồng bộ ngoại tuyến\\ \hline
MSG36 & Bạn thao tác quá nhanh. Vui lòng thử lại sau \{n\} giây. & Lỗi & UC01, UC03, UC11, UC26, UC42 và các thao tác bị giới hạn tần suất\\ \hline
\end{longtable}
}

% ---------------------------------------------------------------------
%  PHỤ LỤC D. MA TRẬN TRUY VẾT YÊU CẦU
% ---------------------------------------------------------------------
\section{Ma trận truy vết yêu cầu}\label{pl:truy-vet}

Ma trận truy vết tại Bảng~\ref{tab:truy-vet} liên kết yêu cầu trong đề bài~\cite{debai} và chức năng mở rộng ở Bảng~\ref{tab:chuc-nang-mo-rong} với nghiệp vụ con, use case, quy tắc nghiệp vụ và yêu cầu phi chức năng tương ứng. Yêu cầu đề bài được mã hóa theo thứ tự xuất hiện (YC); ba đặc điểm nổi bật được ký hiệu ĐĐ. Ma trận hỗ trợ kiểm tra theo hai chiều: mỗi yêu cầu phải được hiện thực bởi ít nhất một use case và mỗi use case phải truy ngược được về ít nhất một yêu cầu.

{\small
\begin{longtable}{|C{1.2cm}|L{\dimexpr\textwidth-12\tabcolsep-7\arrayrulewidth-1.2cm-2.0cm-2.6cm-1.9cm-1.7cm\relax}|L{2.0cm}|L{2.6cm}|L{1.9cm}|L{1.7cm}|}
\caption{Ma trận truy vết từ yêu cầu đến nghiệp vụ và use case}\label{tab:truy-vet}\\
\hline
\hd{Mã} & \hd{Yêu cầu} & \hd{Nghiệp vụ} & \hd{Use case} & \hd{Quy tắc} & \hd{NFR}\\ \hline
\endfirsthead
\multicolumn{6}{l}{\footnotesize\itshape Bảng \thetable\ (tiếp theo)}\\ \hline
\hd{Mã} & \hd{Yêu cầu} & \hd{Nghiệp vụ} & \hd{Use case} & \hd{Quy tắc} & \hd{NFR}\\ \hline
\endhead
\multicolumn{6}{|l|}{\hs{Yêu cầu của đề bài}}\\ \hline
YC00 & Định danh người dùng làm nền tảng cho chia sẻ và phân quyền (yêu cầu ngầm định) & NV01.1–NV01.4 & UC01--UC04 & BR01–BR03 & NFR14, NFR15, NFR17\\ \hline
YC01 & Tạo danh sách mua sắm theo từng ngày, chỉnh sửa, cập nhật và lưu lại cho những lần sau & NV02.1–NV02.3 & UC08, UC08, UC08, UC08 & BR08, BR09 & NFR05, NFR11\\ \hline
YC02 & Tạo nhóm gia đình, chia sẻ danh sách, mọi thành viên cùng đóng góp và cập nhật & NV01.5–NV01.7, NV02.2 & UC05--UC07, UC08 & BR04–BR07, BR23 & NFR10, NFR16\\ \hline
YC03 & Trưởng nhóm điều phối, phân công mua sắm, bảo đảm không quên hoặc trùng lặp & NV02.4, NV02.5 & UC08, UC09 & BR09, BR10 & NFR10\\ \hline
YC04 & Lên kế hoạch bữa sáng, trưa, tối theo ngày hoặc tuần & NV05.1, NV05.2 & UC17, UC17, UC17 & BR18 & NFR05\\ \hline
YC05 & Đề xuất món ăn dựa trên thực phẩm có sẵn trong tủ lạnh & NV05.3 & UC18 & BR19 & NFR09\\ \hline
YC06 & Lưu thông tin thực phẩm: tên, số lượng, ngày hết hạn, vị trí lưu trữ; theo dõi và quản lý & NV03.1–NV03.4 & UC12, UC11, UC13, UC13, UC13 & BR11, BR12, BR14, BR15 & NFR04\\ \hline
YC07 & Tự động nhắc khi thực phẩm sắp hết hạn trước 3 ngày & NV03.5 & UC04, UC12, UC14 & BR12, BR13 & NFR12\\ \hline
YC08 & Lưu, quản lý công thức yêu thích, liên kết nguyên liệu với tủ lạnh & NV04.1, NV04.2, NV04.4 & UC15, UC15, UC16, UC15 & BR16, BR17 & —\\ \hline
YC09 & Tìm kiếm món ăn, thực phẩm trong hệ thống hoặc trong tủ lạnh & NV06.1 & UC20 & BR23 & NFR09\\ \hline
YC10 & Phân loại thực phẩm theo danh mục như thịt, rau củ, đồ hộp & NV06.2, NV08.3 & UC20, UC23 & BR22 & —\\ \hline
YC11 & Gợi ý món ăn phù hợp với thực phẩm sẵn có & NV06.3 & UC12, UC18 & BR19 & NFR09\\ \hline
YC12 & Báo cáo đồ đã mua và thực phẩm đã tiêu thụ trong một khoảng thời gian & NV07.1, NV07.2 & UC21, UC21 & BR21 & —\\ \hline
YC13 & Quản trị viên quản lý tài khoản người dùng & NV08.1 & UC24 & BR23 & NFR16\\ \hline
YC14 & Quản trị viên quản lý danh mục tên món đồ, loại mặt hàng, đơn vị đo và cấu hình hệ thống & NV08.2–NV08.5 & UC23, UC25 & BR16, BR22 & NFR20\\ \hline
\multicolumn{6}{|l|}{\hs{Đặc điểm nổi bật của đề bài}}\\ \hline
ĐĐ01 & Tính tiện lợi cao, kiểm soát toàn bộ quy trình mua sắm và nấu nướng & NV01–NV08 & Toàn bộ & — & NFR05, NFR08, NFR11\\ \hline
ĐĐ02 & Chia sẻ trong gia đình, phân công giữa các thành viên & NV01.5–NV01.7, NV02.4 & UC05--UC07, UC08 & BR04–BR07, BR10 & NFR10, NFR16\\ \hline
ĐĐ03 & Thông báo nhắc nhở thông minh, chọn món từ nguyên liệu sẵn có & NV03.5, NV05.3 & UC04, UC14, UC18 & BR13, BR19 & NFR12\\ \hline
\multicolumn{6}{|l|}{\hs{Chức năng mở rộng đề xuất}}\\ \hline
MR01 & Mời thành viên bằng mã mời hoặc mã QR & NV01.6 & UC06, UC06 & BR06 & NFR17\\ \hline
MR02 & Chuyển mặt hàng đã mua vào tủ lạnh & NV02.6 & UC11 & BR11 & NFR13\\ \hline
MR03 & Sinh danh sách mua sắm từ kế hoạch bữa ăn & NV02.7 & UC10 & BR20 & —\\ \hline
MR04 & Quét mã vạch sản phẩm (định hướng phát triển, không cài đặt) & NV03.1 & — & BR22 & NFR18\\ \hline
MR05 & Gợi ý món có chấm điểm, ưu tiên thực phẩm sắp hết hạn & NV05.3 & UC18 & BR19 & NFR09\\ \hline
MR06 & Xác nhận đã nấu, trừ kho theo FEFO & NV05.4 & UC19 & BR15 & NFR13\\ \hline
MR07 & Bổ sung nguyên liệu còn thiếu từ công thức & NV04.3 & UC10 & BR09 & —\\ \hline
MR08 & Báo cáo chi tiêu, lãng phí, xuất báo cáo & NV07.1, NV07.3, NV07.4 & UC21, UC21, UC22 & BR21 & —\\ \hline
MR09 & Thiết lập nhắc nhở cá nhân hóa & NV01.4 & UC04 & BR12, BR13 & —\\ \hline
MR10 & Đồng bộ thời gian thực và hoạt động ngoại tuyến (định hướng phát triển) & NV02.2, NV02.5 & UC08, UC08, UC09 & — & NFR10, NFR11\\ \hline
\end{longtable}
}

Kiểm tra ngược cho thấy cả 25 use case đều có trong ma trận. Trong số đó, 20 use case hiện thực các yêu cầu YC00–YC14; năm use case còn lại — \uc{UC10}, \uc{UC11}, \uc{UC19}, \uc{UC21} và \uc{UC22} — phát sinh từ chức năng mở rộng. Các chức năng \uc{UC04}, \uc{UC06}, \uc{UC18} và \uc{UC21} vừa đáp ứng yêu cầu đề bài, vừa được mở rộng tương ứng. Kết quả truy vết không ghi nhận use case nằm ngoài phạm vi hoặc yêu cầu nào chưa được hiện thực.

Trong ma trận, \nfr{NFR10} và \nfr{NFR11} xuất hiện ở các dòng YC02, YC03, ĐĐ01 và MR10 nhưng thuộc định hướng phát triển, không phải tiêu chí chấp nhận. Tầng ưu tiên (1, 2, 3) của từng use case được ghi trong Bảng~\ref{tab:danh-sach-uc}. Các yêu cầu YC và đặc điểm đề bài được hiện thực chủ yếu bởi use case tầng 2 và 3, còn chức năng mở rộng MR01--MR09 bởi use case tầng 1.

\begingroup
\usetalldiagrams
% ---------------------------------------------------------------------
%  PHỤ LỤC. HIỆN THỰC HÓA USE CASE (chi tiết của mục Hiện thực hóa các use
%  case trọng tâm trong Chương 3)
% ---------------------------------------------------------------------
\section{Hiện thực hóa use case}\label{pl:hien-thuc-uc}

Phụ lục này là phần chi tiết của mục~\ref{ssec:hien-thuc-uc}, theo cùng khuôn hiện thực đã nêu ở đó: mỗi use case do một dịch vụ ứng dụng của module sở hữu đảm nhận, chạy trong một giao dịch theo thứ tự khóa toàn cục và phát sự kiện miền cho các tác dụng phụ.

Bảng~\ref{tab:c3-uc-thiet-ke} là ma trận truy vết 53 use case, gom theo chín gói như Bảng~\ref{tab:danh-sach-uc}. Cột Module ghi module sở hữu dịch vụ, trong ngoặc là module tham gia cùng giao dịch hoặc được đọc qua cổng. Cột API dùng mã của mục~\ref{sec:api}, trong ngoặc là API hỗ trợ cùng luồng. Tên bảng được ghi không kèm schema vì là duy nhất trong toàn CSDL; schema của bảng suy ra từ module sở hữu nêu ở cột Module. Ở cột Sự kiện, mũi tên chỉ bên nhận, tên trong ngoặc là hàng đợi River, còn “đồng bộ” nghĩa là ghi \texttt{sync.space\_changes} và phát \texttt{space:\{sid\}} sau commit. Cột Sơ đồ ghi số hình biểu đồ tuần tự hoặc máy trạng thái trong chương này. Chỉ năm API không gắn với một use case cụ thể: A05 làm mới phiên cho mọi lời gọi, A16 xóa tài khoản theo \nfr{NFR18}, Y01 và Y02 phục vụ đồng bộ và thời gian thực (\nfr{NFR10}, \nfr{NFR11}), G01 liệt kê các không gian để ứng dụng chuyển không gian làm việc.

{\footnotesize
\begin{longtable}{|C{1.0cm}|L{2.1cm}|L{1.9cm}|L{4.0cm}|L{3.4cm}|C{\dimexpr\textwidth-12\tabcolsep-7\arrayrulewidth-1.0cm-2.1cm-1.9cm-4.0cm-3.4cm\relax}|}
\caption{Ma trận truy vết từ use case tới thiết kế backend}\label{tab:c3-uc-thiet-ke}\\
\hline
\hd{Mã UC} & \hd{Module} & \hd{API} & \hd{Bảng dữ liệu chính} & \hd{Sự kiện, tác vụ nền} & \hd{Sơ đồ}\\ \hline
\endfirsthead
\multicolumn{6}{l}{\footnotesize\itshape Bảng \thetable\ (tiếp theo)}\\ \hline
\hd{Mã UC} & \hd{Module} & \hd{API} & \hd{Bảng dữ liệu chính} & \hd{Sự kiện, tác vụ nền} & \hd{Sơ đồ}\\ \hline
\endhead
\multicolumn{6}{|l|}{\hs{Gói 1. Xác thực và quản lý tài khoản}}\\ \hline
UC01 & identity \newline (+~household) & A01, A02 (A03) & \path{users}, \path{otp_challenges}; \path{spaces}, \path{space_members} & \path{user.registered} $\to$ thư OTP (\path{notify}) & \ref{fig:seq-dang-ky-otp}, \ref{fig:trang-thai-tai-khoan}\\ \hline
UC02 & identity & A02, A03 & \path{otp_challenges}; Redis \path{otp:cooldown}, \path{otp:hourly} & gửi lại thư OTP (\path{notify}) & \ref{fig:seq-dang-ky-otp}\\ \hline
UC03 & identity & A04 (A14) & \path{users}, \path{refresh_sessions}, \path{devices} & — & \ref{fig:seq-dang-nhap-lam-moi}, \ref{fig:seq-dung-lai-ma-lam-moi}\\ \hline
UC04 & identity & A07, A02, A08 & \path{otp_challenges}, \path{users}, \path{refresh_sessions} & \path{password.reset_requested} $\to$ thư OTP & \ref{fig:seq-dang-ky-otp}\\ \hline
UC05 & identity & A06, A15 & \path{refresh_sessions}, \path{devices}; Redis \path{auth:ver} & — & —\\ \hline
UC06 & identity & A09, A10 (Y03, Y04) & \path{users}, \path{media.objects}; \path{space_changes} & \path{media.uploaded} $\to$ xử lý ảnh (\path{media}); đồng bộ mọi không gian & —\\ \hline
UC07 & identity & A11 & \path{users}, \path{refresh_sessions}; Redis \path{auth:ver} & — & —\\ \hline
UC08 & identity & A12, A13 & \path{user_settings}; Redis \path{cfg:all}, \path{suggest:{sid}} & đổi $N$ $\to$ xóa \path{suggest:{sid}} & —\\ \hline
\multicolumn{6}{|l|}{\hs{Gói 2. Quản lý nhóm gia đình}}\\ \hline
UC09 & household & G02 (G04) & \path{spaces}, \path{space_members} & \path{CopySpaceData} (\path{default}); đồng bộ & —\\ \hline
UC10 & household & G05, G06 & \path{invite_codes}, \path{spaces}; Redis \path{cfg:all} & — & —\\ \hline
UC11 & household & G07 & \path{invite_codes}, \path{spaces}, \path{space_members}; \path{users} & \path{member.joined} $\to$ trưởng nhóm (\path{notify}); đồng bộ & \ref{fig:seq-tham-gia-nhom}\\ \hline
UC12 & household & G03 & \path{spaces}, \path{space_members}; \path{users}, \path{space_changes} & — & —\\ \hline
UC13 & household \newline (+~shopping) & G08 & \path{space_members}; \path{shopping_items} & \path{member.removed} $\to$ thông báo; đồng bộ & \ref{fig:seq-roi-nhom-chuyen-quyen}\\ \hline
UC14 & household & G09 & \path{space_members}, \path{invite_codes} & \path{leader.transferred} $\to$ thông báo; đồng bộ & \ref{fig:seq-roi-nhom-chuyen-quyen}\\ \hline
UC15 & household \newline (+~shopping) & G10 & \path{space_members}, \path{spaces}; \path{shopping_items} & \path{member.left} $\to$ thông báo; \path{PurgeSpace} khi giải tán; đồng bộ & \ref{fig:seq-roi-nhom-chuyen-quyen}\\ \hline
\multicolumn{6}{|l|}{\hs{Gói 3. Quản lý danh sách mua sắm}}\\ \hline
UC16 & shopping & S02, S04 (S01, S03, S05) & \path{shopping_lists}, \path{shopping_items} & đồng bộ & —\\ \hline
UC17 & shopping & S07, S08, S09 & \path{shopping_items}, \path{shopping_lists} & đồng bộ & \ref{fig:trang-thai-mat-hang-be}\\ \hline
UC18 & shopping & S06 & \path{shopping_lists}, \path{shopping_items} & đồng bộ & —\\ \hline
UC19 & shopping & S10, S11 & \path{shopping_items}; \path{space_members} & \path{AssignmentDigest} $\to$ \path{item.assigned}, \path{item.unassigned} $\to$ FCM (\path{notify}); đồng bộ & \ref{fig:seq-phan-cong}\\ \hline
UC20 & shopping & S12, S13 & \path{shopping_items}, \path{shopping_lists} & \path{items.purchased} $\to$ đồng bộ, báo cáo & \ref{fig:seq-outbox-mua-hang}\\ \hline
UC21 & shopping \newline (+~pantry) & S14 & \path{shopping_items}; \path{food_lots}, \path{stock_movements} & \path{items.stocked}, \path{lot.changed} $\to$ đồng bộ & \ref{fig:seq-nhap-tu}\\ \hline
UC22 & mealplan \newline (+~recipe, pantry, shopping) & S15 (S01) & \path{meal_entries}; \path{recipe_ingredients}, \path{food_lots}, \path{shopping_lists}, \path{shopping_items} & đồng bộ khi tạo danh sách & \ref{fig:seq-sinh-ds-ke-hoach}\\ \hline
UC23 & shopping & S01, S03 & \path{shopping_lists}, \path{shopping_items} & — & —\\ \hline
\multicolumn{6}{|l|}{\hs{Gói 4. Quản lý thực phẩm trong tủ lạnh}}\\ \hline
UC24 & pantry & P01, P03 & \path{food_lots}; \path{user_settings} & — & —\\ \hline
UC25 & pantry \newline (+~catalog) & P02 (Q02, Q03, Y03, Y04) & \path{food_lots}, \path{stock_movements}; \path{foods}, \path{food_shelf_lives}; \path{media.objects} & \path{lot.changed} $\to$ đồng bộ; \path{food.changed} $\to$ ES (\path{index}) & \ref{fig:seq-them-thuc-pham}\\ \hline
UC26 & catalog & P08 & \path{foods}; Redis \path{barcode:{code}} & gọi Open Food Facts có ngắt mạch & \ref{fig:seq-ma-vach}\\ \hline
UC27 & pantry & P04 & \path{food_lots}, \path{stock_movements} & \path{lot.changed} $\to$ đồng bộ & —\\ \hline
UC28 & pantry & P05 & \path{food_lots}, \path{stock_movements} & \path{lot.changed} $\to$ đồng bộ & \ref{fig:seq-xac-nhan-nau}\\ \hline
UC29 & pantry & P06 & \path{food_lots}, \path{stock_movements} & \path{lot.changed} $\to$ đồng bộ & —\\ \hline
UC30 & pantry & P07 (M05) & \path{food_lots}; \path{user_settings} & — & \ref{fig:seq-goi-y-mon}\\ \hline
UC31 & notifi\-cation & (A14, A17, A18) & \path{notifications}, \path{deliveries}; \path{food_lots}, \path{user_settings}, \path{devices} & tác vụ định kỳ mỗi phút $\to$ \path{UserReminder} $\to$ \path{PushNotification} (\path{notify}) & \ref{fig:seq-nhac-han-dung-be}\\ \hline
\multicolumn{6}{|l|}{\hs{Gói 5. Quản lý công thức nấu ăn}}\\ \hline
UC32 & recipe & R02, R09 (Y03, Y04) & \path{recipes}, \path{recipe_ingredients} & \path{recipe.changed} $\to$ ES (\path{index}) & —\\ \hline
UC33 & recipe \newline (+~mealplan) & R04, R05 & \path{recipes}, \path{recipe_ingredients}; \path{meal_entries} & \path{recipe.changed} $\to$ ES; đồng bộ & —\\ \hline
UC34 & recipe \newline (+~pantry) & R03 (R01) & \path{recipes}, \path{recipe_ingredients}; \path{food_lots} & — & —\\ \hline
UC35 & recipe & R06, R07 & \path{favorite_recipes} & xóa \path{suggest:{sid}} & —\\ \hline
UC36 & recipe \newline (+~pantry, shopping) & R08 & \path{recipe_ingredients}; \path{food_lots}, \path{shopping_items} & đồng bộ & —\\ \hline
\multicolumn{6}{|l|}{\hs{Gói 6. Quản lý kế hoạch bữa ăn}}\\ \hline
UC37 & mealplan \newline (+~recipe) & M02 (R01) & \path{meal_entries}; \path{recipes} & đồng bộ & \ref{fig:seq-lap-lich}\\ \hline
UC38 & mealplan & M01 & \path{meal_entries} & — & —\\ \hline
UC39 & mealplan & M03, M04 & \path{meal_entries} & đồng bộ & —\\ \hline
UC40 & mealplan \newline (+~recipe, pantry) & M05 & \path{recipes}, \path{favorite_recipes}, \path{food_lots}; Redis \path{suggest:{sid}}, \path{cfg:all} & bộ đệm xóa khi có \path{lot.changed}, \path{recipe.changed}, \path{config.changed} & \ref{fig:seq-goi-y-mon}\\ \hline
UC41 & mealplan \newline (+~pantry) & M06 & \path{meal_entries}; \path{food_lots}, \path{stock_movements} & \path{meal.cooked}, \path{lot.changed} $\to$ đồng bộ & \ref{fig:seq-xac-nhan-nau}\\ \hline
\multicolumn{6}{|l|}{\hs{Gói 7. Tìm kiếm và phân loại thực phẩm}}\\ \hline
UC42 & search \newline (+~pantry) & Q01, Q02 & ES \path{foods}, \path{recipes}; \path{food_lots} (\path{StockQuery.SearchLots}) & chỉ mục nhận \path{food.changed}, \path{recipe.changed} (\path{index}) & \ref{fig:seq-tim-kiem-chi-muc}\\ \hline
UC43 & search \newline (+~pantry) & Q01, P01 (Q03) & ES \path{foods}; \path{food_categories}, \path{food_lots} & — & \ref{fig:seq-tim-kiem-chi-muc}\\ \hline
\multicolumn{6}{|l|}{\hs{Gói 8. Thống kê báo cáo}}\\ \hline
UC44 & report & B01 & \path{daily_space_stats}; \path{shopping_items}, \path{foods}, \path{food_categories} & làm mới hằng đêm & —\\ \hline
UC45 & report & B02 & \path{daily_space_stats}; \path{stock_movements}, \path{meal_entries} & làm mới hằng đêm & —\\ \hline
UC46 & report & B03 & \path{daily_space_stats}; \path{stock_movements}, \path{food_lots} & làm mới hằng đêm & \ref{fig:seq-bao-cao-xuat}\\ \hline
UC47 & report & B04, B05 & \path{report.exports}; MinIO \path{exports} & \path{ExportReport} (\path{report}) $\to$ \path{export.completed} $\to$ FCM (\path{notify}) & \ref{fig:seq-bao-cao-xuat}\\ \hline
\multicolumn{6}{|l|}{\hs{Gói 9. Quản trị hệ thống}}\\ \hline
UC48 & admin \newline (+~identity) & D01, D02 & \path{audit_logs}; \path{users}, \path{refresh_sessions}; Redis \path{auth:ver} & PUBLISH \path{user:{uid}}; thư thông báo (\path{notify}) & \ref{fig:trang-thai-tai-khoan}, \ref{fig:seq-quan-tri-cau-hinh}\\ \hline
UC49 & catalog & D03, D04, D11 & \path{foods}, \path{food_shelf_lives}; \path{audit_logs} & \path{food.changed} $\to$ ES (\path{index}); \path{food.merged} $\to$ đổi \path{food_id} ở pantry, shopping, recipe & \ref{fig:seq-duyet-thuc-pham}\\ \hline
UC50 & catalog & D05 & \path{food_categories}; \path{audit_logs} & đổi ETag danh mục & —\\ \hline
UC51 & catalog & D06 & \path{units}; \path{audit_logs} & đổi ETag danh mục & —\\ \hline
UC52 & recipe & D07 & \path{recipes}, \path{recipe_ingredients}; \path{audit_logs} & \path{recipe.changed} $\to$ ES (\path{index}) & —\\ \hline
UC53 & admin & D08, D09 (D10) & \path{system_configs}, \path{audit_logs}; Redis \path{cfg:all} & \path{config.changed} $\to$ Pub/Sub \path{cfg}, xóa \path{suggest:*} & \ref{fig:seq-quan-tri-cau-hinh}\\ \hline
\end{longtable}
}

\InputIfFileExists{chapters/c3/ht-dang-ky.tex}{}{}

% =====================================================================
\subsection{Tham gia nhóm gia đình (UC11)}\label{sssec:ht-tham-gia-nhom}

\uc{UC11} do dịch vụ \texttt{household.Service.JoinSpace} hiện thực qua G07 \texttt{POST} \path{/spaces/join}. Luồng gồm hai lời gọi, ứng với phần kiểm tra và xem trước (bước 4--6) và phần gia nhập sau khi người dùng xác nhận (bước 7--8) trong Bảng~\ref{tab:UC11}; Hình~\ref{fig:seq-tham-gia-nhom} thể hiện cộng tác của cả hai. Mã mời đến từ ô nhập, mã QR hoặc liên kết mời \path{https://dicho.example/join/{code}}; liên kết này được khai báo là Universal Links trên iOS và App Links trên Android nên hệ điều hành mở thẳng ứng dụng với mã điền sẵn (bước 3b).

\diagram[\textwidth]{c3-29-seq-tham-gia-nhom}{Biểu đồ tuần tự hiện thực use case Tham gia nhóm gia đình}{fig:seq-tham-gia-nhom}

Lời gọi thứ nhất mang trường \texttt{dryRun}: dịch vụ chỉ đọc, tra mã bằng một phép nối \texttt{invite\_codes} với \texttt{spaces} theo \texttt{code}, kiểm tra theo đúng thứ tự của Hình~\ref{fig:act-tham-gia-nhom} là \br{BR06}, \br{BR04} rồi \br{BR05} và trả tên nhóm, trưởng nhóm, số thành viên. Lời gọi xác nhận kiểm tra lại toàn bộ trong giao dịch ghi, vì giữa hai lần gọi mã có thể bị thu hồi hoặc nhóm có thể đã đủ người. Tình huống khó nhất là hai người nhập mã cùng lúc khi nhóm còn một chỗ: nếu mỗi giao dịch chỉ đếm rồi chèn, cả hai cùng thấy 9 người và nhóm thành 11 thành viên. Mức cô lập \texttt{READ COMMITTED} không ngăn được hiện tượng lệch ghi (write skew) này~\cite{kleppmann}. Sau khóa lũy đẳng, dịch vụ tra mã để biết \texttt{space\_id} rồi lấy \path{pg_advisory_xact_lock(hashtext(space_id::text))}, một khóa tư vấn gắn với giao dịch (bước 12 của hình). Giao dịch đến sau phải chờ (tối đa bằng \texttt{lock\_timeout}), rồi đọc lại mã và số thành viên bằng câu lệnh mới, thấy người vừa gia nhập và nhận \msg{MSG09}. G05--G10 lấy cùng khóa này nên mọi thay đổi thành viên, vai trò và mã mời của một nhóm được tuần tự hóa. Nhóm không khóa dòng \texttt{household.spaces} bằng \texttt{FOR UPDATE}, vì dòng này giữ bộ đếm \texttt{change\_seq} mà mọi giao dịch ghi của không gian đều cập nhật ở bước cuối; khóa nó từ đầu sẽ xếp hàng cả việc thêm mặt hàng hay nhập lô phía sau việc gia nhập và phá thứ tự khóa toàn cục. Hai không gian trùng giá trị băm chỉ phải chờ nhau một cách không cần thiết, không làm sai kết quả. Khóa tư vấn cũng không chặn được việc một người dùng hai thiết bị vào hai nhóm khác nhau, vì hai nhóm có hai khóa khác nhau; hàng rào cho \br{BR04} là chỉ mục \texttt{uq\_member\_family}, tức \texttt{UNIQUE(user\_id) WHERE space\_kind = 'family'}, với lỗi vi phạm 23505 được ánh xạ thành 422 mã 4004 kèm \msg{MSG08}. Mã mời chỉ hợp lệ khi \texttt{status = 'active'} và \texttt{expires\_at > now()}, không phụ thuộc tác vụ đánh dấu hết hạn hằng đêm. Giới hạn 10 lần mỗi giờ cho mỗi người (\texttt{rl:join:\{uid\}}) áp dụng cả cho lời gọi chạy thử, vì đó cũng là cách dò mã; với $32^8 \approx 1{,}1 \cdot 10^{12}$ tổ hợp, việc dò là không khả thi.

Cùng giao dịch còn cập nhật \texttt{last\_space\_id} qua cổng của identity để chuyển không gian làm việc (bước 8 của use case), chèn tác vụ \texttt{PushNotification} cho sự kiện \texttt{member.joined} và kết thúc bằng \texttt{space\_changes}. Sau commit, dịch vụ xóa khóa đệm \texttt{member:\{sid\}:\{uid\}} và phát \texttt{space:\{sid\}} để danh sách thành viên trên thiết bị khác cập nhật trong 3 giây (\nfr{NFR10}); worker báo cho trưởng nhóm theo cơ chế outbox của Hình~\ref{fig:seq-outbox-mua-hang} nếu người này bật thông báo hoạt động nhóm (\uc{UC08}). Các nhánh lỗi trả 422 mã 4004 kèm \br{BR06} (\msg{MSG07}), \br{BR04} (\msg{MSG08}) hoặc \br{BR05} (\msg{MSG09}); vượt giới hạn tần suất trả 429 mã 4290 kèm \msg{MSG36}. Nhờ khóa lũy đẳng, nếu phản hồi thành công bị mất, lần gửi lại nhận đúng phản hồi 200 đã lưu thay vì lỗi “đã là thành viên”; nếu lần gửi lại đến khi lần đầu chưa xong, ứng dụng nhận 409 mã 4002 và thử lại sau.

\InputIfFileExists{chapters/c3/ht-roi-nhom.tex}{}{}

% =====================================================================
\subsection{Phân công người mua (UC19)}\label{sssec:ht-phan-cong}

Phương thức \texttt{ShoppingService.AssignItem} hiện thực \uc{UC19} qua S10 \texttt{PUT} \path{.../shopping-items/{itemId}/assignee} để giao và S11 \texttt{DELETE} cùng đường dẫn để bỏ phân công (bước 3a), khi đó người được giao là rỗng. Hình~\ref{fig:seq-phan-cong} thể hiện vòng lặp phân công nhiều mặt hàng và nhánh bỏ phân công.

\diagram[\textwidth]{c3-30-seq-phan-cong}{Biểu đồ tuần tự hiện thực use case Phân công người mua}{fig:seq-phan-cong}

Mỗi lời gọi áp dụng cho một mặt hàng với phiên bản riêng trong \texttt{If-Match}, nhất quán với cơ chế đồng bộ theo thực thể (Hình~\ref{fig:seq-dong-bo-ngoai-tuyen}); chọn nhiều mặt hàng (bước 1) tương ứng một chuỗi lời gọi có kết quả độc lập, đúng yêu cầu ở bước 4b. Số việc của từng thành viên (bước 2) tính trên bản sao cục bộ. Quyền trưởng nhóm được kiểm tra hai lớp. Chính sách ở tầng ứng dụng đọc vai trò từ khóa đệm \texttt{member:\{sid\}:\{uid\}} để từ chối nhanh. Bộ đệm có thể chưa kịp phản ánh nhánh 4a, khi người dùng vừa chuyển quyền trên thiết bị khác, nên sau khóa lũy đẳng dịch vụ đọc lại dòng \texttt{space\_members} của người gọi và người được giao với \texttt{FOR SHARE}, theo thứ tự \texttt{user\_id} tăng dần. S10 không đổi thành viên nên không cần khóa tư vấn; G09 cập nhật chính các dòng này với \texttt{FOR UPDATE}, nên phân công và chuyển quyền luôn diễn ra theo một thứ tự xác định. Khóa trên dòng người được giao còn bảo vệ \br{BR07}: giao dịch rời nhóm hoặc xóa thành viên xóa dòng thành viên trước rồi mới gọi \texttt{MemberLeaveHook} để trả mặt hàng về Cần mua, nên phải chờ phân công kết thúc và thấy phân công đó; ngược lại, nếu người được giao đã rời nhóm, phép đọc không thấy dòng nào và lời gọi nhận 422 mã 4004 kèm \br{BR10}, \msg{MSG35}. Mặt hàng và danh sách bị khóa \texttt{FOR UPDATE} rồi mới so phiên bản; phân công là quyết định điều phối như đánh dấu đã mua nên không áp dụng quy tắc lệnh sau thắng, phiên bản lệch trả 409 mã 4002 kèm bản hiện tại và \msg{MSG35}. Điều kiện trạng thái là \texttt{to\_buy} hoặc \texttt{assigned} (\br{BR10}, Hình~\ref{fig:trang-thai-mat-hang-be}) và danh sách chưa Hoàn tất (\br{BR08}). Mỗi giao dịch kết thúc bằng \texttt{space\_changes}, sau commit phát \texttt{space:\{sid\}}; S11 đi qua cùng các bước kiểm tra vai trò và phiên bản.

Thông báo phải gộp: khi trưởng nhóm giao năm mặt hàng liên tiếp, người được giao chỉ nhận một thông báo “Bạn được giao mua $k$ mặt hàng trong danh sách <tên>” (bước 6). Đồng thời, module notification không được đọc bảng của shopping (mục~\ref{sec:kien-truc}), nên $k$ phải nằm sẵn trong nội dung sự kiện. Hai yêu cầu được dung hòa bằng một tác vụ gom của chính shopping. Giao dịch S10 chèn tác vụ \texttt{AssignmentDigest} duy nhất theo cặp người được giao và danh sách, lịch chạy lùi 5 giây; bốn lần chèn sau trùng tác vụ đang chờ nên bị River bỏ qua. Khi chạy, tác vụ đếm các mặt hàng của danh sách được giao cho người đó từ mốc đợt giao, rồi chèn \texttt{PushNotification} mang sự kiện \texttt{item.assigned} với $k$, tên danh sách và mốc đợt. Worker thông báo tôn trọng thiết lập \texttt{assignment\_notify\_on} (\uc{UC08}), lưu thông báo vào hộp thư với \texttt{dedupe\_key} gồm người nhận, danh sách và mốc đợt để lần thử lại không sinh bản trùng, rồi đẩy FCM tới từng thiết bị; lỗi tạm thời được thử lại tối đa ba lần trên hàng đợi \texttt{notify} (bước 6a, \nfr{NFR12}). Tự giao cho chính mình thì không đẩy; bỏ phân công dùng cùng cơ chế với sự kiện \texttt{item.unassigned} để báo người được giao trước đó. Về lỗi, người không phải trưởng nhóm nhận 403 mã 3007 (\msg{MSG13}) ở cả hai lớp kiểm tra; mặt hàng Đã mua, Đã bỏ qua, Đã nhập tủ hoặc thuộc danh sách Hoàn tất nhận 422 mã 4004 kèm \br{BR10} hoặc \br{BR08}. Ứng dụng tổng hợp các phản hồi, hiển thị \msg{MSG10} cùng số mặt hàng bị bỏ qua.

% =====================================================================
\subsection{Sinh danh sách mua sắm từ kế hoạch bữa ăn (UC22)}\label{sssec:ht-sinh-ds}

Theo Bảng~\ref{tab:UC22}, \uc{UC22} có hai pha: tính đề xuất (bước 2--5) và tạo danh sách sau khi người dùng chỉnh sửa (bước 6--7). Cả hai pha đi qua S15 \texttt{POST} \path{.../shopping-lists/generate}, do \texttt{MealHandler} chuyển cho \texttt{MealPlanService.GenerateList} của module mealplan; lời gọi thứ nhất mang \texttt{dryRun}, lời gọi thứ hai mang các dòng đã chọn, như Hình~\ref{fig:seq-sinh-ds-ke-hoach}. Ba cổng của module khác được vẽ gộp thành một bên tham gia.

\diagram[\textwidth]{c3-31-seq-sinh-ds-ke-hoach}{Biểu đồ tuần tự hiện thực use case Sinh danh sách mua sắm từ kế hoạch bữa ăn}{fig:seq-sinh-ds-ke-hoach}

Dù đường dẫn nằm dưới tài nguyên danh sách, dịch vụ thuộc mealplan vì phép tính cần dữ liệu của bốn module. Mealplan được phép gọi cổng \texttt{RecipeQuery} của recipe, \texttt{StockQuery} của pantry và \texttt{ListComposer} của shopping (hai phương thức \texttt{PendingQuantities} và \texttt{Compose}), còn recipe đã phụ thuộc shopping vì \uc{UC36}; đặt \uc{UC22} trong shopping sẽ tạo chu trình shopping $\to$ recipe $\to$ shopping. Khoảng ngày phải thỏa \texttt{from} $\le$ \texttt{to} và nằm trong cửa sổ từ \texttt{\$today} đến số ngày của tham số \texttt{mealplan.horizonDays} (mặc định 28) theo \br{BR18}, nếu không trả 400 mã 2002 (\msg{MSG02}). Pha xem trước không ghi gì: bốn phép đọc gồm món chưa nấu có công thức, nguyên liệu nhân theo tỷ lệ khẩu phần, lượng đang chờ mua và lô còn hạn chạy trong một giao dịch \texttt{REPEATABLE READ READ ONLY}. Nhờ ảnh chụp chung, khi thành viên khác đang xác nhận nấu (\uc{UC41}) hay nhập tủ (\uc{UC21}), phép tính không thấy trạng thái nửa trước nửa sau như lô đã bị trừ mà món vẫn Dự kiến; giao dịch chỉ đọc cũng không giữ khóa dòng nên không cản ai. Lượng đang chờ mua gồm mặt hàng Cần mua, Đã phân công và Đã mua nhưng chưa nhập tủ, thuộc danh sách chưa Hoàn tất có ngày mua không sau ngày nấu (\br{BR20}). Thành phần miền \texttt{ShortageCalculator.Compute} là hàm thuần, kiểm thử đơn vị được trực tiếp bằng ví dụ cà rốt của Chương~\ref{chap:dac-ta}; thuật toán đã trình bày ở mục~\ref{sssec:kt-goi-y}. Món tên tự do được trả trong danh sách riêng để ghi chú (bước 2a). Khi không thiếu nguyên liệu nào, S15 trả 200 mã 1000 kèm \msg{MSG15} và danh sách rỗng (bước 5a), vì đây là thông tin chứ không phải lỗi.

Lời gọi xác nhận không tính lại phần thiếu, vì đề xuất chỉ mang tính tư vấn và người dùng được chỉnh ở bước 6. Trước khi mở giao dịch, dịch vụ kiểm tra số lượng (\br{BR14}, \msg{MSG16}) và ngày dự kiến mua không trước \texttt{\$today} (\br{BR08}, \msg{MSG02}); khi lệnh được phát lại từ hàng đợi ngoại tuyến, mốc so sánh là ngày của header \texttt{X-Client-Time} nếu thời điểm đó trong vòng 7 ngày. Giao dịch ghi bắt đầu bằng khóa lũy đẳng rồi giao việc tạo cho cổng \texttt{ListComposer.Compose}, nhờ đó quy tắc gộp chỉ cài đặt ở module shopping. Tạo mới (bước 7) chèn danh sách với \texttt{origin = meal\_plan}, trạng thái \texttt{drafting} cùng các mặt hàng \texttt{to\_buy}. Với nhánh gộp (bước 6a), ứng dụng lấy các danh sách ứng viên qua S01 lọc theo trạng thái và khoảng ngày; \texttt{Compose} khóa \texttt{FOR UPDATE} dòng danh sách đích để tuần tự hóa với thành viên đang thêm mặt hàng (\uc{UC17}), danh sách đã Hoàn tất làm giao dịch bị hủy với 422 mã 4004 kèm \br{BR08}, còn lại gộp theo \br{BR09} và trả \msg{MSG11}. Giao dịch kết thúc bằng \texttt{space\_changes}, sau commit phát \texttt{space:\{sid\}}; khóa lũy đẳng bảo đảm gửi lại không tạo hai danh sách.

% =====================================================================
\subsection{Thêm thực phẩm vào tủ lạnh (UC25)}\label{sssec:ht-them-thuc-pham}

Phương thức \texttt{PantryService.AddLot} hiện thực \uc{UC25} qua P02 \texttt{POST} \path{/spaces/{sid}/lots}, với các API hỗ trợ trong cùng luồng: P08 cho nhánh quét mã vạch (bước 2a), Q02 \texttt{GET} \path{/spaces/{sid}/foods/suggest} cho gợi ý tên (bước 3--4), Q03 cho danh mục loại, đơn vị và vị trí lưu trữ, Y03 và Y04 cho ảnh của thực phẩm mới (bước 3a). Hình~\ref{fig:seq-them-thuc-pham} thể hiện luồng đầy đủ; nhánh mã vạch được tham chiếu sang Hình~\ref{fig:seq-ma-vach}.

\diagram[\textwidth]{c3-32-seq-them-thuc-pham}{Biểu đồ tuần tự hiện thực use case Thêm thực phẩm vào tủ lạnh}{fig:seq-them-thuc-pham}

Q02 tìm trên Elasticsearch với bộ lọc phạm vi bắt buộc, chỉ gồm thực phẩm dùng chung và thực phẩm riêng của không gian hiện hành (\nfr{NFR09}, \br{BR23}) và trả kèm đơn vị, vị trí mặc định theo loại cùng số ngày hạn dùng theo từng vị trí lấy từ \texttt{catalog.food\_shelf\_lives}, nên ứng dụng điền sẵn ngày hết hạn ở bước 4 mà không cần thêm lượt gọi. P02 dùng được khi ngoại tuyến (\nfr{NFR11}) nhờ id UUIDv7 do thiết bị sinh và \texttt{Idempotency-Key}; khi phát lại, ràng buộc ngày mua không sau hôm nay được xét theo ngày của \texttt{X-Client-Time}. Dịch vụ kiểm tra dữ liệu vào trước mọi truy vấn, đúng bước 6: trường bắt buộc, số lượng (\br{BR14}), ngày mua không sau \texttt{\$today} và ngày hết hạn không trước ngày mua (\br{BR11}); ràng buộc cuối được lặp lại bằng \texttt{CHECK} \texttt{ck\_lot\_dates} trên \texttt{food\_lots}. Trong giao dịch, sau khóa lũy đẳng, thực phẩm phải là thực phẩm dùng chung đang hoạt động hoặc thực phẩm riêng của chính không gian; thực phẩm của không gian khác làm giao dịch bị hủy với 404 như không tồn tại (\br{BR23}). Ngày hết hạn bỏ trống được đề xuất theo \br{BR11} bằng ngày mua cộng số ngày của cặp thực phẩm, vị trí; thực phẩm chưa có hạn dùng gợi ý thì trường này bắt buộc. Ngày hết hạn đã qua (bước 7a) không bị chặn: ứng dụng hỏi xác nhận, máy chủ vẫn lưu và trả trạng thái Đã hết hạn được suy diễn lúc đọc (\br{BR12}).

Thực phẩm mới (bước 3a) được gửi trong trường \texttt{newFood} và tạo trong cùng giao dịch qua cổng \texttt{catalog/api.CreateSpaceFood}: bản ghi \texttt{catalog.foods} có \texttt{scope = space}, chủ sở hữu là không gian hiện hành và \texttt{status = pending\_review}; chỉ mục duy nhất từng phần theo không gian, loại và tên ngăn hai thành viên tạo trùng, khi xung đột thì dùng lại bản ghi đã có. Sự kiện \texttt{food.changed} sinh tác vụ \texttt{IndexFood} để thực phẩm xuất hiện trong gợi ý của riêng không gian này và trong hàng chờ duyệt của \uc{UC49} (Hình~\ref{fig:seq-duyet-thuc-pham}). Ảnh không đi qua API: Y03 nhận \texttt{\{purpose: food, targetId\}}, kiểm tra quyền trên đối tượng đích, tạo \texttt{media.objects} ở trạng thái \texttt{pending} và trả URL \texttt{PUT} ký sẵn (5 phút, tối đa 5 MB, JPEG hoặc PNG). Ứng dụng tải thẳng lên MinIO rồi gọi Y04; API dùng lệnh \texttt{HEAD} để xác nhận đối tượng có thật, đúng kích thước và kiểu, rồi chèn tác vụ xử lý ảnh cho sự kiện \texttt{media.uploaded} để \texttt{MediaWorker} nén ảnh và tạo ảnh thu nhỏ WebP. Giao dịch ghi chèn \texttt{food\_lots} với \texttt{unit\_factor} là ảnh chụp hệ số quy đổi, để việc sửa đơn vị ở \uc{UC51} không làm đổi lượng của lô cũ, cùng biến động \texttt{stock\_in} làm đầu vào tỷ lệ lãng phí của \uc{UC46}; lô nhập tay không có \texttt{total\_cost} vì chỉ S14 biết đơn giá. Giao dịch kết thúc bằng \texttt{space\_changes}; sau commit, dịch vụ xóa bộ đệm gợi ý \texttt{suggest:\{sid\}} và phát \texttt{space:\{sid\}}; phản hồi 201 trả lô kèm trạng thái hạn dùng theo ngưỡng $N$ của người gọi. Thiếu trường trả 400 mã 2001 (\msg{MSG01}); số lượng sai và ngày hết hạn trước ngày mua trả 400 mã 2002 (\msg{MSG16}, \msg{MSG14}). Lỗi tra mã vạch (mã 4001 hoặc 5004, \msg{MSG23}) đưa người dùng về biểu mẫu để nhập tay, đúng bước 4b của \uc{UC26}.

% =====================================================================
\subsection{Gợi ý món ăn (UC40)}\label{sssec:ht-goi-y-mon}

Phương thức chỉ đọc \texttt{MealPlanService.Suggest} hiện thực \uc{UC40} qua M05 \texttt{GET} \path{/spaces/{sid}/recipe-suggestions} với các tham số \texttt{date}, \texttt{mealSlot} (\texttt{breakfast}, \texttt{lunch} hoặc \texttt{dinner}), số khẩu phần và \texttt{expiringOnly}. Cùng một API phục vụ ba điểm vào: thẻ “Gợi ý cho bạn” của \uc{UC37} (bước 2a), màn hình gợi ý độc lập và bước 3 của \uc{UC30}. Ở \uc{UC30}, ứng dụng ghép P07 với M05 có \texttt{expiringOnly = true} vì pantry không được gọi mealplan. Hình~\ref{fig:seq-goi-y-mon} thể hiện luồng xử lý.

\diagram[\textwidth]{c3-33-seq-goi-y-mon}{Biểu đồ tuần tự hiện thực use case Gợi ý món ăn từ thực phẩm sẵn có}{fig:seq-goi-y-mon}

Theo điều kiện sau của Bảng~\ref{tab:UC40}, kết quả gợi ý được lưu 10 phút hoặc đến khi dữ liệu đầu vào đổi, trong Redis hash \texttt{suggest:\{sid\}} với trường gồm người dùng, ngày, bữa, số khẩu phần và cờ \texttt{expiringOnly}, vì tồn kho thuộc không gian còn ngưỡng $N$ và hệ số $F$ thuộc từng người. Hash bị xóa sau commit của mọi giao dịch phát \texttt{lot.changed} hay \texttt{recipe.changed}, khi quản trị viên đổi tham số gợi ý (\texttt{config.changed}) và khi A13 đổi $N$ của một người (xóa ở mọi không gian của người đó); thời hạn sống chặn trên độ cũ nếu lệnh xóa thất bại. Khi bộ đệm trống, dịch vụ đọc $\theta$, trọng số và số món tối đa từ \texttt{cfg:all}, nên thay đổi ở \uc{UC53} có hiệu lực ngay lần tính sau. Ba truy vấn trong một giao dịch chỉ đọc lấy ngưỡng $N$ của người xem, tồn kho khả dụng qua \texttt{StockQuery.Available} (lô \texttt{active} có \texttt{expires\_on} $\ge$ \texttt{\$today}) và qua \texttt{RecipeQuery} các công thức ở trạng thái \texttt{published} có ít nhất một nguyên liệu trong tủ, lọc theo bữa, cùng danh sách yêu thích; công thức mẫu đang ẩn hay bản nháp không bao giờ được gợi ý. Thành phần miền \texttt{SuggestionEngine.Rank} chấm điểm theo \br{BR19} (mục~\ref{sssec:kt-goi-y}) với mục tiêu 2 giây (\nfr{NFR09}); với \texttt{expiringOnly}, chỉ giữ các món có $E > 0$, tức dùng ít nhất một lô sắp hết hạn, và ứng dụng hiển thị ba món đầu. Tủ lạnh trống (bước 2a) trả công thức mẫu được yêu thích nhiều nhất kèm lời nhắc thêm thực phẩm; không món nào đạt ngưỡng (bước 4a) trả ba món có $R$ cao nhất kèm \msg{MSG31}, cả hai đều mang mã 1000. Tham số sai trả 400 mã 2002 (\msg{MSG02}). Là phép đọc thuần, M05 không dùng khóa lũy đẳng và không giữ khóa dòng.

% =====================================================================
\subsection{Lên lịch bữa ăn (UC37)}\label{sssec:ht-lap-lich}

Sau khi người dùng chọn món ở một trong ba thẻ, gồm Gợi ý (M05, Hình~\ref{fig:seq-goi-y-mon}), Yêu thích và Tất cả công thức (R01, chỉ trả công thức \texttt{published}), M02 \texttt{POST} \path{/spaces/{sid}/meals} gọi \texttt{MealPlanService.PlanMeals} để ghi các món vào lịch, như Hình~\ref{fig:seq-lap-lich}.

\diagram[\textwidth]{c3-40-seq-lap-lich}{Biểu đồ tuần tự hiện thực use case Lên lịch bữa ăn}{fig:seq-lap-lich}

M02 nhận một ngày, một bữa và mảng món; mỗi món có id do thiết bị sinh, tham chiếu công thức hoặc tên tự do (bước 3a) và số khẩu phần, nên chọn nhiều món ở bước 3 được ghi trong một giao dịch. Trước mọi truy vấn, dịch vụ kiểm tra \br{BR18}: ngày nằm từ \texttt{\$today} đến hết cửa sổ lập kế hoạch đọc từ tham số \texttt{mealplan.horizonDays} (mặc định 28 ngày, miền 7–60, nếu vi phạm trả \msg{MSG28}), và số khẩu phần từ 1 đến 20. Khi lệnh được phát lại từ hàng đợi ngoại tuyến, \br{BR18} được xét theo ngày của \texttt{X-Client-Time} nếu thời điểm đó trong vòng 7 ngày và không vượt giờ máy chủ; ngoài khoảng đó máy chủ trả 422 mã 4004 và ứng dụng giữ bản nháp để người dùng sửa. Trong giao dịch, sau khóa lũy đẳng, \texttt{RecipeQuery} xác nhận công thức là công thức mẫu đang công bố hoặc thuộc chính không gian; công thức ngoài phạm vi làm giao dịch bị hủy với 404 mã 4001. Món được chèn ở trạng thái \texttt{planned}, giao dịch kết thúc bằng \texttt{space\_changes} và sau commit phát \texttt{space:\{sid\}}. Món trùng trong cùng bữa (bước 5a) được ứng dụng phát hiện trên bản sao cục bộ; nếu người dùng chọn tăng khẩu phần, ứng dụng gọi M03 thay cho M02. Sau commit, dịch vụ đối chiếu nguyên liệu bằng \texttt{RecipeQuery.Ingredients} và \texttt{StockQuery.Available}, cũng là thành phần dùng chung với \uc{UC34}, và trả nhãn đủ hoặc thiếu $k$ nguyên liệu cho từng món (bước 5); món tên tự do không tham gia đối chiếu. Phép đối chiếu nằm ngoài giao dịch ghi để giữ giao dịch ngắn. Ngày ngoài khoảng trả 400 mã 2002 với \msg{MSG28}, khẩu phần sai trả \msg{MSG02}; M03 áp dụng cùng quy tắc khi di chuyển món (\uc{UC39}), còn món đã nấu không được sửa.

% =====================================================================
\subsection{Báo cáo lãng phí và xuất báo cáo (UC46, UC47)}\label{sssec:ht-bao-cao}

\texttt{ReportService} hiện thực \uc{UC46} qua B03 \texttt{GET} \path{.../reports/waste} và \uc{UC47} qua cặp B04, B05: B04 nhận yêu cầu xuất và trả ngay, tác vụ River \texttt{ExportReport} kết xuất tệp vào MinIO, B05 trả trạng thái và cấp URL tải ký sẵn. Report chỉ đọc dữ liệu của module khác và không gọi cổng nào, nên Hình~\ref{fig:seq-bao-cao-xuat} vẽ handler và dịch vụ thành một bên tham gia.

\diagram[\textwidth]{c3-34-seq-bao-cao-xuat}{Biểu đồ tuần tự hiện thực use case Xem báo cáo lãng phí và Xuất báo cáo}{fig:seq-bao-cao-xuat}

Bước 4 của \uc{UC46} kiểm tra \br{BR21}; khoảng sai hoặc dài hơn 12 tháng trả 400 mã 2002 (\msg{MSG25}). Report là module mô hình đọc (read model): ngoại lệ có chủ đích của quy tắc phụ thuộc cho phép nó \texttt{SELECT} qua các khung nhìn \texttt{report.v\_*} đặt trên bảng của shopping, pantry, mealplan và catalog (mục~\ref{sec:kien-truc}). Nguồn số liệu được thiết kế lai để chi phí không tăng theo độ dài khoảng. Phần lịch sử đến hết hôm qua đọc từ \texttt{report.daily\_space\_stats}, khung nhìn cụ thể hóa làm mới hằng đêm, mỗi dòng ứng với một không gian, một ngày, một loại thực phẩm và một nhóm đại lượng, nên khoảng 12 tháng chỉ phải cộng vài nghìn dòng đã tổng hợp thay vì quét toàn bộ biến động kho. Các chỉ số không tổng hợp sẵn gồm số lô bị loại bỏ trên số lô nhập tủ (tỷ lệ lãng phí theo Bảng~\ref{tab:chi-so-bao-cao}), phân bổ theo lý do, năm thực phẩm bị bỏ nhiều nhất và số liệu của \texttt{\$today}; chúng được truy vấn trực tiếp từ \texttt{pantry.stock\_movements} qua chỉ mục (\texttt{space\_id}, \texttt{occurred\_at}). Giá trị lãng phí của một lần loại bỏ bằng \texttt{total\_cost} $\times$ lượng loại bỏ / \texttt{initial\_qty}, trong đó \texttt{total\_cost} của lô là số lượng thực tế nhân đơn giá ghi nhận khi nhập tủ từ mặt hàng nguồn (S14). Cách tính theo tỷ lệ cho đúng đơn vị kể cả khi S14 đổi đơn vị so với mặt hàng; lô không có \texttt{total\_cost} chỉ tính theo số lượng và được đếm riêng (\br{BR21}). Hai phần đọc chung một giao dịch chỉ đọc nên các thẻ chỉ số nhất quán; truy vấn cùng độ dài cho kỳ trước cho mức thay đổi ở bước 6, và cờ \texttt{empty} phục vụ bước 5a. Người không phải thành viên, kể cả quản trị viên, nhận 404 mã 4001 (\br{BR23}).

\uc{UC47} được xử lý bất đồng bộ, vì kết xuất PDF có biểu đồ tốn vài giây CPU, nếu giữ trong yêu cầu sẽ chiếm tiến trình API, vỡ mục tiêu 95\% API dưới 1 giây (\nfr{NFR09}) và dễ đứt trên mạng di động. B04 nhận loại báo cáo, khoảng thời gian và định dạng PDF hoặc CSV (bước 1). Sau khóa lũy đẳng, giao dịch chèn một dòng vào bảng \texttt{report.exports} do module report sở hữu, gồm không gian, người yêu cầu, loại, khoảng, định dạng và trạng thái \texttt{pending}, cùng tác vụ \texttt{ExportReport} trên hàng đợi \texttt{report}, duy nhất theo người dùng, loại, khoảng và định dạng để chạm nhiều lần không sinh nhiều tệp; API trả ngay 202 mã 1001 kèm \texttt{exportId}. Bảng \texttt{media.objects} chỉ dành cho ảnh. Worker gọi lại đúng truy vấn của B01--B03 nên số liệu trong tệp trùng với màn hình, kết xuất PDF hoặc CSV mã hóa UTF-8 (bước 2). Khi thành công, tệp được ghi vào bucket \texttt{exports} có vòng đời 7 ngày; trong một giao dịch, dòng xuất chuyển sang \texttt{ready} với \texttt{object\_key}, \texttt{expires\_at}, và tác vụ \texttt{PushNotification} cho sự kiện \texttt{export.completed} được chèn để module notification báo “Báo cáo đã sẵn sàng”. Kết xuất lỗi (bước 2a) được River thử lại tối đa ba lần; sau lần cuối, dòng xuất chuyển sang \texttt{failed} kèm lỗi và người dùng cũng được báo qua \texttt{PushNotification}. Ứng dụng gọi B05 khi nhận thông báo hoặc thăm dò thưa dần khi màn hình còn mở. Sau bước kiểm tra thành viên chung, B05 lọc theo \texttt{id}, \texttt{space\_id} và \texttt{requested\_by}, nên \texttt{exportId} không tồn tại, thuộc không gian khác hoặc của người khác đều nhận 404 mã 4001. Phản hồi phân theo trạng thái: \texttt{pending} để ứng dụng tiếp tục chờ; \texttt{ready} kèm URL \texttt{GET} ký sẵn 15 phút, ứng dụng tải thẳng từ MinIO và mở bảng chia sẻ của hệ điều hành (bước 3); \texttt{failed} để ứng dụng hiện nút thử lại, tức gửi một yêu cầu B04 mới.

\InputIfFileExists{chapters/c3/ht-quan-tri.tex}{}{}

% =====================================================================
\subsection{Các use case còn lại}\label{sssec:ht-con-lai}

Đăng nhập và làm mới phiên (\uc{UC03}) có biểu đồ ở mục~\ref{ssec:xac-thuc-phien} (Hình~\ref{fig:seq-dang-nhap-lam-moi}). Các use case khác đã có biểu đồ ở mục~\ref{ssec:ky-thuat} gồm \uc{UC20} (Hình~\ref{fig:seq-outbox-mua-hang}), \uc{UC21} (Hình~\ref{fig:seq-nhap-tu}), \uc{UC26} (Hình~\ref{fig:seq-ma-vach}), \uc{UC31} (Hình~\ref{fig:seq-nhac-han-dung-be}), \uc{UC41} (Hình~\ref{fig:seq-xac-nhan-nau}), \uc{UC42} và \uc{UC43} (Hình~\ref{fig:seq-tim-kiem-chi-muc}); trong đó \uc{UC21} gọi cổng \texttt{StockIn.AddLots} và \uc{UC41} gọi \texttt{StockConsumer.ConsumeForMeal} của pantry trong cùng giao dịch với module gọi. \uc{UC28} dùng chung cơ chế FEFO của Hình~\ref{fig:seq-xac-nhan-nau} cho một thực phẩm: P05 gọi \texttt{PantryService.Consume}, khóa các lô \texttt{active} có \texttt{expires\_on} $\ge$ \texttt{\$today} theo thứ tự (\texttt{expires\_on}, \texttt{id}) rồi trừ dần theo \br{BR15}; khi chọn đích danh một lô (bước 3a) chỉ lô đó bị khóa, lô hết hạn chỉ bị trừ khi được chọn tường minh (bước 4b), và lượng vượt lượng còn trả 422 mã 2003 (\msg{MSG17}). P04 của \uc{UC27} không áp dụng lệnh sau thắng cho lượng còn lại: \texttt{If-Match} lệch trả 409 mã 4002, thay đổi lượng ghi biến động \texttt{adjust}, lượng lớn hơn lượng ban đầu trả 400 mã 2002 và lượng bằng 0 chuyển lô sang Đã dùng hết; vị trí, ngày hết hạn và ghi chú thì hợp nhất theo trường kèm \texttt{overwrittenFields}. P06 của \uc{UC29} bắt buộc lý do và ghi biến động \texttt{discard} cho \uc{UC46}.

Ở gói tài khoản và nhóm, A15 \texttt{DELETE} \path{/me/devices/{deviceId}} gỡ mã FCM khi đăng xuất (\uc{UC05}). A10 cập nhật hồ sơ (\uc{UC06}) ghi \texttt{space\_changes} với thực thể thành viên ở mọi không gian của người dùng để tên và ảnh đại diện mới đến được các thành viên khác. A11 đổi mật khẩu (\uc{UC07}) thu hồi mọi phiên khác, tăng \texttt{token\_version} rồi cấp cặp mã mới cho phiên hiện tại, nên thiết bị khác bị đăng xuất ngay còn người dùng không phải đăng nhập lại (bước 3). A12 trả thiết lập kèm khung giờ nhắc \texttt{reminder.runAt} và ngưỡng $N$ mặc định ở dạng chỉ đọc; A13 đổi $N$ thì xóa bộ đệm gợi ý của mọi không gian của người đó (\uc{UC08}). A16 xóa tài khoản theo \nfr{NFR18}: người dùng nhập lại mật khẩu, trưởng nhóm của nhóm còn thành viên khác bị từ chối với 422 mã 4004 kèm \msg{MSG24} (\br{BR07}); ngay lập tức mọi phiên bị thu hồi, \texttt{deleted\_at} được đặt, đăng nhập bị chặn và thiết bị bị xóa; sau 30 ngày, tác vụ nền ẩn danh email và họ tên, xóa ảnh, dữ liệu không gian cá nhân và tài liệu Elasticsearch, còn dữ liệu đóng góp trong nhóm được giữ với tên “Người dùng đã xóa”. G02 tạo nhóm (\uc{UC09}) dùng chung chỉ mục \br{BR04} với \uc{UC11}; trường \texttt{copyPersonalData} gồm danh sách, lô và công thức làm giao dịch chèn tác vụ \texttt{CopySpaceData} lũy đẳng theo không gian để sao chép sau commit, còn G04 đổi tên, ảnh nhóm. G05 tạo mã mời (\uc{UC10}) chỉ dành cho trưởng nhóm, lấy khóa tư vấn của không gian, từ chối không gian cá nhân với 422 mã 4004, thu hồi mã cũ và chèn mã mới hết hạn sau \texttt{invite.ttlHours} trong một giao dịch (\br{BR06}); Caddy phục vụ hai tệp \path{/.well-known/apple-app-site-association} và \path{/.well-known/assetlinks.json} để liên kết mời mở được ứng dụng (mục~\ref{sec:trien-khai}).

G03 (\uc{UC12}) trả thành viên, vai trò và \texttt{recentActivity} gồm 20 thay đổi gần nhất trong \texttt{sync.space\_changes} kèm tên người thực hiện (bước 3); \uc{UC13}--\uc{UC15} đã được trình bày ở Hình~\ref{fig:seq-roi-nhom-chuyen-quyen}.

Ở các gói còn lại, S01 (\uc{UC16}, \uc{UC23}) lọc theo trạng thái, khoảng ngày và gom theo tháng, phản hồi có tổng chi tiêu và người mua. S07 khóa danh sách trước khi gộp theo \br{BR09} (\msg{MSG11}); bỏ qua mặt hàng là lệnh S08 đặt \texttt{status = skipped}; mặt hàng Đã nhập tủ chỉ đọc (\msg{MSG32}). Danh sách Hoàn tất chặn sửa nội dung với 422 mã 4004 (\br{BR08}), trừ ba lệnh: S08 sửa số lượng thực tế và đơn giá của mặt hàng Đã mua, S13 hoàn tác đưa danh sách về Đang mua, và S14 nhập tủ. Công thức thuộc không gian: ở không gian cá nhân, đây là “bộ sưu tập cá nhân” của Chương~\ref{chap:dac-ta}. Chỉ người tạo hoặc trưởng nhóm được sửa, xóa công thức (\uc{UC33}); sửa công thức mẫu trả 403 mã 3007. R09 sao chép một công thức mẫu thành bản của không gian, ghi \texttt{origin\_recipe\_id} (\uc{UC32}, \br{BR17}). Trạng thái nháp, công bố và ẩn chỉ dùng cho công thức mẫu (\uc{UC52}); công thức của không gian luôn ở trạng thái công bố. R06 (\uc{UC35}) chỉ nhận công thức mẫu đã công bố hoặc công thức của không gian mà người gọi là thành viên, ngược lại trả 404 mã 4001. Xóa công thức chuyển các món Dự kiến thành món tên tự do qua \texttt{RecipeDeleteHook}. R03 (\uc{UC34}) và R08 (\uc{UC36}) đối chiếu nguyên liệu với tồn kho qua \texttt{StockQuery}; R08 giao phần thiếu cho \texttt{ListComposer}.

P01 và Q01 (\uc{UC43}) nhận bộ lọc loại, vị trí, trạng thái hạn dùng và từ khóa, phản hồi kèm số lượng theo loại (facets); Q01 tìm danh mục chung, thực phẩm riêng và công thức của không gian trên Elasticsearch với bộ lọc phạm vi bắt buộc, tìm lô trong không gian qua \texttt{StockQuery.SearchLots} trên PostgreSQL và chuyển hẳn sang PostgreSQL khi Elasticsearch lỗi. B01 (\uc{UC44}) gom theo ngày, tuần hoặc tháng; B01, B02 (\uc{UC45}) dùng cùng nguồn lai như B03. Ngoài \uc{UC48}, \uc{UC49} và \uc{UC53} đã trình bày ở trên, các use case quản trị theo cùng khuôn: mọi API yêu cầu \texttt{system\_role = admin} và ghi \texttt{admin.audit\_logs} gồm giá trị cũ và mới trong cùng giao dịch. D05 và D06 (\uc{UC50}, \uc{UC51}) xóa được bản ghi chưa được tham chiếu; bản ghi đang được tham chiếu trả 422 mã 4004 kèm \br{BR22} và \msg{MSG22} để ứng dụng đề nghị chuyển sang Ngừng sử dụng, mọi thay đổi đều làm đổi ETag của Q03. Đơn vị cơ sở là duy nhất trong mỗi nhóm đại lượng và nhóm đại lượng của đơn vị đang được tham chiếu không đổi được (\br{BR16}). D07 (\uc{UC52}) quản lý công thức mẫu; mỗi lần công bố hay ẩn phát \texttt{recipe.changed} để đánh lại chỉ mục.

\endgroup
% ---------------------------------------------------------------------
%  PHỤ LỤC. ĐẶC TẢ CHI TIẾT API BỔ SUNG (tiếp mục Thiết kế API của Chương 3)
% ---------------------------------------------------------------------
\section{Đặc tả chi tiết API bổ sung}\label{pl:dac-ta-api}

Phụ lục này bổ sung cho mục~\ref{sec:api}: các API dùng chung quy ước, mã phản hồi và macro bảng đặc tả của mục đó. Phần đầu ghi tham số và hành vi của các API không đặc tả chi tiết; phần sau đặc tả đầy đủ tám API có rủi ro thiết kế cao, mỗi API kèm test case dự kiến.

\subsection{Tham số và hành vi bổ sung}\label{pl:api-bo-sung}

Phần lớn API ngoài mục~\ref{ssec:dac-ta-api} là thao tác thêm, sửa, xóa theo đúng quy ước ở trên. Bảng~\ref{tab:c3-api-bo-sung} ghi những API có tham số hoặc hành vi không suy ra được từ quy ước, để hợp đồng OpenAPI và kiểm thử tích hợp có căn cứ.

{\scriptsize\renewcommand{\arraystretch}{1.05}
\begin{longtable}{|C{1.15cm}|L{5.0cm}|L{\dimexpr\textwidth-6\tabcolsep-4\arrayrulewidth-1.15cm-5.0cm\relax}|}
\caption{Tham số và hành vi bổ sung của các API không đặc tả chi tiết}\label{tab:c3-api-bo-sung}\\
\hline
\hd{Mã} & \hd{Yêu cầu} & \hd{Hành vi, phản hồi và lỗi}\\ \hline
\endfirsthead
\multicolumn{3}{l}{\footnotesize\itshape Bảng \thetable\ (tiếp theo)}\\ \hline
\hd{Mã} & \hd{Yêu cầu} & \hd{Hành vi, phản hồi và lỗi}\\ \hline
\endhead
\multicolumn{3}{|l|}{\hs{Tài khoản}}\\ \hline
A02 & \texttt{email}; \texttt{purpose}: \texttt{register} hoặc \texttt{reset\_password}; \texttt{code}: 6 chữ số & \texttt{register}: kích hoạt tài khoản, \texttt{ActivationHook} tạo không gian cá nhân trong cùng giao dịch, cấp phiên như A04. \texttt{reset\_password}: trả \texttt{resetToken} 256 bit, chỉ lưu băm, sống 10 phút, dùng một lần. Mã sai: 400, 3006 (MSG04); lần sai thứ năm vô hiệu mã (BR02)\\ \hline
A08 & \texttt{email}, \texttt{resetToken}, \texttt{newPassword} (BR01) & Đặt mật khẩu mới, đưa \texttt{failed\_login\_count} về 0, xóa \texttt{locked\_until}, thu hồi mọi phiên, tăng \texttt{token\_version} (BR03). Mã hết hạn, đã dùng hoặc không khớp email: 400, 3006\\ \hline
A10 & \texttt{fullName} (2--50 ký tự), \texttt{avatarObjectId}, \texttt{locale} & Ghi một dòng \texttt{space\_changes} (thực thể \texttt{member}) cho mọi không gian của người dùng, để thành viên khác thấy tên và ảnh mới (UC06)\\ \hline
A11 & \texttt{currentPassword}, \texttt{newPassword} & Mật khẩu hiện tại sai: 400, 2002 (MSG02), tăng bộ đếm sai (BR03). Thành công: thu hồi mọi phiên khác, tăng \texttt{token\_version} và trả cặp mã mới cho phiên hiện tại (UC07 bước 3)\\ \hline
A12, A13 & A13: \texttt{reminderDays} (1--7 hoặc \texttt{null} để dùng mặc định), ba công tắc thông báo & A12 trả thêm \texttt{reminder.runAt} và $N$ mặc định của hệ thống ở dạng chỉ đọc (UC08 bước 2). A13 đổi $N$ thì xóa \texttt{suggest:\{sid\}} của mọi không gian của người dùng\\ \hline
A16 & \texttt{password} & Sai mật khẩu: 400, 2002. Là trưởng nhóm của nhóm còn thành viên khác: 422, 4004, BR07 (MSG24). Ngay lập tức thu hồi mọi phiên, đặt \texttt{deleted\_at}, chặn đăng nhập, xóa thiết bị; sau 30 ngày tác vụ nền ẩn danh email, họ tên, xóa ảnh, dữ liệu không gian cá nhân và tài liệu Elasticsearch. Dữ liệu đóng góp trong nhóm được giữ, hiển thị “Người dùng đã xóa” (NFR18)\\ \hline
\multicolumn{3}{|l|}{\hs{Không gian và nhóm}}\\ \hline
G02 & \texttt{name} (3--50 ký tự, BR05); \texttt{copyPersonalData}: \texttt{\{lists, lots, recipes\}} kiểu boolean & 201 kèm không gian mới, người gọi là trưởng nhóm; đã thuộc nhóm gia đình: 422, 4004, BR04 (MSG08). Sau commit chèn tác vụ \texttt{CopySpaceData}, lũy đẳng theo không gian, để sao chép phần dữ liệu được chọn (UC09)\\ \hline
G03 & -- & Thông tin nhóm, thành viên (trưởng nhóm đứng đầu), \texttt{changeSeq}, và \texttt{recentActivity}: 20 thay đổi gần nhất từ \texttt{sync.space\_changes} kèm tên người thực hiện (UC12 bước 3)\\ \hline
G05 & -- & Thời hạn đọc từ \texttt{invite.ttlHours}; mã cũ bị thu hồi trong cùng giao dịch (BR06); không gian cá nhân: 422, 4004. Trả \texttt{code}, \texttt{expiresAt} và \texttt{link} dạng \path{https://dicho.example/join/{code}}, mở thẳng ứng dụng qua Universal Links, App Links\\ \hline
\multicolumn{3}{|l|}{\hs{Mua sắm, tủ lạnh, tìm kiếm}}\\ \hline
S01 & \texttt{status}, \texttt{from}, \texttt{to}, \texttt{groupBy = month}, phân trang & Mỗi danh sách kèm số mặt hàng, tổng chi tiêu (mặt hàng có đơn giá) và những người đã mua; \texttt{groupBy = month} trả tổng theo tháng cho thẻ Lịch sử (UC23). Lọc \texttt{status} chưa hoàn tất và ngày trong khoảng cho danh sách đích ở UC22 bước 6a\\ \hline
P01, Q01 & \texttt{categoryId}, \texttt{location}, \texttt{expiryState} nhận nhiều giá trị; \texttt{q} & Cùng tham số là “hoặc”, khác tham số là “và”; phản hồi kèm \texttt{facets} đếm số kết quả theo từng loại, vị trí, trạng thái để vẽ bảng lọc (UC43); \texttt{expiryState} tính theo $N$ của người xem. Q01 tìm danh mục, công thức trên Elasticsearch với bộ lọc bắt buộc \texttt{scope = global} hoặc \texttt{owner\_space\_id = sid}; lô trong tủ tìm bằng PostgreSQL\\ \hline
P04 & Trường bất kỳ của lô; \texttt{If-Match} & \texttt{remainingQty} không theo LWW: lệch phiên bản trả 409, 4002; lớn hơn \texttt{initial\_qty} trả 400, 2002; bằng 0 chuyển \texttt{consumed}; mỗi lần đổi lượng ghi \texttt{stock\_movements} loại \texttt{adjust}. Vị trí, ngày hết hạn, ghi chú hợp nhất theo trường và trả \texttt{overwrittenFields}\\ \hline
Q02 & \texttt{q} (từ 2 ký tự), \texttt{limit} $\le$ 10 & Nằm dưới \path{/spaces/{sid}} để gồm thực phẩm riêng của không gian; 20 yêu cầu/giây/người; mỗi gợi ý kèm đơn vị, vị trí mặc định và hạn dùng gợi ý theo vị trí; Elasticsearch lỗi hoặc quá 300~ms thì chuyển sang PostgreSQL trigram\\ \hline
\multicolumn{3}{|l|}{\hs{Công thức}}\\ \hline
R01 & \texttt{source}: \texttt{space}, \texttt{template} hoặc \texttt{favorite}; \texttt{q}; phân trang & Chỉ trả công thức \texttt{published}: công thức của không gian và công thức mẫu đã công bố (UC34, UC37)\\ \hline
R04, R05 & R04: toàn bộ công thức (BR17), \texttt{If-Match} & Người tạo hoặc trưởng nhóm; người khác và mọi lệnh sửa công thức mẫu: 403, 3007. R05 gọi \texttt{RecipeDeleteHook}: món Dự kiến dùng công thức thành món tự do giữ tên, món Đã nấu giữ tham chiếu (UC33)\\ \hline
R06 & -- & Chỉ nhận công thức mẫu đã công bố hoặc công thức thuộc không gian mà người gọi là thành viên; ngược lại 404, 4001\\ \hline
R09 & \texttt{name} (tùy chọn) & Sao chép công thức mẫu đã công bố thành công thức của không gian, ghi \texttt{origin\_recipe\_id}; 201 kèm công thức mới (BR17). Công thức không phải mẫu: 404, 4001\\ \hline
\multicolumn{3}{|l|}{\hs{Báo cáo và tệp}}\\ \hline
B01 & \texttt{from}, \texttt{to} ($\le$ 12 tháng, BR21); \texttt{groupBy}: \texttt{day}, \texttt{week} hoặc \texttt{month} & Tổng chi tiêu, cơ cấu theo loại, chuỗi theo \texttt{groupBy}; không có dữ liệu: 200, 1000 (MSG26); khoảng sai: 400, 2002 (MSG25)\\ \hline
B04, B05 & B04: \texttt{reportType}, \texttt{from}, \texttt{to}, \texttt{format} (\texttt{pdf}, \texttt{csv}) & B04 chèn \texttt{report.exports} (\texttt{pending}) và tác vụ xuất trong cùng giao dịch, trả 202, 1001 kèm \texttt{exportId}. B05 lọc theo \texttt{space\_id} và \texttt{requested\_by} sau khi kiểm tra thành viên, nên tệp của người khác trả 404, 4001; trả \texttt{pending}, \texttt{ready} kèm URL \texttt{GET} ký sẵn 15 phút, hoặc \texttt{failed} kèm lỗi. Khi xong, \texttt{export.completed} sinh thông báo đẩy\\ \hline
Y03 & \texttt{purpose}: \texttt{avatar}, \texttt{space}, \texttt{recipe} hoặc \texttt{food}; \texttt{targetId}; \texttt{contentType}; \texttt{size} & Kiểm tra quyền trên đối tượng đích (chính mình, trưởng nhóm, người tạo công thức, thực phẩm riêng của không gian mình); trả URL \texttt{PUT} ký sẵn 5 phút, tối đa 5~MB, JPEG hoặc PNG\\ \hline
\multicolumn{3}{|l|}{\hs{Quản trị}}\\ \hline
D05, D06 & Mục danh mục; \texttt{DELETE} trên một mục & Chưa được tham chiếu thì xóa; đang được tham chiếu: 422, 4004, BR22 (MSG22), ứng dụng đề nghị chuyển \texttt{inactive}. D06 không cho đổi nhóm đại lượng của đơn vị đang được tham chiếu; mỗi nhóm có đúng một đơn vị cơ sở (BR16)\\ \hline
D09 & \texttt{values}: ánh xạ khóa tham số sang giá trị & Kiểm tra miền và tổng trọng số bằng 1 (MSG27), rồi ghi mọi khóa trong một giao dịch với một dòng \texttt{audit\_logs} và một sự kiện \texttt{config.changed}; không áp dụng một phần. Sau commit \texttt{PUBLISH cfg} để mọi bản xóa \texttt{cfg:all}; đổi \texttt{suggest.*} thì xóa cả \texttt{suggest:*} (UC53)\\ \hline
D11 & \texttt{action}: \texttt{approve} hoặc \texttt{merge}; \texttt{targetFoodId} khi gộp & Như đoạn dưới đây (UC49)\\ \hline
\end{longtable}
}

D11 khép lại vòng đời của thực phẩm riêng mà người dùng tạo ở \uc{UC25} bước 3a và \uc{UC26} bước 4a. Thực phẩm phải có \texttt{scope = space} và \texttt{status = pending\_review}; nếu quản trị viên khác vừa xử lý, yêu cầu nhận 409, mã 4002. Nhánh \texttt{approve} chuyển thực phẩm thành bản dùng chung (\texttt{scope = global}, bỏ \texttt{owner\_space\_id}, \texttt{status = active}). Nếu tên trong cùng loại hoặc mã vạch trùng một thực phẩm dùng chung, chỉ mục duy nhất theo phạm vi trả 409, mã 4003, và quản trị viên chỉ còn cách gộp. Nhánh \texttt{merge} yêu cầu đích là thực phẩm dùng chung đang hoạt động và khác chính nó; thực phẩm nguồn nhận \texttt{status = inactive}, \texttt{merged\_into = targetFoodId}, rồi sự kiện \texttt{food.merged} sinh các tác vụ River của pantry, shopping và recipe đổi \texttt{food\_id} của lô, mặt hàng, nguyên liệu sang đích. Các tác vụ này lũy đẳng vì câu \texttt{UPDATE ... WHERE food\_id = <nguồn>} chạy lại không đổi thêm gì; tài liệu Elasticsearch của nguồn bị xóa. Cả hai nhánh ghi \texttt{audit\_logs} trong cùng giao dịch. Nhờ cột \texttt{merged\_into}, lệnh ngoại tuyến còn mang id cũ vẫn được máy chủ chuyển sang thực phẩm đích thay vì nhận 404.



% ---------------------------------------------------------------------
\subsection{A05 – Làm mới phiên}\label{sssec:api-a05}

A05 cấp cặp mã mới khi mã truy cập hết hạn (\br{BR03}); cơ chế xoay vòng, phát hiện dùng lại và cửa sổ ân hạn 30 giây được lập luận ở mục~\ref{sec:bao-mat} (Hình~\ref{fig:seq-dang-nhap-lam-moi}). Dòng phiên được khóa theo băm SHA-256 của mã bằng \texttt{FOR UPDATE}; dùng lại trong ân hạn chỉ thu hồi phiên kế nhiệm (\texttt{superseded}), ngoài ân hạn thì thu hồi cả họ và tăng \texttt{token\_version}. A05 không dùng \texttt{Idempotency-Key} vì cửa sổ ân hạn đã đảm nhiệm việc gửi lại, tránh lưu mã làm mới dạng rõ (Bảng~\ref{tab:c3-api-a05}, Mã nguồn~\ref{code:c3-api-a05}).

\begin{apispec}{Đặc tả và test case dự kiến của API A05 -- Làm mới phiên}{tab:c3-api-a05}{tab:c3-tc-a05}
\apirow{URL, quyền}{\texttt{POST} \path{/it4788/v1/auth/refresh}; công khai, không cần \texttt{Authorization}.}
\apirow{Yêu cầu}{\texttt{refreshToken}: chuỗi base64url 43 ký tự, bắt buộc; \texttt{deviceId}: UUID, bắt buộc, trùng \texttt{device\_id} của phiên.}
\apirow{Phản hồi}{200, 1000: \texttt{accessToken}, \texttt{expiresIn}, \texttt{refreshToken} mới, \texttt{refreshExpiresAt} (30 ngày kể từ lần làm mới).}
\apirow{Mã lỗi}{2001; 3001 mã không tồn tại, hết hạn hoặc đã thu hồi; 3005 tài khoản \texttt{disabled} (khóa tạm \texttt{locked} chỉ chặn A04, không chặn làm mới phiên đang có); 3008 dùng lại mã đã xoay hoặc sai thiết bị.}
\tchead
\tcrow{TC-A05-01}{Làm mới hợp lệ}{Mã còn hạn, chưa xoay, đúng thiết bị}{200, 1000; phiên cũ có \texttt{rotated\_at}; thêm kế nhiệm cùng \texttt{family\_id}}
\tcrow{TC-A05-02}{Sai dữ liệu, biên thời hạn}{(a) Thiếu \texttt{refreshToken}; (b) \texttt{expires\_at} sớm hơn hiện tại 1 giây}{(a) 400, 2001; (b) 401, 3001; không tạo phiên}
\tcrow{TC-A05-03}{Dùng lại ngoài ân hạn}{Gửi lại mã đã xoay 31 giây trước}{401, 3008; cả họ bị thu hồi, \texttt{token\_version} tăng 1; mã kế nhiệm cũng bị từ chối}
\tcrow{TC-A05-04}{Gửi lại trong ân hạn}{Mã đã xoay 5 giây trước, cùng thiết bị, kế nhiệm chưa dùng}{200; kế nhiệm bị thu hồi (\texttt{superseded}); họ còn đúng một phiên hoạt động}
\tcrow{TC-A05-05}{Sai thiết bị, trạng thái tài khoản}{(a) \texttt{deviceId} khác; (b) \texttt{status = disabled}; (c) \texttt{locked} do người khác đoán sai mật khẩu}{(a) 401, 3008, thu hồi cả họ; (b) 403, 3005; (c) 200, phiên của chủ tài khoản không bị cắt}
\end{apispec}

\apiNeed{3.4cm}\codecaption{Ví dụ yêu cầu và phản hồi của A05 khi mã cũ bị dùng lại}{code:c3-api-a05}
\begin{Verbatim}[fontsize=\scriptsize]
POST /it4788/v1/auth/refresh                                 X-Platform: ios
{ "refreshToken": "q3v9Xb7tLk2...wYc8", "deviceId": "0192f3a1-7c2e-..." }
HTTP/1.1 401 Unauthorized                 (mã đã xoay từ 31 giây trước)
{ "code": 3008, "message": "Refresh token reused; session family revoked",
  "traceId": "8e1c0d2a9b3f4e5d6c7b8a9f0e1d2c3b" }
\end{Verbatim}

% ---------------------------------------------------------------------
\subsection{G07 – Tham gia nhóm bằng mã mời}\label{sssec:api-g07}

G07 hiện thực \uc{UC11} (Bảng~\ref{tab:c3-api-g07}). Lần gọi \texttt{dryRun = true} kiểm tra \br{BR06}, \br{BR04}, \br{BR05} theo thứ tự của Hình~\ref{fig:act-tham-gia-nhom} và trả thông tin nhóm để xác nhận. Lần gọi xác nhận kiểm tra lại trong giao dịch ghi theo thứ tự khóa toàn cục: chèn khóa lũy đẳng, tra mã mời để lấy \texttt{space\_id}, lấy khóa tư vấn \path{pg_advisory_xact_lock(hashtext(space_id::text))}, rồi đọc lại trạng thái mã và đếm thành viên, cuối cùng mới tăng \texttt{change\_seq}. Khóa tư vấn này cũng được G05, G06, G08--G10 lấy, nên hai người tranh chỗ cuối, hay một người tham gia đúng lúc trưởng nhóm đổi mã mời, được tuần tự hóa mà không phải khóa dòng \texttt{household.spaces} từ đầu giao dịch. Mã sai, hết hạn hay bị thu hồi nhận cùng \msg{MSG07}; với giới hạn 10 lần/giờ/người (\path{rl:join:<userId>}) trên $32^8$ mã, dò vét cạn là không khả thi (Mã nguồn~\ref{code:c3-api-g07}).

\begin{apispec}{Đặc tả và test case dự kiến của API G07 -- Tham gia nhóm bằng mã mời}{tab:c3-api-g07}{tab:c3-tc-g07}
\apirow{URL, quyền}{\texttt{POST} \path{/it4788/v1/spaces/join}; người dùng đã đăng nhập; 10 lần/giờ/người, tính cả lần gọi \texttt{dryRun}; \texttt{Idempotency-Key} khi \texttt{dryRun = false}.}
\apirow{Yêu cầu}{\texttt{code}: 8 ký tự A--Z, 2--9, không có O, I, 0, 1; máy chủ bỏ khoảng trắng và chuyển chữ in hoa (BR06); \texttt{dryRun}: boolean, mặc định \texttt{false}.}
\apirow{Phản hồi}{200, 1000. Xem trước: \texttt{space}, \texttt{leaderName}, \texttt{memberCount}, \texttt{maxMembers}. Tham gia: \texttt{space}, \texttt{role = "member"}, \texttt{joinedAt}; người đã là thành viên của nhóm nhận lại thông tin hiện có.}
\apirow{Mã lỗi}{2001; 2002 sai định dạng (MSG02); 4004 với \texttt{rule} BR06 (MSG07), BR04 (MSG08, kèm tên nhóm hiện tại) hoặc BR05 (MSG09).}
\tchead
\tcrow{TC-G07-01}{Tham gia hợp lệ}{Mã còn hiệu lực, nhóm 3/10, người dùng chưa thuộc nhóm gia đình}{200, 1000; thêm một thành viên; \texttt{change\_seq} tăng 1; một tác vụ thông báo cho trưởng nhóm}
\tcrow{TC-G07-02}{Xem trước, chuẩn hóa mã}{\texttt{dryRun = true}, mã gõ “ k7mx 9qpa”}{200; tên nhóm, trưởng nhóm, 3/10; CSDL không đổi}
\tcrow{TC-G07-03}{Sai định dạng, vi phạm quy tắc}{(a) Mã 7 hoặc 9 ký tự, có chữ O; (b) \texttt{expires\_at} sớm hơn hiện tại 1 giây; (c) đang ở nhóm khác}{(a) 400, 2002; (b) 422, 4004, \texttt{rule = "BR06"}; (c) 422, 4004, \texttt{rule = "BR04"}, MSG08 có tên nhóm}
\tcrow{TC-G07-04}{Tranh chấp chỗ cuối}{Nhóm 9/10, hai người gửi đồng thời}{Một yêu cầu 200, yêu cầu kia 422, 4004, \texttt{rule = "BR05"}; nhóm có đúng 10 thành viên}
\tcrow{TC-G07-05}{Gửi lại lũy đẳng}{Lặp lại TC-G07-01 với cùng \texttt{Idempotency-Key}}{200, \texttt{Idempotent-Replayed: true}; không thêm thành viên, không thêm tác vụ}
\end{apispec}

\apiNeed{3.4cm}\codecaption{Ví dụ yêu cầu và phản hồi của G07}{code:c3-api-g07}
\begin{Verbatim}[fontsize=\scriptsize]
POST /it4788/v1/spaces/join         Idempotency-Key: 0192f3d0-8a1b-7c2d-...
{ "code": "K7MX9QPA", "dryRun": false }
HTTP/1.1 200 OK
{ "code": 1000, "message": "Đã lưu thay đổi.",
  "data": { "space": { "id": "0192e8b1-...", "name": "Nhà mình" },
            "role": "member", "joinedAt": "2026-10-25T03:20:11Z" } }
\end{Verbatim}

% ---------------------------------------------------------------------
\subsection{S12, S14 – Đánh dấu đã mua và chuyển vào tủ lạnh}\label{sssec:api-s12-s14}

S12, S14 hiện thực \uc{UC20}, \uc{UC21} (Bảng~\ref{tab:c3-api-s12-s14}, Hình~\ref{fig:seq-outbox-mua-hang} và Hình~\ref{fig:seq-nhap-tu}). Hai API được tách vì đánh dấu diễn ra tại chợ, thường khi ngoại tuyến, còn nhập tủ diễn ra ở nhà. Cả hai nhận mảng mặt hàng và trả 200 với \texttt{results[i].outcome} cho từng mặt hàng, trong đó hai kết quả \texttt{already\_purchased}, \texttt{already\_stocked} là ca lũy đẳng chứ không phải lỗi. Lỗi dữ liệu vào 2xxx, lỗi quy tắc 4004 và xung đột 4002 làm hỏng cả lô: \texttt{errors} chỉ rõ \texttt{items[i]} và không mặt hàng nào được ghi.

Mua trùng được giải theo quy tắc lệnh được máy chủ ghi nhận trước thắng. Đây là cách hiện thực “giữ bản ghi sớm hơn” ở bước 6a của \uc{UC20} (\nfr{NFR11}): đồng hồ thiết bị không đáng tin, nên \texttt{purchasedAt} do ứng dụng gửi chỉ để hiển thị và báo cáo. Một câu \texttt{UPDATE ... WHERE status IN ('to\_buy', 'assigned') RETURNING} trên cả lô quyết định ai thắng; mặt hàng không được trả về là đã có người mua trước, nhận \texttt{already\_purchased} kèm \texttt{current}. Nếu không mặt hàng nào được ghi, như khi phát lại một mặt hàng duy nhất, phản hồi là 409, mã 4002 (\msg{MSG35}) kèm \texttt{current.purchasedBy}. Mua thay người được giao cần \texttt{onBehalf = true} sau hộp thoại \msg{MSG12} (\br{BR10}). Danh sách Hoàn tất chặn sửa nội dung (\br{BR08}) trừ ba lệnh: S08 sửa \texttt{actualQuantity}, \texttt{unitPrice} của mặt hàng đã mua, S13 hoàn tác (danh sách trở về Đang mua) và S14 nhập tủ, vì người dùng thường cất hàng sau khi danh sách đã hoàn tất. Mặt hàng đã nhập tủ được trả lại lô cũ (Mã nguồn~\ref{code:c3-api-s14}).

\begin{apispec}{Đặc tả và test case dự kiến của API S12, S14 -- Đánh dấu đã mua, chuyển vào tủ lạnh}{tab:c3-api-s12-s14}{tab:c3-tc-s12-s14}
\apisub{S12 -- Đánh dấu đã mua}
\apirow{URL, quyền}{\texttt{POST} \path{/it4788/v1/spaces/{sid}/shopping-items/purchase}; mọi thành viên (BR10); \texttt{Idempotency-Key}; \texttt{X-Client-Time} khi phát lại ngoại tuyến.}
\apirow{Yêu cầu}{\texttt{items} (1--50): \texttt{itemId}; \texttt{actualQuantity} ($>$ 0, tối đa hai chữ số thập phân, mặc định lượng cần mua); \texttt{unitPrice} (số nguyên $\ge$ 0, đồng); \texttt{onBehalf} (mặc định \texttt{false}); \texttt{purchasedAt} (thời điểm thao tác trên thiết bị, chỉ để hiển thị và báo cáo).}
\apirow{Phản hồi}{200, 1000: \texttt{results} (\texttt{itemId}, \texttt{outcome} là \texttt{purchased} hoặc \texttt{already\_purchased} kèm \texttt{current} gồm \texttt{purchasedBy}, \texttt{purchasedAt}; \texttt{status}, \texttt{version}); \texttt{list} với trạng thái tính lại theo BR08.}
\apirow{Mã lỗi}{2002 (MSG16); 4001; 4002 khi mọi mặt hàng đã được người khác mua (MSG35, \texttt{current.purchasedBy}); 4004 với \texttt{rule = "BR10"} khi mặt hàng đang giao cho người khác mà thiếu \texttt{onBehalf} (MSG12 có tên người được giao), \texttt{rule = "BR08"} khi danh sách đã Hoàn tất.}
\apisub{S14 -- Chuyển vào tủ lạnh}
\apirow{URL, quyền}{\texttt{POST} \path{/it4788/v1/spaces/{sid}/shopping-items/stock}; như S12; nhận cả mặt hàng của danh sách đã Hoàn tất.}
\apirow{Yêu cầu}{\texttt{items} (1--50): \texttt{itemId}; \texttt{foodId} (bắt buộc với mặt hàng tên tự do); \texttt{storageLocation} (\texttt{chill}, \texttt{freezer}, \texttt{crisper}, \texttt{door}, \texttt{pantry\_shelf}); \texttt{quantity}, \texttt{unitId} (mặc định lượng thực tế); \texttt{purchasedOn} (mặc định ngày đánh dấu đã mua); \texttt{expiresOn} ($\ge$ \texttt{purchasedOn}, bỏ trống thì tính theo BR11).}
\apirow{Phản hồi}{200, 1000: \texttt{results} (\texttt{itemId}, \texttt{outcome} là \texttt{stocked} hoặc \texttt{already\_stocked}, \texttt{lotId}, \texttt{expiresOn}, \texttt{expiryState} theo ngưỡng $N$ của người gọi). Lô mới ghi \texttt{source\_item\_id} và \texttt{total\_cost} bằng lượng thực tế nhân đơn giá, nên giá trị lãng phí vẫn đúng khi đơn vị đổi lúc nhập tủ.}
\apirow{Mã lỗi}{2002 ngày hết hạn trước ngày mua (MSG14), số lượng sai (MSG16); 4002 mặt hàng vừa bị hoàn tác, không còn Đã mua (MSG35); 4004 mặt hàng tên tự do chưa chọn thực phẩm (bước 2a, \texttt{field = "items[i].foodId"}).}
\tchead
\tcrow{TC-S12-01}{Đánh dấu hợp lệ}{Hai mặt hàng \texttt{to\_buy}, một mặt hàng có đơn giá}{200; hai \texttt{outcome = purchased}; danh sách chuyển \texttt{shopping}}
\tcrow{TC-S12-02}{Trùng một phần}{Hai mặt hàng, B đã đánh dấu mặt hàng thứ hai trước}{200; \texttt{purchased} và \texttt{already\_purchased} với \texttt{current.purchasedBy} là B; bản ghi của B giữ nguyên}
\tcrow{TC-S12-03}{Xung đột ngoại tuyến}{A phát lại một mặt hàng mà B đã mua trong lúc A mất mạng}{409, 4002 (MSG35 có tên B), \texttt{current.purchasedBy} là B; không ghi}
\tcrow{TC-S12-04}{Mua thay, biên số lượng}{(a) Mặt hàng giao cho C, thiếu \texttt{onBehalf}; (b) \texttt{actualQuantity} bằng 0 hoặc 1,005}{(a) 422, 4004, \texttt{rule = "BR10"}, MSG12 có tên C; (b) 400, 2002; không mặt hàng nào được ghi}
\tcrow{TC-S14-01}{Cất hàng của danh sách Hoàn tất}{Ba mặt hàng \texttt{purchased} của danh sách \texttt{completed}, bỏ trống \texttt{expiresOn}}{200; ba lô có hạn theo hạn dùng gợi ý tại vị trí; ba dòng \texttt{stock\_in}; \texttt{total\_cost} bằng lượng nhân đơn giá}
\tcrow{TC-S14-02}{Biên ngày, tính nguyên tử}{Phần tử thứ hai có \texttt{expiresOn} trước ngày mua 1 ngày}{400, 2002 (MSG14), \texttt{field = "items[1].expiresOn"}; không lô nào được tạo}
\tcrow{TC-S14-03}{Gửi lại với khóa mới}{Mặt hàng đã \texttt{stocked} bởi thiết bị khác}{200, \texttt{outcome = already\_stocked}, \texttt{lotId} cũ; số lô không đổi}
\end{apispec}

\apiNeed{3.4cm}\codecaption{Ví dụ yêu cầu và phản hồi của S14}{code:c3-api-s14}
\begin{Verbatim}[fontsize=\scriptsize]
POST /it4788/v1/spaces/0192e8b1-.../shopping-items/stock
{ "items": [ { "itemId": "0192f101-aa01-...", "storageLocation": "chill" },
    { "itemId": "0192f101-aa02-...", "storageLocation": "freezer",
      "expiresOn": "2026-12-20" } ] }
HTTP/1.1 200 OK
{ "code": 1000, "message": "Đã lưu thay đổi.", "data": { "results": [
  { "itemId": "0192f101-aa01-...", "outcome": "stocked",
    "lotId": "0192f3e2-77a0-...", "expiresOn": "2026-10-28",
    "expiryState": "expiring" },
  { "itemId": "0192f101-aa02-...", "outcome": "already_stocked",
    "lotId": "0192f3e2-77a1-...", "expiresOn": "2026-12-20",
    "expiryState": "fresh" } ] } }
\end{Verbatim}

% ---------------------------------------------------------------------
\subsection{S15 – Sinh danh sách mua sắm từ kế hoạch}\label{sssec:api-s15}

S15 hiện thực \uc{UC22} theo Hình~\ref{fig:seq-sinh-ds-ke-hoach} (Bảng~\ref{tab:c3-api-s15}). Hai pha của use case dùng chung một API. Lần gọi \texttt{dryRun = true} tính phần thiếu theo \br{BR20} trong một giao dịch chỉ đọc; đầu vào là món chưa nấu có công thức trong khoảng \texttt{[from, to]}, lô còn hạn đến ngày nấu và lượng đang chờ mua, gồm mặt hàng Cần mua, Đã phân công, Đã mua chưa nhập tủ, trong đó hàng chưa mua chỉ được tính khi ngày dự kiến mua không sau ngày cần dùng. Mỗi dòng thiếu nêu các món dùng nó và ngày cần sớm nhất để người dùng hiểu vì sao phải mua. Món tên tự do không tính được nên được trả riêng (bước 2a); không thiếu gì thì trả 200 kèm \msg{MSG15} (bước 5a), vì đó là thông tin chứ không phải lỗi.

Lần gọi xác nhận mang các dòng mà người dùng giữ lại hoặc chỉnh ở bước 6. Máy chủ không tính lại phần thiếu, vì đề xuất chỉ mang tính tư vấn, mà kiểm tra dữ liệu rồi giao cổng \texttt{ListComposer} của shopping tạo danh sách \texttt{origin = meal\_plan} hoặc gộp vào danh sách có sẵn theo \br{BR09}, nhờ đó quy tắc gộp chỉ được cài đặt một lần. Danh sách đích được khóa \texttt{FOR UPDATE} trước khi gộp để tuần tự hóa với thành viên đang thêm mặt hàng; danh sách đích ứng viên ở bước 6a lấy bằng S01 (Mã nguồn~\ref{code:c3-api-s15}).

\begin{apispec}{Đặc tả và test case dự kiến của API S15 -- Sinh danh sách mua sắm từ kế hoạch}{tab:c3-api-s15}{tab:c3-tc-s15}
\apirow{URL, quyền}{\texttt{POST} \path{/it4788/v1/spaces/{sid}/shopping-lists/generate}; thành viên không gian; \texttt{Idempotency-Key} khi \texttt{dryRun = false}; \texttt{X-Client-Time} khi phát lại ngoại tuyến.}
\apirow{Yêu cầu}{\texttt{from}, \texttt{to}: ngày, bắt buộc, hôm nay $\le$ \texttt{from} $\le$ \texttt{to} $\le$ hôm nay + 28 (BR18); \texttt{dryRun}: boolean, mặc định \texttt{false}. Khi xác nhận: \texttt{lines} (1--100) gồm \texttt{foodId}, \texttt{quantity} ($>$ 0, tối đa hai chữ số thập phân), \texttt{unitId}; \texttt{name} (1--100 ký tự, mặc định theo ngày mua); đúng một trong hai \texttt{plannedDate} ($\ge$ hôm nay, BR08) để tạo mới hoặc \texttt{targetListId} để gộp (bước 6a).}
\apirow{Phản hồi}{Xem trước: 200, 1000: \texttt{lines} (\texttt{foodId}, \texttt{name}, \texttt{shortQty}, \texttt{unitId}, \texttt{neededBy}, \texttt{meals}); \texttt{skippedMeals} (món tên tự do); \texttt{suggestedPlannedDate} là ngày liền trước \texttt{neededBy} sớm nhất nhưng không trước hôm nay; \texttt{lines} rỗng thì kèm MSG15. Xác nhận: 201, 1000: \texttt{listId}, \texttt{createdItemIds}, \texttt{merged} (\texttt{itemId}, lượng sau gộp; kèm MSG11).}
\apirow{Mã lỗi}{2001 thiếu \texttt{from}, \texttt{to}, hoặc \texttt{lines} khi xác nhận (MSG01); 2002 khoảng ngày sai, \texttt{plannedDate} ở quá khứ, có cả hoặc không có \texttt{plannedDate} và \texttt{targetListId} (MSG02), số lượng sai (MSG16); 4001 \texttt{foodId} hoặc \texttt{targetListId} ngoài phạm vi không gian; 4004 \texttt{rule = "BR08"} khi danh sách đích đã Hoàn tất.}
\tchead
\tcrow{TC-S15-01}{Ví dụ kiểm chứng chương 2}{D là ngày mai; bữa D cần 0,4 kg cà rốt, bữa D+2 cần 0,6 kg; lô A 0,5 kg hết hạn D+1, lô B 0,2 kg hết hạn D+3; 0,1 kg chờ mua hôm nay; \texttt{dryRun = true}}{200; một dòng cà rốt \texttt{shortQty} 0,3 kg, \texttt{neededBy} D+2, \texttt{meals} là bữa D+2; \texttt{suggestedPlannedDate} D+1; CSDL không đổi}
\tcrow{TC-S15-02}{Không thiếu, món tự do}{Tủ đủ mọi nguyên liệu; một món tên tự do trong khoảng}{200, 1000; \texttt{lines} rỗng, MSG15; \texttt{skippedMeals} có món tự do (bước 2a, 5a)}
\tcrow{TC-S15-03}{Biên khoảng ngày, chọn đích}{(a) \texttt{to} bằng hôm nay + 28, đối chứng + 29; (b) \texttt{from} là hôm qua hoặc sau \texttt{to}; (c) có cả \texttt{plannedDate} và \texttt{targetListId}}{(a) 200, đối chứng 400, 2002; (b), (c) 400, 2002, \texttt{field} đúng tham số sai; CSDL không đổi}
\tcrow{TC-S15-04}{Xác nhận tạo mới}{Bỏ một dòng, sửa cà rốt thành 0,5 kg; \texttt{plannedDate} hôm nay}{201; một danh sách \texttt{meal\_plan}, \texttt{drafting}; mặt hàng \texttt{to\_buy} đúng các dòng gửi lên; \texttt{change\_seq} tăng}
\tcrow{TC-S15-05}{Gộp vào danh sách có sẵn}{\texttt{targetListId} là danh sách Đang lập có 2 lạng cà rốt; dòng 0,3 kg cà rốt}{201; mặt hàng cũ thành 5 lạng (BR09) và có trong \texttt{merged}, MSG11; không tạo danh sách mới}
\tcrow{TC-S15-06}{Đích đã Hoàn tất}{\texttt{targetListId} là danh sách \texttt{completed}}{422, 4004, \texttt{rule = "BR08"}; không mặt hàng nào được thêm}
\tcrow{TC-S15-07}{Gửi lại lũy đẳng}{(a) Lặp lại TC-S15-04 cùng \texttt{Idempotency-Key}; (b) cùng khóa nhưng \texttt{lines} khác}{(a) 201, \texttt{Idempotent-Replayed: true}, không có danh sách thứ hai; (b) 400, 2002}
\end{apispec}

\apiNeed{3.4cm}\codecaption{Ví dụ lần gọi xem trước và lần gọi xác nhận của S15}{code:c3-api-s15}
\begin{Verbatim}[fontsize=\scriptsize]
POST /it4788/v1/spaces/0192e8b1-.../shopping-lists/generate
{ "from": "2026-10-26", "to": "2026-11-01", "dryRun": true }
HTTP/1.1 200 OK
{ "code": 1000, "message": "Thành công", "data": {
  "lines": [ { "foodId": "0192a0d3-...", "name": "Cà rốt", "shortQty": 0.3,
      "unitId": "<kg>", "neededBy": "2026-10-28",
      "meals": [ { "mealId": "0192f5c1-...", "mealDate": "2026-10-28",
                   "recipeName": "Canh rau củ" } ] } ],
  "skippedMeals": [ { "mealId": "0192f5c7-...", "name": "Phở bò" } ],
  "suggestedPlannedDate": "2026-10-27" } }
POST /it4788/v1/spaces/0192e8b1-.../shopping-lists/generate
Idempotency-Key: 0192f6a0-3b1c-7d2e-...
{ "from": "2026-10-26", "to": "2026-11-01", "dryRun": false,
  "lines": [ { "foodId": "0192a0d3-...", "quantity": 0.5, "unitId": "<kg>" }],
  "name": "Đi chợ cuối tuần", "plannedDate": "2026-10-27" }
HTTP/1.1 201 Created
{ "code": 1000, "message": "Đã lưu thay đổi.", "data": {
  "listId": "0192f6a1-...", "createdItemIds": [ "0192f6a2-..." ],
  "merged": [] } }
\end{Verbatim}

% ---------------------------------------------------------------------
\subsection{P05 – Ghi nhận sử dụng thực phẩm}\label{sssec:api-p05}

P05 hiện thực \uc{UC28} (Bảng~\ref{tab:c3-api-p05}). Với \texttt{strategy = fefo}, lô trong đường dẫn chỉ xác định thực phẩm; máy chủ khóa mọi lô còn hạn (\texttt{expires\_on >= \$today}) của thực phẩm đó rồi trừ dần theo ngày hết hạn (\br{BR15}), quy đổi qua đơn vị cơ sở (\br{BR16}). \texttt{strategy = lot} chỉ trừ lô được chọn và là cách duy nhất để trừ lô quá hạn. Mỗi lô bị trừ có một dòng \texttt{stock\_movements} loại \texttt{consume}, lô về 0 chuyển \texttt{consumed} (\br{BR14}). Lệnh là phép trừ tương đối nên không cần \texttt{If-Match} (Mã nguồn~\ref{code:c3-api-p05}).

\begin{apispec}{Đặc tả và test case dự kiến của API P05 -- Ghi nhận sử dụng thực phẩm}{tab:c3-api-p05}{tab:c3-tc-p05}
\apirow{URL, quyền}{\texttt{POST} \path{/it4788/v1/spaces/{sid}/lots/{lotId}/consume}; thành viên không gian; \texttt{Idempotency-Key}.}
\apirow{Yêu cầu}{\texttt{quantity}: $>$ 0, tối đa hai chữ số thập phân, bắt buộc; \texttt{unitId}: bắt buộc, cùng nhóm đại lượng với lô (BR16); \texttt{strategy}: \texttt{fefo} hoặc \texttt{lot}, mặc định \texttt{fefo}.}
\apirow{Phản hồi}{200, 1000: \texttt{foodId}; \texttt{consumed} (\texttt{lotId}, lượng đã trừ, \texttt{remainingQty}, \texttt{lifecycle}, \texttt{version}); \texttt{totalRemaining}.}
\apirow{Mã lỗi}{2001; 2002 (MSG16, hoặc \texttt{rule = "BR16"}); 2003 vượt lượng còn lại (MSG17); 4001; 4004 \texttt{rule = "BR15"} khi dùng \texttt{fefo} mà chỉ còn lô quá hạn, ứng dụng hỏi người dùng theo bước 4b.}
\tchead
\tcrow{TC-P05-01}{Trừ FEFO qua hai lô}{Lô A 0,3 kg hết hạn D+1, lô B 1 kg hết hạn D+5; dùng 500 g}{200; A về 0, \texttt{consumed}; B còn 0,8 kg; hai dòng \texttt{consume}}
\tcrow{TC-P05-02}{Biên trên của lượng dùng}{Như TC-P05-01; dùng 1,3 kg; dùng 1,31 kg}{1,3 kg: 200, hai lô \texttt{consumed}; 1,31 kg: 422, 2003 (MSG17 kèm 1,3 kg), không lô nào đổi}
\tcrow{TC-P05-03}{Dữ liệu sai}{\texttt{quantity} bằng 0 hoặc 0,005; đơn vị “quả” cho lô tính theo kg}{400, 2002; đơn vị sai kèm \texttt{rule = "BR16"}}
\tcrow{TC-P05-04}{Chỉ còn lô quá hạn}{Lần 1 \texttt{fefo}; lần 2 \texttt{lot} trên lô quá hạn}{Lần 1: 422, 4004, \texttt{rule = "BR15"}, không ghi; lần 2: 200, chỉ lô được chọn bị trừ}
\tcrow{TC-P05-05}{Đồng thời, lũy đẳng}{(a) Tổng còn 1,3 kg, hai yêu cầu 0,8 kg cùng lúc; (b) gửi lại TC-P05-01 cùng khóa}{(a) Một 200, một 422, 2003, không có lượng âm; (b) 200, \texttt{Idempotent-Replayed: true}, tồn kho không đổi}
\end{apispec}

\apiNeed{3.4cm}\codecaption{Ví dụ yêu cầu và phản hồi của P05}{code:c3-api-p05}
\begin{Verbatim}[fontsize=\scriptsize]
POST /it4788/v1/spaces/0192e8b1-.../lots/0192f3e2-77a0-.../consume
{ "quantity": 500, "unitId": "<g>", "strategy": "fefo" }
HTTP/1.1 200 OK
{ "code": 1000, "message": "Đã lưu thay đổi.",
  "data": { "foodId": "0192a011-...", "consumed": [
    { "lotId": "0192f3e2-77a0-...", "quantity": 0.3, "remainingQty": 0,
      "lifecycle": "consumed" },
    { "lotId": "0192f3f9-1c20-...", "quantity": 0.2, "remainingQty": 0.8,
      "lifecycle": "active" } ], "totalRemaining": { "quantity": 800 } } }
\end{Verbatim}

% ---------------------------------------------------------------------
\subsection{M05 – Gợi ý món ăn}\label{sssec:api-m05}

M05 hiện thực \uc{UC40} bằng thuật toán ở mục~\ref{sec:thiet-ke-chi-tiet} (Bảng~\ref{tab:c3-api-m05}). Chỉ công thức \texttt{published} được xét, gồm công thức của không gian và công thức mẫu đã công bố. Kết quả phụ thuộc ngưỡng $N$ và danh sách yêu thích của người gọi, nên bộ đệm 10 phút là một hash Redis \texttt{suggest:\{sid\}} với mỗi trường ứng với một người xem và bộ tham số. Khóa bị xóa khi có \texttt{lot.changed}, \texttt{recipe.changed}, \texttt{config.changed} với tham số \texttt{suggest.*}, hoặc khi một thành viên đổi $N$ qua A13, nên một lệnh \texttt{DEL} làm mới kết quả cho mọi người xem. Không có món đạt chuẩn vẫn trả 200: tủ trống trả công thức mẫu, không món nào đạt $\theta$ trả ba món có $R$ cao nhất kèm \msg{MSG31}, \texttt{meta.fallback} cho biết nhánh nào. Mỗi món kèm $R$, $E$, $F$ để kiểm thử đối chiếu với \br{BR19} (Mã nguồn~\ref{code:c3-api-m05}); mục tiêu 2 giây (\nfr{NFR09}).

\begin{apispec}{Đặc tả và test case dự kiến của API M05 -- Gợi ý món ăn}{tab:c3-api-m05}{tab:c3-tc-m05}
\apirow{URL, quyền}{\texttt{GET} \path{/it4788/v1/spaces/{sid}/recipe-suggestions}; thành viên không gian.}
\apirow{Yêu cầu}{\texttt{date}: ngày nấu dự kiến, mặc định hôm nay, trong cửa sổ của BR18; lô còn hạn và $E$ được tính tại ngày này (UC37); \texttt{mealSlot}: \texttt{breakfast}, \texttt{lunch}, \texttt{dinner}, tùy chọn; \texttt{expiringOnly}: boolean, mặc định \texttt{false}, chỉ giữ món dùng ít nhất một lô sắp hết hạn (UC30 bước 3); \texttt{servings}: 1--20, mặc định bằng số thành viên; \texttt{limit}: 1 đến \texttt{suggest.maxResults} (mặc định 10).}
\apirow{Phản hồi}{200, 1000: \texttt{items} giảm dần theo \texttt{score} (\texttt{recipe}, \texttt{score}, \texttt{r}, \texttt{e}, \texttt{f}, \texttt{missing}, \texttt{expiringUsed}); \texttt{meta.fallback} là \texttt{none}, \texttt{empty\_pantry} hoặc \texttt{below\_threshold} (kèm MSG31).}
\apirow{Mã lỗi}{2002 tham số ngoài miền (MSG02; \texttt{date} ngoài cửa sổ kèm MSG28); 4001 không phải thành viên.}
\tchead
\tcrow{TC-M05-01}{Ví dụ kiểm chứng chương 2}{Cá lóc còn 1 ngày, dứa 2, cà chua 5, trứng 10; $N = 3$; không có me}{200; canh chua cá lóc ($S = 0{,}75$) đứng trước trứng xào cà chua ($0{,}6$); \texttt{missing} là [Me]}
\tcrow{TC-M05-02}{Biên ngưỡng $\theta$}{Công thức A có $R = 0{,}6$ (3/5), B có $R = 0{,}5$}{200; A có mặt, B bị loại}
\tcrow{TC-M05-03}{Hai nhánh dự phòng}{(a) Mọi công thức có $R < 0{,}6$; (b) không có lô \texttt{active}}{(a) \texttt{below\_threshold}, MSG31, đúng 3 món; (b) \texttt{empty\_pantry}, công thức mẫu phổ biến}
\tcrow{TC-M05-04}{Tham số ngoài miền}{\texttt{servings} bằng 0, 21; \texttt{mealSlot = snack}; \texttt{limit = 11}; \texttt{date} bằng hôm nay + 29; đối chứng \texttt{servings} bằng 1, 20}{400, 2002, \texttt{field} đúng tham số sai; đối chứng: 200}
\tcrow{TC-M05-05}{Khác người xem, phân quyền}{(a) Hai thành viên có $N = 1$ và $N = 7$; (b) người ngoài không gian}{(a) $E$ khác nhau theo $N$, mỗi người một trường bộ đệm riêng; (b) 404, 4001}
\tcrow{TC-M05-06}{Chỉ món dùng lô sắp hết hạn, làm mới bộ đệm}{(a) \texttt{expiringOnly = true}, \texttt{limit = 3} như UC30; (b) gọi lại sau khi P05 dùng hết lô cá lóc}{(a) Tối đa 3 món, mọi món có \texttt{expiringUsed} khác rỗng; (b) kết quả được tính lại, canh chua cá lóc không còn $E = 1$}
\end{apispec}

\apiNeed{3.4cm}\codecaption{Ví dụ yêu cầu và phản hồi của M05}{code:c3-api-m05}
\begin{Verbatim}[fontsize=\scriptsize]
GET /it4788/v1/spaces/0192e8b1-.../recipe-suggestions
    ?date=2026-10-25&mealSlot=dinner&servings=4
HTTP/1.1 200 OK
{ "code": 1000, "message": "Thành công", "data": { "items": [
  { "recipe": { "id": "0192c3d4-...", "name": "Canh chua cá lóc" },
    "score": 0.75, "r": 0.75, "e": 1.0, "f": 0,
    "missing": [ { "foodId": "0192a0f1-...", "name": "Me" } ],
    "expiringUsed": [ { "name": "Cá lóc", "daysLeft": 1 },
                      { "name": "Dứa", "daysLeft": 2 } ] },
  { "recipe": { "id": "0192c3e5-...", "name": "Trứng xào cà chua" },
    "score": 0.6, "r": 1.0, "e": 0, "f": 0, "missing": [] } ] },
  "meta": { "fallback": "none", "theta": 0.6 } }
\end{Verbatim}

% ---------------------------------------------------------------------
\subsection{Y01 – Lấy thay đổi đồng bộ}\label{sssec:api-y01}

Y01 là đường dữ liệu duy nhất để thiết bị bắt kịp thay đổi, cho cả đồng bộ ngoại tuyến (\nfr{NFR11}) lẫn gợi ý thời gian thực (\nfr{NFR10}) (Bảng~\ref{tab:c3-api-y01}). \texttt{seq} là bộ đếm riêng từng không gian, tăng theo thứ tự commit, nên thiết bị chỉ lưu một số nguyên; phần tử mang ảnh chụp thực thể nên áp lại là vô hại (Mã nguồn~\ref{code:c3-api-y01}). Con trỏ cũ hơn 90 ngày nhận 410, mã 4100; thiết bị khi đó đọc \texttt{changeSeq} từ G03, tải lại toàn bộ rồi tiếp tục từ số đó.

\begin{apispec}{Đặc tả và test case dự kiến của API Y01 -- Lấy thay đổi đồng bộ}{tab:c3-api-y01}{tab:c3-tc-y01}
\apirow{URL, quyền}{\texttt{GET} \path{/it4788/v1/spaces/{sid}/changes}; thành viên không gian, kiểm tra lại ở mỗi lần gọi.}
\apirow{Yêu cầu}{\texttt{since}: số nguyên $\ge$ 0, bắt buộc, không lớn hơn \texttt{change\_seq} hiện tại; \texttt{limit}: 1--500, mặc định 500.}
\apirow{Phản hồi}{200, 1000: \texttt{changes} tăng theo \texttt{seq}, mỗi phần tử gồm \texttt{seq}, \texttt{entityType}, \texttt{op} (\texttt{upsert}, \texttt{delete}), \texttt{entityId}, \texttt{version}, \texttt{entity} (vắng mặt khi xóa); \texttt{nextSince}; \texttt{hasMore}.}
\apirow{Mã lỗi}{2001; 2002 \texttt{since} âm hoặc vượt \texttt{change\_seq}, \texttt{limit} ngoài 1--500; 4001; 4100 \texttt{since} cũ hơn 90 ngày.}
\tchead
\tcrow{TC-Y01-01}{Gộp thay đổi, bia mộ}{\texttt{since = 1840}; 1841 và 1843 cùng một mặt hàng; 1845 là xóa mềm}{200; mặt hàng xuất hiện một lần với \texttt{seq} 1843; 1845 có \texttt{op = delete}, không có \texttt{entity}; \texttt{nextSince = 1845}}
\tcrow{TC-Y01-02}{Biên phân trang}{501 thay đổi sau \texttt{since}, \texttt{limit} mặc định}{200; 500 phần tử, \texttt{hasMore = true}; lần gọi tiếp trả đúng 1 phần tử}
\tcrow{TC-Y01-03}{Không đổi, ngoài miền}{(a) \texttt{since} bằng \texttt{change\_seq}; (b) \texttt{since} bằng $-$1 hoặc \texttt{change\_seq} + 1; \texttt{limit} bằng 0, 501}{(a) 200, \texttt{changes} rỗng, \texttt{nextSince} bằng \texttt{since}; (b) 400, 2002}
\tcrow{TC-Y01-04}{Con trỏ quá cũ}{\texttt{since} nhỏ hơn \texttt{seq} nhỏ nhất còn giữ}{410, 4100; không trả dữ liệu}
\tcrow{TC-Y01-05}{Mất tư cách thành viên}{Người dùng vừa bị xóa khỏi nhóm (UC13)}{404, 4001, kể cả khi \texttt{since} hợp lệ}
\end{apispec}

\apiNeed{3.4cm}\codecaption{Ví dụ yêu cầu và phản hồi của Y01}{code:c3-api-y01}
\begin{Verbatim}[fontsize=\scriptsize]
GET /it4788/v1/spaces/0192e8b1-.../changes?since=1840&limit=500
HTTP/1.1 200 OK
{ "code": 1000, "message": "Thành công", "data": { "changes": [
    { "seq": 1843, "entityType": "shopping_item", "op": "upsert",
      "entityId": "01929f4e-6c1a-...", "version": 5,
      "entity": { "status": "assigned", "assigneeId": "0192e7c0-..." } },
    { "seq": 1845, "entityType": "shopping_item", "op": "delete",
      "entityId": "01929f4e-80b4-...", "version": 3 } ],
  "nextSince": 1845, "hasMore": false } }
\end{Verbatim}

% =====================================================================

% ---------------------------------------------------------------------
%  PHỤ LỤC. ĐẶC TẢ CHI TIẾT CÁC BẢNG CƠ SỞ DỮ LIỆU (tiếp mục Thiết kế CSDL)
% ---------------------------------------------------------------------
\section{Đặc tả chi tiết các bảng cơ sở dữ liệu}\label{pl:dac-ta-bang}

Bảng~\ref{tab:c3-dac-ta-bang} đặc tả các bảng trọng tâm (PK, FK, NN, UQ, CK, DF: khóa chính, khóa ngoại, NOT NULL, duy nhất, CHECK, mặc định). Cột chung không lặp lại: \texttt{id} uuid PK; \texttt{created\_at}, \texttt{updated\_at} timestamptz NN DF \texttt{now()}; \texttt{version} int NN DF 1; \texttt{deleted\_at} timestamptz. \texttt{pantry.food\_lots} được đặc tả bằng Mã nguồn~\ref{code:c3-ddl-food-lots}, các bảng còn lại ở Bảng~\ref{tab:c3-bang-con-lai}.

{\footnotesize\setlength{\tabcolsep}{4pt}\renewcommand{\arraystretch}{1.1}
\begin{longtable}{|L{4.85cm}|L{2.5cm}|L{1.6cm}|L{\dimexpr\textwidth-8\tabcolsep-5\arrayrulewidth-4.85cm-2.5cm-1.6cm\relax}|}
\caption{Đặc tả các bảng trọng tâm}\label{tab:c3-dac-ta-bang}\\
\hline
\hd{Tên cột} & \hd{Kiểu dữ liệu} & \hd{Ràng buộc} & \hd{Mô tả}\\ \hline
\endfirsthead
\multicolumn{4}{l}{\footnotesize\itshape Bảng \thetable\ (tiếp theo)}\\ \hline
\hd{Tên cột} & \hd{Kiểu dữ liệu} & \hd{Ràng buộc} & \hd{Mô tả}\\ \hline
\endhead
\multicolumn{4}{|l|}{\hs{\texttt{identity.users} -- cột chung trừ \texttt{version}}}\\ \hline
\path{email} & citext & NN, UQ, CK & \texttt{uq\_users\_email}; tối đa 100 ký tự, không phân biệt hoa thường (\br{BR01})\\ \hline
\path{password_hash} & text & NN & Băm Argon2id kèm tham số (\nfr{NFR14})\\ \hline
\path{full_name} & varchar(50) & NN, CK & 2--50 ký tự (\br{BR01})\\ \hline
\path{system_role} & enum & NN, DF & \texttt{user}, \texttt{admin}\\ \hline
\path{status} & enum & NN, DF & \texttt{pending} (chờ OTP), \texttt{active}, \texttt{locked} (khóa tạm), \texttt{disabled} (quản trị viên khóa)\\ \hline
\path{failed_login_count} & smallint & NN, DF, CK & Số lần sai liên tiếp, 0--5 (\br{BR03})\\ \hline
\path{locked_until}, \path{lock_reason} & timestamptz; varchar(200) & -- & Hết khóa tạm 15 phút (\br{BR03}); lý do quản trị viên khóa (\uc{UC24})\\ \hline
\path{last_login_at} & timestamptz & -- & Lần đăng nhập thành công gần nhất (\uc{UC24})\\ \hline
\path{locale}, \path{avatar_key} & enum; text & NN; -- & \texttt{vi} hoặc \texttt{en}; khóa ảnh trong MinIO\\ \hline
\path{last_space_id} & uuid & -- & Không gian mở khi đăng nhập\\ \hline
\path{token_version} & int & NN, DF 0 & Tăng khi đăng xuất, đổi mật khẩu, bị khóa\\ \hline
\multicolumn{4}{|l|}{\hs{\texttt{identity.refresh\_sessions} -- cột chung \texttt{id}, \texttt{created\_at}}}\\ \hline
\path{user_id}, \path{family_id} & uuid & NN, FK & Chủ phiên; họ phiên (thu hồi cả họ)\\ \hline
\path{token_hash} & bytea & NN, UQ & SHA-256 của refresh token\\ \hline
\path{device_id}, \path{platform} & uuid; enum & NN & Bản cài đặt; \texttt{android} hoặc \texttt{ios}\\ \hline
\path{user_agent}, \path{ip} & text; inet & -- & Phục vụ quản lý phiên, điều tra\\ \hline
\path{expires_at}, \path{rotated_at} & timestamptz & NN; -- & Hết hạn 30 ngày; lúc xoay vòng\\ \hline
\path{revoked_at}, \path{revoke_reason} & timestamptz; text & -- & Thời điểm, lý do thu hồi\\ \hline
\multicolumn{4}{|l|}{\hs{\texttt{household.spaces} -- cột chung trừ \texttt{deleted\_at}; UQ (\texttt{id}, \texttt{kind})}}\\ \hline
\path{kind}, \path{name} & enum; varchar(50) & NN; CK & \texttt{personal}, \texttt{family} (\br{BR04}); 3--50 ký tự\\ \hline
\path{change_seq} & bigint & NN, DF 0 & Bộ đếm thay đổi của không gian\\ \hline
\path{avatar_key}, \path{created_by} & text; uuid & --; FK & Ảnh nhóm; người tạo\\ \hline
\path{dissolved_at} & timestamptz & -- & Giải tán nhóm (\br{BR07}), thay xóa mềm\\ \hline
\multicolumn{4}{|l|}{\hs{\texttt{household.space\_members} -- không dùng cột chung}}\\ \hline
\path{space_id}, \path{user_id} & uuid & PK, FK & Khóa chính phức hợp\\ \hline
\path{space_kind} & enum & NN, FK & Sao chép \texttt{kind}; FK ghép tới (\texttt{id}, \texttt{kind})\\ \hline
\path{role}, \path{joined_at} & enum; timestamptz & NN & \texttt{leader}, \texttt{member} (\br{BR05}); ngày vào\\ \hline
\multicolumn{4}{|l|}{\hs{\texttt{shopping.shopping\_items} -- đủ cột chung}}\\ \hline
\path{list_id}, \path{space_id} & uuid & NN, FK & Danh sách; không gian (lọc quyền)\\ \hline
\path{food_id}, \path{custom_name} & uuid; varchar(100) & FK; CK & Thực phẩm hoặc tên tự do\\ \hline
\path{quantity}, \path{unit_id} & numeric(12,2); uuid & NN, CK; FK & Lượng $>$ 0 (\br{BR14}); đơn vị (\br{BR16})\\ \hline
\path{status} & enum & NN, DF & \texttt{to\_buy}, \texttt{assigned}, \texttt{purchased}, \texttt{stocked}, \texttt{skipped}\\ \hline
\path{assignee_id}, \path{assigned_by} & uuid & FK & Người được giao, tối đa một (\br{BR10})\\ \hline
\path{purchased_by}, \path{purchased_at} & uuid; timestamptz & FK; CK & Bắt buộc khi đã mua, đã nhập tủ; người được máy chủ ghi nhận trước\\ \hline
\path{actual_quantity} & numeric(12,2) & CK & Lượng thực mua, $>$ 0\\ \hline
\path{unit_price} & bigint & CK & Đơn giá thực mua (đồng), $\ge$ 0; nhân lượng thực mua thành \texttt{total\_cost} của lô (\br{BR21})\\ \hline
\path{lot_id} & uuid & -- & Lô tạo ra khi nhập tủ (logic)\\ \hline
\path{note}, \path{position} & varchar(200); int & --; NN & Ghi chú; thứ tự hiển thị\\ \hline
\multicolumn{4}{|l|}{\hs{\texttt{pantry.stock\_movements} -- cột chung \texttt{id}; chỉ ghi thêm}}\\ \hline
\path{lot_id}, \path{space_id} & uuid & NN, FK & Lô; không gian (cho báo cáo)\\ \hline
\path{kind} & enum & NN & \texttt{stock\_in}, \texttt{consume}, \texttt{cook}, \texttt{discard}, \texttt{adjust}\\ \hline
\path{quantity}, \path{unit_id} & numeric(12,2); uuid & NN, CK; FK & Theo đơn vị lô, $\neq$ 0; dương trừ \texttt{adjust} (chênh lệch có dấu khi P04 sửa lượng còn)\\ \hline
\path{reason} & enum & CK & \texttt{expired}, \texttt{spoiled}, \texttt{other}; bắt buộc khi loại bỏ\\ \hline
\path{meal_entry_id} & uuid & CK & Bắt buộc khi \texttt{cook} (\uc{UC19}); logic\\ \hline
\path{actor_id}, \path{occurred_at} & uuid; timestamptz & NN, FK; DF & Người thực hiện; thời điểm\\ \hline
\multicolumn{4}{|l|}{\hs{\texttt{mealplan.meal\_entries} -- đủ cột chung}}\\ \hline
\path{space_id}, \path{meal_date} & uuid; date & NN, FK; NN & Không gian; ngày ăn\\ \hline
\path{meal_slot}, \path{status} & enum & NN & \texttt{breakfast}, \texttt{lunch}, \texttt{dinner}; \texttt{planned}, \texttt{cooked}\\ \hline
\path{recipe_id}, \path{custom_name} & uuid; varchar(100) & CK & Công thức (logic) hoặc tên tự do\\ \hline
\path{servings}, \path{position} & smallint & NN, CK & 1--20 khẩu phần; thứ tự trong bữa\\ \hline
\path{cooked_at}, \path{cooked_by} & timestamptz; uuid & CK; FK & Bắt buộc khi đã nấu (\uc{UC19}, \br{BR18})\\ \hline
\path{created_by} & uuid & NN, FK & Người lập kế hoạch\\ \hline
\multicolumn{4}{|l|}{\hs{\texttt{sync.space\_changes} -- không dùng cột chung; giữ 90 ngày}}\\ \hline
\path{space_id}, \path{seq} & uuid; bigint & PK, FK & \texttt{seq} lấy từ \texttt{spaces.change\_seq}\\ \hline
\path{entity_type}, \path{entity_id} & varchar(40); uuid & NN & Loại, định danh thực thể thay đổi\\ \hline
\path{op}, \path{entity_version} & enum; int & NN & \texttt{upsert} hoặc \texttt{delete}; phiên bản mới\\ \hline
\path{changed_by}, \path{changed_at} & uuid; timestamptz & NN, FK; DF & Người thay đổi; thời điểm (cũng là nguồn hoạt động gần đây của G03)\\ \hline
\multicolumn{4}{|l|}{\hs{\texttt{sync.idempotency\_keys} -- không dùng cột chung; giữ 7 ngày}}\\ \hline
\path{user_id}, \path{key} & uuid & PK, FK & Người dùng; khóa lũy đẳng\\ \hline
\path{request_hash} & bytea & NN & SHA-256 của yêu cầu; khác băm thì trả 400 mã 2002\\ \hline
\path{status} & enum & NN, DF & \texttt{in\_progress} khi chèn; \texttt{completed} ngay trước commit\\ \hline
\path{status_code}, \path{response_body} & smallint; jsonb & CK & Bắt buộc khi \texttt{completed}; phản hồi phát lại nguyên văn\\ \hline
\path{created_at}, \path{expires_at} & timestamptz & NN & Hết hạn sau 7 ngày\\ \hline
\end{longtable}
}

{\footnotesize\setlength{\tabcolsep}{4pt}\renewcommand{\arraystretch}{1.1}
\begin{longtable}{|L{3.9cm}|L{2.8cm}|L{3.3cm}|L{\dimexpr\textwidth-8\tabcolsep-5\arrayrulewidth-3.9cm-2.8cm-3.3cm\relax}|}
\caption{Tóm tắt các bảng còn lại}\label{tab:c3-bang-con-lai}\\
\hline
\hd{Bảng} & \hd{Mục đích} & \hd{Cột chính} & \hd{Ràng buộc đáng chú ý}\\ \hline
\endfirsthead
\multicolumn{4}{l}{\footnotesize\itshape Bảng \thetable\ (tiếp theo)}\\ \hline
\hd{Bảng} & \hd{Mục đích} & \hd{Cột chính} & \hd{Ràng buộc đáng chú ý}\\ \hline
\endhead
\path{user_settings} & Thiết lập nhắc nhở & \path{reminder_days} & CK 1--7, rỗng thì dùng cấu hình (\br{BR12})\\ \hline
\path{otp_challenges} & Mã OTP, mã đặt lại mật khẩu & \path{code_hash}, \path{reset_token_hash} & Sai $\le$ 5 lần (\br{BR02}); mã đặt lại UQ, 10 phút, dùng một lần\\ \hline
\path{devices} & Thiết bị nhận đẩy & \path{fcm_token} & UQ mã FCM; \texttt{id} là khóa của A15\\ \hline
\path{invite_codes} & Mã mời & \path{space_id}, \path{code} & UQ từng phần (\br{BR06})\\ \hline
\path{food_categories} & Loại thực phẩm & \path{default_storage} & UQ \texttt{lower(name)} (\br{BR22})\\ \hline
\path{units} & Đơn vị đo & \path{factor_to_base}, \path{is_base} & CK hệ số $>$ 0; một đơn vị cơ sở mỗi nhóm (\br{BR16})\\ \hline
\path{foods} & Thực phẩm chung và riêng & \path{scope}, \path{merged_into} & UQ từng phần theo \texttt{scope} (\br{BR22}); CK gộp thì \texttt{inactive}\\ \hline
\path{food_shelf_lives} & Hạn dùng gợi ý & \path{storage_location} & PK phức hợp\\ \hline
\path{shopping_lists} & Danh sách mua sắm & \path{planned_date}, \path{status} & \texttt{status} theo \br{BR08}\\ \hline
\path{recipes} & Công thức & \path{kind}, \path{status}, \path{steps} & CK \texttt{space\_id} theo \texttt{kind}; công thức không gian luôn \texttt{published}\\ \hline
\path{recipe_ingredients} & Nguyên liệu & \path{food_id}, \path{required} & PK phức hợp\\ \hline
\path{favorite_recipes} & Công thức yêu thích & \path{user_id}, \path{recipe_id} & PK phức hợp\\ \hline
\path{notifications} & Hộp thư thông báo & \path{dedupe_key} & UQ khóa chống trùng (\br{BR13})\\ \hline
\path{deliveries} & Lần gửi đẩy, thư & \path{device_id}, \path{status} & UQ theo thông báo, thiết bị, kênh\\ \hline
\path{system_configs} & Tham số (\uc{UC25}) & \path{key}, \path{value} & PK \path{key}; miền ở Bảng~\ref{tab:c3-tham-so}\\ \hline
\path{audit_logs} & Nhật ký quản trị & \path{old_value}, \path{new_value} & Chỉ ghi thêm\\ \hline
\path{objects} & Ảnh tải lên (Y03, Y04) & \path{purpose}, \path{target_id} & \texttt{pending}, \texttt{ready}, \texttt{failed}, \texttt{deleted}\\ \hline
\path{exports} & Tệp báo cáo xuất (\uc{UC22}) & \path{requested_by}, \path{status} & \texttt{pending}, \texttt{ready}, \texttt{failed}; hạn 7 ngày\\ \hline
\path{daily_space_stats} & Tổng hợp báo cáo & \path{stat_date} & Materialized view\\ \hline
\path{report.v_*} & Mô hình đọc & -- & View chỉ đọc trên schema khác\\ \hline
\path{river.*} & Hàng đợi~\cite{river} & \path{river_job} & Do River di trú\\ \hline
\end{longtable}
}

% ---------------------------------------------------------------------
%  PHỤ LỤC H. DANH MỤC SƠ ĐỒ VÀ MÃ NGUỒN PLANTUML
% ---------------------------------------------------------------------
\section{Danh mục sơ đồ và mã nguồn PlantUML}\label{pl:so-do}

Các sơ đồ trong báo cáo được viết bằng PlantUML~\cite{plantuml}, lưu dưới dạng văn bản trong \texttt{diagrams/} và kết xuất thành ảnh PNG trong \texttt{figures/diagrams/}. Lưu mã nguồn giúp quản lý phiên bản và so sánh thay đổi; kiểu dáng chung được khai báo trong \texttt{\_style.iuml}. Bảng~\ref{tab:danh-muc-so-do} liệt kê sơ đồ và tệp nguồn tương ứng.

{\small
\begin{longtable}{|C{1.3cm}|L{\dimexpr\textwidth-8\tabcolsep-5\arrayrulewidth-1.3cm-5.6cm-2.4cm\relax}|L{5.6cm}|L{2.4cm}|}
\caption{Danh mục sơ đồ và tệp nguồn PlantUML}\label{tab:danh-muc-so-do}\\
\hline
\hd{Hình} & \hd{Tên sơ đồ} & \hd{Tệp nguồn} & \hd{Loại sơ đồ}\\ \hline
\endfirsthead
\multicolumn{4}{l}{\footnotesize\itshape Bảng \thetable\ (tiếp theo)}\\ \hline
\hd{Hình} & \hd{Tên sơ đồ} & \hd{Tệp nguồn} & \hd{Loại sơ đồ}\\ \hline
\endhead
\ref{fig:quy-trinh-thu-cong} & Quy trình đi chợ và sử dụng thực phẩm theo cách thủ công & \path{c1-01-quy-trinh-thu-cong.puml} & Hoạt động\\ \hline
\ref{fig:phan-ra-chuc-nang} & Sơ đồ phân rã chức năng của ứng dụng & \path{c1-03-phan-ra-chuc-nang.puml} & Sơ đồ tư duy\\ \hline
\ref{fig:chu-trinh} & Chu trình nghiệp vụ khép kín đề xuất cho ứng dụng & \path{c1-02-chu-trinh-nghiep-vu.puml} & Khối quy trình\\ \hline
\ref{fig:ngu-canh} & Sơ đồ ngữ cảnh của hệ thống & \path{c1-04-ngu-canh.puml} & Ngữ cảnh\\ \hline
\ref{fig:ipo-tong-quat} & Mô hình Đầu vào – Xử lý – Đầu ra tổng quát của hệ thống & \path{c1-05-ipo-tong-quat.puml} & Khối IPO\\ \hline
\ref{fig:quy-trinh-mua-sam} & Quy trình mua sắm có phân công trong nhóm gia đình & \path{c1-06-quy-trinh-mua-sam-nhom.puml} & Hoạt động\\ \hline
\ref{fig:trang-thai-mat-hang} & Vòng đời của một mặt hàng trong danh sách mua sắm & \path{c1-07-trang-thai-mat-hang.puml} & Trạng thái\\ \hline
\ref{fig:trang-thai-thuc-pham} & Vòng đời của một lô thực phẩm trong tủ lạnh & \path{c1-08-trang-thai-thuc-pham.puml} & Trạng thái\\ \hline
\ref{fig:quy-trinh-ke-hoach} & Quy trình lập kế hoạch bữa ăn gắn với tồn kho & \path{c1-09-quy-trinh-ke-hoach-bua-an.puml} & Hoạt động\\ \hline
\ref{fig:mo-hinh-khai-niem} & Mô hình khái niệm nghiệp vụ & \path{c1-10-mo-hinh-khai-niem.puml} & Lớp khái niệm\\ \hline
\ref{fig:quan-he-tac-nhan} & Biểu đồ quan hệ giữa các tác nhân & \path{c2-01-quan-he-tac-nhan.puml} & Use case\\ \hline
\ref{fig:uc-tong-quan} & Biểu đồ use case tổng quan của hệ thống & \path{c2-02-uc-tong-quan.puml} & Use case\\ \hline
\ref{fig:uc-tai-khoan} & Biểu đồ use case phân rã gói Xác thực và quản lý tài khoản & \path{c2-03-uc-tai-khoan.puml} & Use case\\ \hline
\ref{fig:uc-nhom} & Biểu đồ use case phân rã gói Quản lý nhóm gia đình & \path{c2-05-uc-nhom.puml} & Use case\\ \hline
\ref{fig:act-tham-gia-nhom} & Biểu đồ hoạt động của use case Tham gia nhóm gia đình & \path{c2-06-act-tham-gia-nhom.puml} & Hoạt động\\ \hline
\ref{fig:uc-mua-sam} & Biểu đồ use case phân rã gói Quản lý danh sách mua sắm & \path{c2-07-uc-mua-sam.puml} & Use case\\ \hline
\ref{fig:act-danh-dau-da-mua} & Biểu đồ hoạt động của use case Đánh dấu đã mua và Chuyển vào tủ lạnh & \path{c2-08-act-danh-dau-da-mua.puml} & Hoạt động\\ \hline
\ref{fig:act-sinh-ds} & Biểu đồ hoạt động của use case Sinh danh sách mua sắm từ kế hoạch bữa ăn & \path{c2-09-act-sinh-ds-tu-ke-hoach.puml} & Hoạt động\\ \hline
\ref{fig:uc-tu-lanh} & Biểu đồ use case phân rã gói Quản lý thực phẩm trong tủ lạnh & \path{c2-10-uc-tu-lanh.puml} & Use case\\ \hline
\ref{fig:act-nhac-nho} & Biểu đồ hoạt động của use case Gửi nhắc nhở thực phẩm sắp hết hạn & \path{c2-11-act-nhac-nho.puml} & Hoạt động\\ \hline
\ref{fig:seq-nhac-nho} & Biểu đồ tuần tự của luồng nhắc nhở hạn dùng & \path{c2-12-seq-nhac-nho.puml} & Tuần tự\\ \hline
\ref{fig:uc-cong-thuc} & Biểu đồ use case phân rã gói Quản lý công thức nấu ăn & \path{c2-13-uc-cong-thuc.puml} & Use case\\ \hline
\ref{fig:uc-ke-hoach} & Biểu đồ use case phân rã gói Quản lý kế hoạch bữa ăn & \path{c2-14-uc-ke-hoach.puml} & Use case\\ \hline
\ref{fig:act-goi-y} & Biểu đồ hoạt động của use case Gợi ý món ăn từ thực phẩm sẵn có & \path{c2-15-act-goi-y-mon.puml} & Hoạt động\\ \hline
\ref{fig:act-xac-nhan-nau} & Biểu đồ hoạt động của use case Xác nhận đã nấu món ăn & \path{c2-16-act-xac-nhan-nau.puml} & Hoạt động\\ \hline
\ref{fig:uc-tim-kiem} & Biểu đồ use case phân rã gói Tìm kiếm và phân loại thực phẩm & \path{c2-17-uc-tim-kiem.puml} & Use case\\ \hline
\ref{fig:uc-thong-ke} & Biểu đồ use case phân rã gói Thống kê báo cáo & \path{c2-18-uc-thong-ke.puml} & Use case\\ \hline
\ref{fig:uc-quan-tri} & Biểu đồ use case phân rã gói Quản trị hệ thống & \path{c2-19-uc-quan-tri.puml} & Use case\\ \hline
\multicolumn{4}{|l|}{\hs{Chương 3. Thiết kế backend}}\\ \hline
\ref{fig:c4-ngu-canh} & Sơ đồ ngữ cảnh hệ thống theo mô hình C4 (mức 1) & \path{c3-01-c4-ngu-canh.puml} & C4\\ \hline
\ref{fig:c4-container} & Sơ đồ vùng chứa theo mô hình C4 (mức 2) & \path{c3-02-c4-container.puml} & C4\\ \hline
\ref{fig:giao-tiep-su-kien} & Luồng sự kiện miền: hộp thư đi trong giao dịch và phát tán thời gian thực & \path{c3-03-giao-tiep-su-kien.puml} & Thành phần\\ \hline
\ref{fig:phu-thuoc-module} & Sơ đồ gói thể hiện phụ thuộc giữa các module & \path{c3-04-phu-thuoc-module.puml} & Gói\\ \hline
\ref{fig:c4-thanh-phan-api-vao} & Sơ đồ thành phần của vùng chứa API theo mô hình C4 (mức 3): phía vào & \path{c3-05a-c4-thanh-phan-api-vao.puml} & C4\\ \hline
\ref{fig:c4-thanh-phan-api-ra} & Sơ đồ thành phần của vùng chứa API theo mô hình C4 (mức 3): phía ra & \path{c3-05b-c4-thanh-phan-api-ra.puml} & C4\\ \hline
\ref{fig:c4-thanh-phan-worker} & Sơ đồ thành phần của vùng chứa Worker theo mô hình C4 (mức 3) & \path{c3-06-c4-thanh-phan-worker.puml} & C4\\ \hline
\ref{fig:kien-truc-ung-dung} & Kiến trúc tầng của ứng dụng Flutter và các điểm tiếp xúc với backend & \path{c3-41-kien-truc-ung-dung.puml} & Thành phần\\ \hline
\ref{fig:vong-doi-dong-bo} & Biểu đồ trạng thái vòng đời ứng dụng và kênh đồng bộ trên Android và iOS & \path{c3-42-vong-doi-dong-bo.puml} & Trạng thái\\ \hline
\ref{fig:phan-tang-module} & Phân tầng lục giác bên trong một module & \path{c3-07-phan-tang-module.puml} & Thành phần\\ \hline
\ref{fig:lop-tu-lanh} & Biểu đồ lớp module tủ lạnh (pantry) & \path{c3-08-lop-tu-lanh.puml} & Lớp\\ \hline
\ref{fig:lop-mua-sam-ke-hoach} & Biểu đồ lớp module mua sắm (shopping) & \path{c3-09-lop-mua-sam-ke-hoach.puml} & Lớp\\ \hline
\ref{fig:lop-ke-hoach} & Biểu đồ lớp module kế hoạch bữa ăn (mealplan) & \path{c3-09b-lop-ke-hoach.puml} & Lớp\\ \hline
\ref{fig:pipeline-yeu-cau} & Chuỗi xử lý một yêu cầu HTTP trong tiến trình api & \path{c3-10-pipeline-yeu-cau.puml} & Hoạt động\\ \hline
\ref{fig:trang-thai-tai-khoan} & Biểu đồ trạng thái tài khoản người dùng & \path{c3-11-trang-thai-tai-khoan.puml} & Trạng thái\\ \hline
\ref{fig:trang-thai-mat-hang-be} & Biểu đồ trạng thái mặt hàng mua sắm ở phía máy chủ & \path{c3-12-trang-thai-mat-hang-be.puml} & Trạng thái\\ \hline
\ref{fig:trang-thai-tac-vu} & Vòng đời tác vụ nền trong hàng đợi River & \path{c3-13-trang-thai-tac-vu.puml} & Trạng thái\\ \hline
\ref{fig:seq-xac-nhan-nau} & Biểu đồ tuần tự xác nhận đã nấu với giao dịch nguyên tử và khóa dòng & \path{c3-14-seq-xac-nhan-nau.puml} & Tuần tự\\ \hline
\ref{fig:seq-outbox-mua-hang} & Biểu đồ tuần tự đánh dấu đã mua với hộp thư đi giao dịch & \path{c3-15-seq-outbox-mua-hang.puml} & Tuần tự\\ \hline
\ref{fig:seq-nhap-tu} & Biểu đồ tuần tự chuyển mặt hàng đã mua vào tủ lạnh qua cổng StockIn & \path{c3-15b-seq-nhap-tu.puml} & Tuần tự\\ \hline
\ref{fig:seq-dong-bo-ngoai-tuyen} & Biểu đồ tuần tự đồng bộ ngoại tuyến với phát lại lũy đẳng & \path{c3-16-seq-dong-bo-ngoai-tuyen.puml} & Tuần tự\\ \hline
\ref{fig:seq-thoi-gian-thuc} & Biểu đồ tuần tự đồng bộ thời gian thực qua WebSocket, Redis Pub/Sub và FCM & \path{c3-17-seq-thoi-gian-thuc.puml} & Tuần tự\\ \hline
\ref{fig:seq-nhac-han-dung-be} & Biểu đồ tuần tự nhắc hạn dùng phía backend với ba lớp chống trùng & \path{c3-18-seq-nhac-han-dung-be.puml} & Tuần tự\\ \hline
\ref{fig:seq-tim-kiem-chi-muc} & Biểu đồ tuần tự đánh chỉ mục và truy vấn tìm kiếm có dự phòng & \path{c3-37-seq-tim-kiem-chi-muc.puml} & Tuần tự\\ \hline
\ref{fig:seq-tham-gia-nhom} & Biểu đồ tuần tự hiện thực use case Tham gia nhóm gia đình & \path{c3-29-seq-tham-gia-nhom.puml} & Tuần tự\\ \hline
\ref{fig:seq-phan-cong} & Biểu đồ tuần tự hiện thực use case Phân công người mua & \path{c3-30-seq-phan-cong.puml} & Tuần tự\\ \hline
\ref{fig:seq-sinh-ds-ke-hoach} & Biểu đồ tuần tự hiện thực use case Sinh danh sách mua sắm từ kế hoạch bữa ăn & \path{c3-31-seq-sinh-ds-ke-hoach.puml} & Tuần tự\\ \hline
\ref{fig:seq-them-thuc-pham} & Biểu đồ tuần tự hiện thực use case Thêm thực phẩm vào tủ lạnh & \path{c3-32-seq-them-thuc-pham.puml} & Tuần tự\\ \hline
\ref{fig:seq-goi-y-mon} & Biểu đồ tuần tự hiện thực use case Gợi ý món ăn từ thực phẩm sẵn có & \path{c3-33-seq-goi-y-mon.puml} & Tuần tự\\ \hline
\ref{fig:seq-lap-lich} & Biểu đồ tuần tự hiện thực use case Lên lịch bữa ăn & \path{c3-40-seq-lap-lich.puml} & Tuần tự\\ \hline
\ref{fig:seq-bao-cao-xuat} & Biểu đồ tuần tự hiện thực use case Xem báo cáo lãng phí và Xuất báo cáo & \path{c3-34-seq-bao-cao-xuat.puml} & Tuần tự\\ \hline
\ref{fig:seq-dang-ky-otp} & Biểu đồ tuần tự hiện thực use case Đăng ký tài khoản, Xác thực mã OTP và Khôi phục mật khẩu & \path{c3-35-seq-dang-ky-otp.puml} & Tuần tự\\ \hline
\ref{fig:seq-roi-nhom-chuyen-quyen} & Biểu đồ tuần tự hiện thực use case Chuyển quyền trưởng nhóm, Rời nhóm và Xóa thành viên & \path{c3-36-seq-roi-nhom-chuyen-quyen.puml} & Tuần tự\\ \hline
\ref{fig:seq-quan-tri-cau-hinh} & Biểu đồ tuần tự hiện thực use case Cấu hình tham số hệ thống và thao tác khóa tài khoản & \path{c3-38-seq-quan-tri-cau-hinh.puml} & Tuần tự\\ \hline
\ref{fig:seq-duyet-thuc-pham} & Biểu đồ tuần tự hiện thực thao tác duyệt và gộp thực phẩm riêng của use case Quản lý danh mục thực phẩm & \path{c3-39-seq-duyet-thuc-pham.puml} & Tuần tự\\ \hline
\ref{fig:erd-tong-quan} & Sơ đồ thực thể -- liên kết tổng quan theo schema & \path{c3-20-erd-tong-quan.puml} & Thực thể liên kết\\ \hline
\ref{fig:erd-tai-khoan-nhom} & ERD schema identity và household & \path{c3-21-erd-tai-khoan-nhom.puml} & Thực thể liên kết\\ \hline
\ref{fig:erd-mua-sam-tu-lanh} & ERD schema shopping và pantry & \path{c3-22-erd-mua-sam-tu-lanh.puml} & Thực thể liên kết\\ \hline
\ref{fig:erd-cong-thuc-ke-hoach} & ERD schema catalog, recipe và mealplan & \path{c3-23-erd-cong-thuc-ke-hoach.puml} & Thực thể liên kết\\ \hline
\ref{fig:erd-thong-bao-he-thong} & ERD các schema notification, sync, media, admin và report & \path{c3-24-erd-thong-bao-he-thong.puml} & Thực thể liên kết\\ \hline
\ref{fig:seq-dang-nhap-lam-moi} & Biểu đồ tuần tự đăng nhập và làm mới phiên xoay vòng & \path{c3-25-seq-dang-nhap-lam-moi.puml} & Tuần tự\\ \hline
\ref{fig:seq-dung-lai-ma-lam-moi} & Biểu đồ tuần tự phát hiện dùng lại mã làm mới & \path{c3-25b-seq-dung-lai-ma-lam-moi.puml} & Tuần tự\\ \hline
\ref{fig:trien-khai} & Sơ đồ triển khai giai đoạn 1 trên một VPS & \path{c3-26-trien-khai.puml} & Triển khai\\ \hline
\ref{fig:cicd} & Biểu đồ hoạt động của pipeline CI/CD trên GitLab & \path{c3-27-cicd.puml} & Hoạt động\\ \hline
\ref{fig:mo-rong} & Sơ đồ triển khai giai đoạn 2 với nhiều nút & \path{c3-28-mo-rong.puml} & Triển khai\\ \hline
\end{longtable}
}

Sau khi chỉnh sửa sơ đồ, chạy \texttt{make diagrams} tại thư mục gốc để kết xuất lại các tệp \texttt{diagrams/*.puml} thành ảnh PNG. Các sơ đồ use case, trạng thái, lớp và ngữ cảnh dùng bộ dàn trang Smetana tích hợp trong PlantUML, không cần Graphviz. Hướng dẫn chi tiết có trong \texttt{README.md}.

Tệp kiểu dáng dùng chung được trình bày dưới đây. Mỗi sơ đồ nạp tệp này bằng chỉ thị \texttt{!include \_style.iuml}; các sơ đồ UML của Chương~\ref{chap:backend} nạp \texttt{\_style-c3.iuml}, tệp kế thừa \texttt{\_style.iuml} và tăng cỡ chữ để chữ khi in không nhỏ hơn 6,5 pt, còn sơ đồ C4 dùng thư viện chuẩn C4 có sẵn trong PlantUML.

\codecaption{Tệp kiểu dáng dùng chung \_style.iuml}{lst:style}
\VerbatimInput{diagrams/_style.iuml}

Mã nguồn dưới đây là biểu đồ phân rã gói Xác thực và quản lý tài khoản (Hình~\ref{fig:uc-tai-khoan}), minh họa cách khai báo tác nhân, use case, quan hệ bao gồm và tác nhân phụ.

\codecaption{Tệp c2-03-uc-tai-khoan.puml}{lst:uc-tai-khoan}
\VerbatimInput{diagrams/c2-03-uc-tai-khoan.puml}

Biểu đồ hoạt động Gợi ý món ăn (Hình~\ref{fig:act-goi-y}) minh họa cách mô tả làn, vòng lặp và nhánh xử lý trong PlantUML.

\codecaption{Tệp c2-15-act-goi-y-mon.puml}{lst:act-goi-y}
\VerbatimInput{diagrams/c2-15-act-goi-y-mon.puml}

